\documentclass[11pt]{article}

% Core packages
\usepackage[utf8]{inputenc}
\usepackage[T1]{fontenc}
\usepackage{lmodern}
\usepackage[margin=25mm]{geometry}
\usepackage{setspace}
\usepackage{xcolor}
\usepackage{textcomp}
\usepackage[version=4]{mhchem}
\usepackage{xurl}

% Unicode guard: quoted source text is fatal if unhandled (latex-arxiv-tooling lesson)
\usepackage{newunicodechar}
\newunicodechar{’}{'}
\newunicodechar{‘}{`}
\newunicodechar{“}{{\tt``}}
\newunicodechar{”}{{\tt''}}
\newunicodechar{–}{--}
\newunicodechar{—}{---}
\newunicodechar{…}{\ldots}
\newunicodechar{≈}{\ensuremath{\approx}}
\newunicodechar{≤}{\ensuremath{\leq}}
\newunicodechar{≥}{\ensuremath{\geq}}
\newunicodechar{→}{\ensuremath{\rightarrow}}
\newunicodechar{°}{\textdegree}
\newunicodechar{×}{\ensuremath{\times}}
\newunicodechar{²}{\textsuperscript{2}}
\newunicodechar{³}{\textsuperscript{3}}
\newunicodechar{µ}{\textmu}
\newunicodechar{±}{\ensuremath{\pm}}
\newunicodechar{−}{\ensuremath{-}}
\newunicodechar{Δ}{\ensuremath{\Delta}}
\newunicodechar{™}{\texttrademark}
\newunicodechar{α}{\ensuremath{\alpha}}
\newunicodechar{β}{\ensuremath{\beta}}
\newunicodechar{γ}{\ensuremath{\gamma}}
\newunicodechar{κ}{\ensuremath{\kappa}}
\newunicodechar{μ}{\textmu}
\newunicodechar{ζ}{\ensuremath{\zeta}}
\newunicodechar{∞}{\ensuremath{\infty}}
\newunicodechar{√}{\ensuremath{\surd}}
\newunicodechar{′}{$'$}
\newunicodechar{‐}{-}
\newunicodechar{‑}{-}
\newunicodechar{·}{\textperiodcentered}
\newunicodechar{®}{\textregistered}
\newunicodechar{ĸ}{k}
\newunicodechar{黄}{Huang}
\newunicodechar{芩}{qin}

% Math, tables, layout
\usepackage{amsmath,amssymb}
\usepackage{booktabs}
\usepackage{longtable}
\usepackage{array}
\usepackage{tabularx}
\usepackage{graphicx}
\usepackage{caption}
\usepackage{enumitem}
\usepackage{titlesec}
\usepackage{parskip}

% Hyperlinks and metadata
\usepackage{hyperref}
\hypersetup{
    pdftitle={Scutellaria baicalensis as a Longevity Botanical Medicine: Epigenetics, Gene Expression, and the Longevity Substrate},
    pdfauthor={Walker J. Kirkpatrick, ND},
    pdfkeywords={Scutellaria baicalensis, Chinese skullcap, huangqin, baicalin, baicalein, wogonin, epigenetics, DNA methylation, histone deacetylase, microRNA, longevity, sirtuins, AMPK, senescence, naturopathic medicine},
    pdfsubject={Botanical medicine evidence review for naturopathic practice},
    colorlinks=true,
    linkcolor=blue!60!black,
    citecolor=blue!60!black,
    urlcolor=blue!60!black
}

% Citation macro: every citation passes through \src{} so the cited-id set is
% a single regex. ids come from the retrieval ledger, never from memory.
% \allowbreak lets multi-citation clusters wrap instead of overflowing.
\newcommand{\src}[1]{\allowbreak[#1]}
\setlength{\emergencystretch}{3em}

\titleformat{\section}{\large\bfseries}{\thesection}{1em}{}
\titleformat{\subsection}{\normalsize\bfseries}{\thesubsection}{1em}{}
\setstretch{1.18}

\title{\Large \textbf{Scutellaria baicalensis (Chinese Skullcap): A Longevity Botanical Medicine \\ A Grounded Review of Epigenetics, Gene Expression, and the Longevity Substrate \\ for Naturopathic Practice}}
\author{\normalsize Walker J. Kirkpatrick, ND}
\date{\normalsize November 2026}

\begin{document}

\maketitle

\begin{abstract}
\noindent
\textit{Scutellaria baicalensis} Georgi --- Chinese skullcap, huangqin, the ``golden root'' of the classical Chinese materia medica --- has entered the longevity conversation on the strength of a remarkable mechanistic literature. This review synthesizes a retrieval ledger of 293 sources fetched and registered during the project, weighted toward work published between 2019 and 2026, and it organizes the answer the way the evidence itself is organized: epigenetics and gene expression first. That foundation is real. Baicalein induces DNMT1-mediated promoter demethylation with gene re-expression in three independent systems, inhibits HDAC-1/8 with proteasomal degradation of HDAC1, suppresses histone H3K4me3 through direct inhibition of adenosylhomocysteinase, and down-regulates the m$^6$A demethylase FTO. Baicalin is a characterized LSD1/KDM1A inhibitor (IC$_{50}$ 3.01\,$\mu$M) whose target has been engaged in human metastatic tissue ex vivo. Wogonin down-regulates the acetyltransferase p300, stabilizes TET2, and demethylates the KEAP1 promoter. MicroRNA remodeling (miR-21 down, miR-139-3p up), whole-transcriptome programs (an Nrf2/ARE hepatic program; a JAK--STAT program in SAMP8 aging mice), and the founding 2001 demonstration of inducible iNOS and COX-2 gene-expression suppression complete a genuinely multi-layer epigenetic pharmacology. Downstream, the longevity circuitry behaves as that model predicts: baicalein extends lifespan in \textit{C. elegans} in five independent studies (up to 45/57/24\% mean/median/maximum in the first, a clean $\sim$10\% median effect in the most rigorous) and is the only constituent with lifespan data --- baicalin and wogonin failed their one direct test; AMPK, mTOR/autophagy, sirtuins, FOXO/insulin signaling, Nrf2, and NF-$\kappa$B are each modulated in preclinical models; senescent-cell pharmacology is senomorphic (SASP suppression, strongest from minor constituents), not senolytic; and the one lifespan screen that asked which epigenetic enzyme the lifespan benefit requires answered cbp-1, the worm ortholog of the p300/CBP acetyltransferase that wogonin down-regulates in mammals. The human evidence is honest about its ceiling: purified baicalein completed phase I at 100--2800 mg in Chinese volunteers; a 374-patient trial of baicalin 500 mg/day improved lipids and hs-CRP; flavocoxid, a baicalin--catechin blend, matched naproxen in two osteoarthritis randomized trials; and no human study has measured a longevity endpoint, nor has any study of any \textit{Scutellaria} constituent touched an epigenetic aging clock --- the single most consequential gap in this literature. Safety resolves into a rare, idiosyncratic liver-injury signal concentrated in concentrated blends and suprapharmacopoeial doses, a germander-adulteration history that owns most of ``skullcap hepatotoxicity,'' wide preclinical margins (26-week rat NOAEL 2{,}000 mg/kg/day), and one genotype-dependent human drug interaction that matters (rosuvastatin AUC reduced 47\% via OATP1B1). The paper closes with the formulary view: what a naturopathic physician can defensibly say, prescribe, and monitor today, and the five studies that would move this plant from mechanistically credible toward clinically proven as a longevity botanical.
\end{abstract}

\section{Introduction: The Root That Cools Fire}

There are botanicals a physician reaches for because tradition says so, and botanicals whose mechanistic literature has quietly outgrown their reputation. \textit{Scutellaria baicalensis} Georgi --- huangqin, the ``golden herb'' of Chinese medicine --- is increasingly the second kind.\src{81} Its dried root entered the written materia medica in the \textit{Shennong Ben Cao Jing} around 200 AD, giving this plant a documented continuous chart of more than two thousand years.\src{101} The classical physician Zhang Zhongjing built it into the \textit{Shanghan Lun} formulas that Chinese and Japanese medicine still prescribe today, most famously Xiao Chai Hu Tang (Japanese Sho-saiko-to), the seven-herb decoction in which huangqin is paired with bupleurum.\src{81}\src{93} Traditional indications --- ``all kinds of fever, jaundice, diarrhea, edema, and sores'' --- read, to a modern eye, as an inflammatory-conditions list.\src{54}

The question this review asks is narrower and harder than ``is huangqin useful.'' It is whether this plant earns a place in a longevity formulary --- and if so, on what evidentiary basis. The organizing premise, at the client's direction, is that the epigenetic and gene-expression literature is the foundation: not an interesting aside, but the load-bearing story that explains why one root touches inflammation, fibrosis, senescence, and lifespan at once. A plant whose constituents inhibit DNMT1, HDAC1/8, LSD1, and the one-carbon enzyme AHCY, remodel microRNA profiles, and shift whole-transcriptome programs is not acting the way single-target drugs act.\src{142}\src{131}\src{117}\src{156} It is acting like a rheostat on the cell's gene-expression machinery --- turning entire programs up and down --- and that is precisely the pharmacologic shape a longevity botanical would need to have, because aging itself is not one target but a drift of many programs.\src{109}

The method is deliberately paranoid about evidence. Every numbered claim below resolves to one of 293 sources fetched and registered during this review, weighted toward work published between 2019 and 2026 with seminal older anchors retained; most cited sources carry verbatim sentences captured from the fetched text itself, so a colleague can check any sentence against its number and find the trail leads somewhere real (see the Sources section). Six parallel research streams produced the evidence base --- phytochemistry and identity, epigenetics, longevity signaling pathways, senescence and model organisms, human trials and pharmacokinetics, and safety --- and each stream was instructed to flag, rather than round off, every lead it could not verify. Those flags survive into this paper: several plausible-sounding claims (a thyroid-cancer demethylation story, a ``Korean'' phase I trial, lifespan data for baicalin) died under verification and are recorded here as not found, because the negative space is as informative as the findings.

Three registers of evidence run through what follows, and the paper keeps them separate: a two-thousand-year clinical tradition (real data, but of a kind that cannot speak to mechanism); a deep preclinical mechanism literature (mostly cells and rodents, oncology-skewed, with dose-translation problems the paper quantifies rather than hides); and a human trial record that is young, small, and --- for longevity endpoints --- empty. Where the registers disagree, the disagreement is stated. The thesis, previewed honestly: \textit{S. baicalensis} is a mechanistically credible longevity candidate whose epigenetic pharmacology is unusually well mapped for a botanical, whose organismal lifespan evidence rests on a single constituent in a single invertebrate, and whose human longevity evidence does not yet exist. What a clinician does with that asymmetry is a judgment call this paper equips rather than makes.

\section{Identity, Sourcing, and the Two Skullcaps}
\label{sec:identity}

Before pharmacology comes botany, and skullcap's botany has a practical edge. The species is \textit{Scutellaria baicalensis} Georgi (Lamiaceae), validly published in 1775, accepted by Kew's Plants of the World Online with a native range from Siberia through the Russian Far East to Vietnam.\src{76}\src{98} The medicinal organ is the dried root, which the Chinese Pharmacopoeia requires to contain no less than 9\% baicalin (2015 edition), with comparable monographs in the Japanese ($\geq$10\% total flavones), European ($\geq$9\%), US ($\geq$11.0\%), WHO, and ISO standards.\src{101}\src{100} (One 2024 review states an 8\% minimum --- an internal inconsistency in the secondary literature worth knowing about when reading label claims.\src{54}) Content is not a fixed constant: flavonoid composition varies with cultivation region, growth year --- the optimum harvest window is two to three years --- and post-harvest processing, with drying and cutting-drying the most composition-preserving methods studied.\src{73}\src{104}

The distinction that matters clinically is between two different plants both sold as ``skullcap.'' American skullcap, \textit{S. lateriflora}, uses the aerial parts, and the flavonoid geography is organ-specific and species-specific: the four major root flavonoids of \textit{S. baicalensis} are absent from its own aerial parts, while in \textit{S. lateriflora} they concentrate in the leaves --- baicalin localizes to roots in the Chinese species but to leaves in the American one.\src{94}\src{75} A common botanical-medicine claim that American skullcap contains ``little baicalin'' did not survive verification and is better stated as a difference of organ and species accumulation pattern.\src{94}\src{72} The two are not interchangeable, and Japanese pharmacopoeial practice treats the roots of both species as pharmaceutical raw materials while classifying aerial parts as non-pharmaceutical.\src{94}

The hepatotoxicity history attached to the word ``skullcap'' belongs largely to this sourcing problem, not to the plant itself. LiverTox attributes the classic skullcap liver-injury reports to multiple co-ingested herbals and, in some cases, to phytochemical analysis identifying significant adulteration with germander (\textit{Teucrium} spp.) or outright mislabeling.\src{64} Germander adulteration is detectable and documented: LC-UV/MS and microscopy distinguish \textit{S. lateriflora} from \textit{T. canadense} and \textit{T. chamaedrys} by phenylpropanoid markers absent from \textit{Scutellaria} itself, and HPTLC methods detect \textit{T. chamaedrys} in \textit{S. baicalensis} raw material and finished extracts down to 0.5\% w/w.\src{80}\src{82} Section~\ref{sec:safety} returns to the liver question in full; here the sourcing lesson is that authentication is the safety intervention.

Commercial quality is, bluntly, uneven --- and this is the kind of number a formulary remembers. In one published analysis of commercial \textit{S. baicalensis} extracts, four of ten contained no baicalin at all, and potency varied more than eleven-fold across products.\src{102} Modern authentication is mature even where commerce is not: four variable chloroplast DNA regions differentiate true \textit{S. baicalensis} from adulterant species, and whole-chloroplast-genome super-barcoding distinguishes cultivated varieties that morphology cannot separate before flowering.\src{89}\src{84} The practical rule falls out of the literature directly: prescribe material with a species identity, a stated baicalin content against a pharmacopoeial marker, and third-party contaminant testing --- and treat any unlabeled ``skullcap'' product as a question, not an answer.

\section{Phytochemistry: The Flavone Substrate}
\label{sec:phyto}

The root's defining chemistry is a family of root-specific 4$'$-deoxyflavones --- more than thirty flavone types have been catalogued --- of which six are the major bioactives: three aglycones (baicalein, wogonin, oroxylin A) and their glucuronide glycosides (baicalin, wogonoside, oroxyloside), all pharmacologically active in their own right.\src{88}\src{54} Baicalin is simply baicalein-7-O-glucuronide (\ce{C21H18O11}); the glycoside--aglycone pair is interconverted by the body and, decisively, by the gut microbiome --- a pharmacokinetic fact that shapes everything in section~\ref{sec:human}.\src{54}\src{234} Beyond the six majors, the root carries skullcapflavones I and II, scutellarin and its aglycone scutellarein, apigenin, and chrysin, among others.\src{99}\src{73} Minor does not mean inert: the strongest senescence-associated secretory phenotype (SASP) inhibitors in the root are minor constituents, a fact with real clinical implications (section~\ref{sec:senes}).\src{211}

Traditional properties frame the same chemistry in a different language. In TCM theory the root is bitter and cold, entering the lung, gallbladder, spleen, large intestine, and small intestine meridians; it clears heat, dries dampness, purges fire, detoxifies, stops bleeding, and --- notably for a longevity conversation --- is traditionally credited with stabilizing pregnancy.\src{79} Processing is pharmacologically real: wine-processing (pao zhi) measurably changes compound dissolution, and growth-year and drying decisions move flavonoid content in documented, quantifiable ways.\src{92}\src{73} The root-specific accumulation pattern itself reflects an evolutionary divergence of the flavonoid pathway within the genus, which is why \textit{S. baicalensis} root is chemically singular rather than merely abundant.\src{75}

The three constituents this paper follows are baicalein, baicalin, and wogonin, because the mechanistic literature follows them. Where oroxylin A and the minor flavones carry unique evidence --- SASP suppression, blood--brain barrier penetration, a completed phase I trial --- they are named explicitly.

\section{The Epigenetic Foundation}
\label{sec:epigen}

\subsection{Why epigenetics is the right first chapter}

A botanical cannot be a longevity medicine by acting on one target, because aging is not one target. The geroscience consensus picture is a drift of interconnected programs --- inflammatory gene expression, nutrient sensing, mitochondrial quality control, senescence signaling --- and a credible longevity botanical must therefore act \textit{programmatically}. The \textit{S. baicalensis} literature has quietly become one of the best-mapped examples of exactly that in the botanical world: its constituents have documented, mechanistically dissected interactions with DNA methylation enzymes, histone acetyltransferases and deacetylases, histone methyltransferases and demethylases, the ten-eleven translocation dioxygenase TET2, the m$^6$A RNA demethylase FTO, microRNA processing, and whole-transcriptome regulatory programs.\src{142}\src{143}\src{156}\src{153} The founding observation is old enough to vote: in 2001, baicalein, baicalin, and wogonin were shown to suppress inducible nitric oxide synthase gene expression in LPS-stimulated macrophages --- the first modern gene-expression finding for this herb, and one that already displayed the family signature of program-level action.\src{105}

Table~\ref{tab:epigen} maps the foundation. The subsections walk it layer by layer, and the final subsection is the honesty box: where this evidence comes from, at what concentrations, and what it does not yet show.

\begin{table}[ht]
\caption{The epigenetic enzyme targets of \textit{S. baicalensis} constituents, as verified in this review's citation ledger. Context matters: most rows were demonstrated in oncology or fibrosis models (section~\ref{sec:limits}).}
\label{tab:epigen}
\footnotesize
\setlength{\tabcolsep}{4pt}
\begin{tabular}{>{\raggedright\arraybackslash}p{1.4cm} >{\raggedright\arraybackslash}p{2.6cm} >{\raggedright\arraybackslash}p{5.2cm} >{\raggedright\arraybackslash}p{4.3cm}}
\toprule
Flavone & Epigenetic target & Demonstrated effect & Model context \\
\midrule
Baicalein & DNMT1 & Promoter demethylation with gene re-expression (SHP-1, GABRD, SCARA5) & Imatinib-resistant CML CD34$^+$ cells; status-epilepticus rat hippocampus in vivo; hepatic stellate cells\src{142}\src{154}\src{157} \\
Baicalein & HDAC-1/8 & Direct inhibition plus ubiquitin--proteasome degradation of HDAC1; $\uparrow$ histone H3 acetylation; p53 re-acetylation & Core-binding-factor AML; cutaneous T-cell lymphoma\src{131}\src{126} \\
Baicalein & AHCY & Direct inhibition (residues D134/D303) $\rightarrow$ S-adenosylhomocysteine accumulation $\rightarrow$ $\downarrow$ H3K4me3, loss of OCT4/Nanog stemness & MSI colorectal cancer; validated in patient tissues and PDX\src{156} \\
Baicalein & FTO (m$^6$A) & $\downarrow$ ZEB1 mRNA stability via enhanced m$^6$A modification & Pancreatic cancer EMT\src{153} \\
Baicalein & KDM4E & Enhanced demethylase activity; $\downarrow$ H3K9me3 at BICD1 promoter & Triple-negative breast cancer\src{147} \\
Baicalein & microRNAs & miR-21 $\downarrow$ (fibroblasts; bleomycin-fibrotic rat lung in vivo); miR-139-3p $\uparrow$; miR-29 $\uparrow$; miR-424-3p $\downarrow$ (PTEN $\uparrow$) & Fibrotic lung; pancreatic, hepatoma lines\src{121}\src{111}\src{136}\src{123}\src{120} \\
Baicalin & LSD1/KDM1A & Direct inhibition, IC$_{50}$ 3.01\,$\mu$M; reduced CD206/PD-L1 in tumor-associated macrophages on pharmacological inhibition & Enzyme kinetics; patient-derived colorectal liver metastasis specimens ex vivo\src{117}\src{162} \\
Baicalin & non-coding RNAs & miR-139-3p $\uparrow$ (CDK16); miR-124 $\uparrow$ (TLR4/NF-$\kappa$B); lncRNA MALAT1 $\downarrow$; circ\_0000734 axis & Colon cancer; diabetic-nephropathy mice; RA synoviocytes and patient PBMCs\src{141}\src{140}\src{159} \\
Wogonin & p300 (HAT) & $\downarrow$ p300 expression; $\downarrow$ histone H3/H4 crotonylation at the HIV-1 promoter; blocked latent viral reactivation & Latency models; primary CD4$^+$ T cells from ART-suppressed humans ex vivo\src{143} \\
Wogonin & HDAC1/2 & $\downarrow$ protein levels & Lung adenocarcinoma\src{113} \\
Wogonin & TET2 & Stabilization with $\uparrow$ 5hmC via an AMPK--TET2 axis & Ovarian cancer\src{223} \\
Wogonin & KEAP1 promoter & $\downarrow$ CpG methylation, $\uparrow$ Keap1 protein & Radioresistant breast cancer\src{145} \\
Wogonin & LSD1 & Blocks LSD1-mediated NF-$\kappa$B p65 demethylation; macrophage M1$\rightarrow$M2 shift & Osteoarthritis mice, an aging-relevant joint model\src{155} \\
\bottomrule
\end{tabular}
\end{table}

\subsection{DNA methylation: a convergent DNMT1 story}

The methylation literature converges on one enzyme: DNMT1, the maintenance methyltransferase --- not the de novo enzymes DNMT3a/3b, for which no \textit{Scutellaria} evidence was found despite targeted searching. In imatinib-resistant chronic myeloid leukemia CD34$^+$ cells, baicalein induced DNMT1-dependent demethylation of the SHP-1 promoter with re-expression of SHP-1 and downstream inhibition of JAK2/STAT5.\src{142} In a lithium-pilocarpine rat model of status epilepticus --- a non-cancer, neurological model, and therefore important for the aging framing --- baicalein interacted with DNMT1 to suppress methylation of the GABRD promoter and restore its protein in hippocampus.\src{154} In hepatic stellate cells and mouse fibrosis models, baicalein targeted DNMT1-facilitated methylation of the SCARA5 promoter, derepressing SCARA5 and driving ferroptosis of the activated cells.\src{157} For wogonin, the demethylation target sits squarely in the longevity circuitry: reduced CpG methylation of the KEAP1 promoter, raising Keap1 protein and dampening the Nrf2/HIF-1$\alpha$ axis in radioresistant breast cancer.\src{145} One demethylation-adjacent finding is a combination, not a single constituent --- baicalin with rosmarinic acid derepressed PPAR$\gamma$ in stellate cells via reduced MeCP2 recruitment and EZH2 suppression --- and it is the only EZH2/H3K27 signal in this literature.\src{108} The convergence is the point: three independent groups, three different promoters, one enzyme.\src{142}\src{154}\src{157}

\subsection{Histone chemistry: acetylation, methylation, and the demethylases}

Baicalein is a characterized natural HDAC-1/8 inhibitor: in core-binding-factor acute myeloid leukemia it inhibited both isoforms, combined mild direct inhibition with ubiquitin--proteasome degradation of HDAC1, and raised histone H3 acetylation; the same group extended the finding to cutaneous T-cell lymphoma, restoring p53 acetylation through HDAC8 inhibition.\src{131}\src{126} Wogonin works the other side of the same ledger, down-regulating the histone acetyltransferase p300 --- and in the most striking paper in this layer, it suppressed latent HIV-1 reactivation in primary CD4$^+$ T cells from ART-suppressed individuals ex vivo by decreasing histone H3/H4 \emph{crotonylation} at the viral promoter.\src{143} Wogonin also reduces HDAC1 and HDAC2 protein levels in lung carcinoma cells, so both directions of the acetylation rheostat are documented for the genus.\src{113}

The demethylase story has two arms. Baicalin was characterized as a natural LSD1 (KDM1A) inhibitor with an IC$_{50}$ of 3.01\,$\mu$M --- among the tighter enzymatic constants in this literature --- and pharmacological LSD1 inhibition with baicalin reduced CD206 and PD-L1 in tumor-associated macrophages from patient-derived colorectal liver metastasis specimens, the strongest human-material epigenetic evidence for any constituent here, though ex vivo.\src{117}\src{162} Wogonin binds LSD1 in osteoarthritic joints, blocking LSD1-mediated demethylation of NF-$\kappa$B p65 and thereby inhibiting JNK/NF-$\kappa$B activation --- an epigenetic mechanism demonstrated inside an aging-relevant disease model, not just a cancer line.\src{155} Baicalein, for its part, suppresses histone methylation indirectly through one-carbon metabolism: it directly inhibits adenosylhomocysteinase (AHCY), causing S-adenosylhomocysteine accumulation, decreased H3K4me3, and loss of OCT4/Nanog stemness in microsatellite-instability colorectal cancer --- validated in patient tissues and patient-derived xenografts.\src{156} In the opposite direction it enhances the H3K9 demethylase KDM4E in triple-negative breast cancer.\src{147} And wogonin stabilizes TET2, increasing 5-hydroxymethylcytosine via an AMPK--TET2 axis --- active DNA demethylation, the layer most often missed in botanical reviews.\src{223}

\subsection{Non-coding RNA}

The microRNA layer is where single sentences need care, because direction depends on constituent and tissue --- that is data, not contradiction. Baicalein suppresses miR-21 --- the profibrotic, proinflammatory oncomiR --- in TGF-$\beta$1-treated human lung fibroblasts, with miR-21 mimics blunting its antifibrotic effect, and in vivo in bleomycin-induced pulmonary fibrosis in rats.\src{121}\src{111} Both baicalein (pancreatic cancer) and baicalin (colon cancer) up-regulate miR-139-3p --- note the 3p arm; the 5p arm that circulates in review folklore was not what the primary studies found.\src{136}\src{141} Baicalein broadly remodels miRNA profiles in hepatoma cells with targets enriched in MAPK, PI3K--Akt, Wnt, Hippo, and mTOR pathways --- the longevity circuitry again --- and separately up-regulates miR-29 in fibroblasts and down-regulates miR-424-3p to raise PTEN.\src{119}\src{123}\src{120} Baicalin up-regulates miR-124, which targets TLR4 in diabetic-nephropathy mice, and down-regulates lncRNA-MALAT1 with compensatory miR-200c induction in breast cancer cells.\src{141}\src{140} The long non-coding RNA AK021954 mediates baicalein's inhibition of vascular smooth-muscle proliferation --- an atherosclerosis-relevant, non-cancer model.\src{115} In human rheumatoid material, baicalin acted through the circ\_0000734/miR-197-5p/IKBKB axis in synoviocytes, and circ\_0000734 levels in patient PBMCs correlated with RF, anti-CCP, CRP, ESR, and DAS28.\src{159} A systematic review has now catalogued the baicalein--miRNA literature in cancer, for readers who want the full map.\src{152}

\subsection{Whole-transcriptome programs}

When the transcriptome is read whole, the programs that surface are the aging-relevant ones. In HepG2 cells, a 42K-probe microarray found baicalein up-regulating 440 and down-regulating 254 probes, with the Nrf2/ARE antioxidant program the top hit and ARE-luciferase transactivation confirmed.\src{110} In SAMP8 senescence-accelerated mice --- the most directly aging-relevant transcriptome study in this literature --- baicalein regulated fifty hub genes with JAK--STAT the most significant pathway, prevented cortical JAK2/STAT1 activation, and lowered A$\beta$1-42 and RAGE.\src{124} In \textit{C. elegans}, microarray work accompanying the lifespan studies found baicalein down-regulating the age-related genes \textit{mtor} and \textit{parp-1} --- a direct transcriptional read on two canonical aging genes.\src{212} Gastric-cancer RNA-seq after baicalin is dominated by the p53 pathway; single-cell and spatial transcriptomics have defined a ten-gene baicalein-response signature in lung cancer.\src{138}\src{158} And at the RNA-modification layer, baicalein down-regulates the m$^6$A demethylase FTO, destabilizing ZEB1 mRNA --- epitranscriptomic regulation, a fifth layer of the gene-expression story.\src{153}

\subsection{The honesty box: concentrations, contexts, and the one lifespan bridge}

Three cautions keep this foundation honest. First, provenance: this literature is oncology-skewed, because epigenetic enzyme assays live in cancer pharmacology; the non-cancer models (epileptic brain, fibrotic lung, arthritic joint, SAMP8 mouse, \textit{C. elegans}) are the exceptions a longevity review must lean on and label.\src{154}\src{124}\src{212} Second, concentration: most in-vitro epigenetic readouts were obtained at 25--100\,$\mu$M, while oral baicalein at 200--800 mg/day produces total plasma peaks of roughly 280--845 ng/mL in humans, mostly as conjugates --- a two-order-of-magnitude gap this paper does not paper over (section~\ref{sec:human} quantifies it, and section~\ref{sec:limits} keeps it).\src{37} Third, direction: epigenetic modulation by these flavones is not uniformly benign --- baicalin acted as an independent fibrogenic factor and \emph{aggravated} isoproterenol-induced myocardial fibrosis via LSD1/AP-1 signaling with increased histone H3 serine-10 phosphorylation.\src{160} A rheostat turns both ways.

And one finding deserves the last word in this section because it is the whole paper in miniature. In a screen of glucose-toxicity-protective drugs, baicalein extended \textit{C. elegans} lifespan --- and the protection was abolished by RNAi of \textit{cbp-1}, the worm ortholog of the CBP/p300 acetyltransferase.\src{109} The lifespan effect of this plant's best-studied constituent ran through an epigenetic enzyme. That is the foundation's foundation: not a metaphor, a gene.

\section{From Epigenome to Longevity Pathways}
\label{sec:pathways}

\subsection{The insulin/IGF-1 axis: the one direct lifespan experiment}

The single most consequential fact in this literature is also its cleanest: baicalein extends lifespan in \textit{Caenorhabditis elegans}, and no other constituent of the plant has published lifespan data at all. Five studies, 2013--2025, tell a consistent effect with a contested mechanism (Table~\ref{tab:lifespan}). The 2013 founding study reported extensions of 45\%, 57\%, and 24\% in mean, median, and maximum lifespan, activating the Nrf2 homolog SKN-1 in intestinal nuclei.\src{200} A 2016 study found SKN-1 required and DAF-16 dispensable; a 2021 head-to-head comparison of six related flavones at 100\,$\mu$M found baicalein best at an 18.6\% extension, dependent on SKN-1/Nrf2, with \textit{mtor} and \textit{parp-1} down-regulated; a 2023 study implicated ELT-3 and a klotho homolog; and the rigorous 2025 MedComm study --- wild-type N2 worms, seven mutant and three reporter strains, dietary-restriction confounds ruled out --- localized the median-lifespan gain (20 to 22 days at 100\,$\mu$M) to suppression of DAF-2 insulin/IGF-1 signaling with nuclear translocation of the FOXO homolog DAF-16 and \emph{no} SKN-1 translocation.\src{203}\src{212}\src{218}\src{167} The pathway attribution is genuinely unresolved --- SKN-1 without DAF-16 in two studies, DAF-16 without SKN-1 in one, both in another --- but the effect itself replicates in every laboratory that has looked.\src{167}

Two boundary results complete the frame. In \textit{Drosophila}, baicalein at 0.04--0.2 mg/mL in food extended mean and maximum lifespan and reduced oxidative damage in induced-aging models.\src{209} But the only direct lifespan test of baicalin and wogonin --- a 2026 \textit{Drosophila} study feeding American skullcap (\textit{S. lateriflora}) extracts --- found improved mobility (attributed to baicalin) and sleep (a wogonin mix) with no lifespan extension at all; the authors concluded the compounds improve neuronal health ``and do not generally slow down the aging process.''\src{226} For a formulary, that asymmetry matters: the longevity claim belongs to baicalein, and to the aglycone side of the glycoside pair.

\begin{table}[ht]
\caption{Published lifespan studies of \textit{S. baicalensis} constituents. The 2026 \textit{Drosophila} result bounds what can be claimed for baicalin and wogonin.}
\label{tab:lifespan}
\footnotesize
\setlength{\tabcolsep}{4pt}
\begin{tabular}{>{\raggedright\arraybackslash}p{1.8cm} >{\raggedright\arraybackslash}p{2.7cm} >{\raggedright\arraybackslash}p{2.1cm} >{\raggedright\arraybackslash}p{3.9cm} >{\raggedright\arraybackslash}p{3.9cm}}
\toprule
Year, organism & Design & Dose & Lifespan effect & Pathway attribution \\
\midrule
2013, \textit{C. elegans} & N2 wild type & not stated (abstract) & $+$45\% mean, $+$57\% median, $+$24\% maximum & SKN-1/Nrf2 activation\src{200} \\
2016, \textit{C. elegans} & stress-resistance panel & dose--response & extension; heat and arsenite resistance & SKN-1 required; DAF-16 independent\src{203} \\
2021, \textit{C. elegans} & head-to-head, six flavones, Lifespan Machine & 100\,$\mu$M & $+$18.6\% (best of six) & SKN-1 dependent; \textit{mtor}, \textit{parp} down-regulated\src{212} \\
2023, \textit{C. elegans} & Wnt/GATA reporter work & --- & extension; stress amelioration & ELT-3 $+$ SKN-1 (lifespan); DAF-16 (stress)\src{218} \\
2025, \textit{C. elegans} & N2 $+$ 7 mutants $+$ 3 reporters; dietary-restriction confounds ruled out & 100\,$\mu$M & median 20 to 22 days ($\sim$10\%) & DAF-2 $\downarrow$ / DAF-16 nuclear translocation; no SKN-1 shift\src{167} \\
2017, \textit{Drosophila} & natural and induced aging & 0.04--0.2 mg/mL food & mean and maximum extended & antioxidation\src{209} \\
2026, \textit{Drosophila} & \textit{S. lateriflora} extracts (baicalin, wogonin arms) & 1--5 mg/mL food & \textbf{null} --- no extension; mobility and sleep improved & neuronal function, not slowed aging\src{226} \\
\bottomrule
\end{tabular}
\end{table}

\subsection{Sirtuins}

The sirtuin evidence is real but direction-dependent, and saying so plainly is more useful to a clinician than the usual blanket ``activates SIRT1.'' In fasted and high-fat-diet mice, baicalin restrained \emph{excessive} hepatic SIRT1 induction to check glucose overproduction; in olfactory-bulbectomized rats it \emph{raised} SIRT1 in brain while deacetylating NF-$\kappa$B p65 and lowering hippocampal IL-1$\beta$, IL-6, and TNF-$\alpha$ --- one constituent, two directions, two tissues, both mechanistically documented.\src{192}\src{122} SIRT6 is up-regulated by baicalein in cerebral ischemia--reperfusion, deacetylating FOXA2 to promote the cystine transporter SLC7A11 and limit ferroptosis; baicalin in a \textit{Scutellaria}--\textit{Coptidis} combination enhanced ACSL5 via SIRT6 deacetylation in NAFLD models.\src{178}\src{194} SIRT3 restoration is the proven mechanism of baicalein's anti-senescence effect in bleomycin-fibrotic lung --- Sirt3 siRNA abolishes the benefit --- making the mitochondrial sirtuin the load-bearing member of that story.\src{177} A hepatic review catalogs baicalin among SIRT1-active hepatoprotective flavonoids.\src{193}

\subsection{AMPK, mTOR, and autophagy}

Baicalin's AMPK literature is large enough to have earned its own 2025 review, spanning metabolic and inflammatory contexts.\src{183} The anchors: baicalein activated AMPK and suppressed SREBP1 cleavage to reduce hepatic fat in cells and mice; baicalin activated hippocampal AMPK/PGC-1$\alpha$ with NIX-mediated mitophagy in chronically stressed mice, with molecular docking supporting direct AMPK binding; baicalein activated AMPK and inhibited mTOR from 5\,$\mu$M upward in pancreatic neuroendocrine tumor models.\src{163}\src{182}\src{164} Autophagy dependence is proven the way it should be --- by rescue: baicalein induced autophagic cell death through AMPK/ULK1 with mTORC1 down-regulation, and Beclin-1/Vps34/Atg5/Atg7 knockdown completely reversed the effect; baicalin activated protective autophagy through Akt/mTOR/ULK1 in pulmonary fibrosis, with 3-methyladenine and hydroxychloroquine abolishing the antifibrotic benefit; baicalein drove SIRT1/AMPK/mTOR mitophagy in parkinsonian rats, with 3-methyladenine blocking neuroprotection.\src{165}\src{188}\src{172} No retrieved study names an upstream kinase such as LKB1, so the activation mechanism upstream of AMPK remains open.\src{183}

\subsection{FOXO and insulin signaling in mammals}

Downstream of insulin signaling, the mammalian data parallel the worm story: baicalein inhibited AKT to up-regulate FOXO1 in glucocorticoid-treated osteoblasts; baicalin raised AKT/AS160 signaling with GLUT4 translocation and PGC-1$\alpha$ in obese-mouse muscle while lowering HOMA-IR; and baicalein ameliorated insulin resistance through hepatic and muscular PI3K/AKT in a GLP-1R-dependent manner, the effect disappearing in Glp1r-null mice.\src{170}\src{180}\src{175} Insulin/IGF-1 suppression with FOXO activation is the most replicated longevity pathway in experimental biology; it is reassuring, and no more, that a skullcap flavone reads as a mild modulator of it in mammals.\src{167}\src{175}

\subsection{Nrf2: the non-canonical activation and the two directions}

The Nrf2/ARE axis carries some of the best functional data: baicalin protected rat stomach against stress-, ethanol-, and indomethacin-induced ulcers with 72--80\% efficacy attributed to Nrf2-mediated antisecretory action; a neurodegeneration review documents baicalein induction of HO-1, NQO1, GPx, and SOD across models.\src{190}\src{174} Mechanistically, baicalein's Hct116 activation is \emph{non-canonical}: no change in NRF2 transcription, ubiquitinylation, or proteasomal turnover, and no redox dependence --- instead decreased Nrf2 phosphorylation (IC$_{50}$ $\approx$ 20\,$\mu$M), a kinase-inhibition profile.\src{116} The pathway also runs in the epigenetic layer from section~\ref{sec:epigen}: wogonin's KEAP1-promoter demethylation.\src{145} And direction is context-dependent: baicalein activated Nrf2/SLC7A11/GPX4 to inhibit ferroptosis in cardiomyocyte models, while baicalin in cisplatin-treated lung cancer \emph{accumulated} KEAP1 and reduced NRF2/HO-1 to aid chemosensitization --- oncology's opposite polarity, stated here so the reader is not surprised by it in the primary literature.\src{161}\src{149}

\subsection{NF-$\kappa$B inflammaging}

The inflammaging axis is the most consistently suppressed pathway in the corpus. Baicalin dose-dependently blocked LPS-induced CD14/TLR4/MyD88/NF-$\kappa$B p65 activation and IL-6 release in macrophages and mice; baicalin and baicalein reduced PI3K/AKT/p65 phosphorylation in hyperuricemic nephropathy; and baicalin's SIRT1 effect deacetylates p65 at lysine 310 in the depression model --- the sirtuin--NF-$\kappa$B crosstalk that geroscience readers will recognize as the canonical inflammaging brake.\src{189}\src{186}\src{122} The epigenetic layer connects here too: wogonin's LSD1 binding blocks p65 demethylation in arthritic joints, keeping NF-$\kappa$B signaling in check through a histone-demethylase mechanism.\src{155}

\subsection{Mitochondrial quality control}

The mitochondrial file is coherent: baicalin directly bound p38 MAPK and shifted skeletal-muscle fibers from glycolytic toward oxidative through PGC-1$\alpha$; baicalein restored mitochondrial function and biogenesis in rotenone-parkinsonian rats; baicalin reversed A$\beta$-oligomer memory deficits and hippocampal mitochondrial fragmentation through PDE4/PKA/Drp1; the AMPK/PGC-1$\alpha$/NIX mitophagy axis in stressed hippocampus and the SIRT3 restoration in aging fibrotic lung complete a biogenesis-plus-mitophagy signature.\src{191}\src{176}\src{179}\src{182}\src{177}

\subsection{CNS aging pharmacology: MAO-B, BACE1}

For neuroaging specifically, the enzymology is concrete. Against human recombinant MAO-B, baicalin (IC$_{50}$ 6.262\,$\mu$M) and baicalein (21.01\,$\mu$M) are inhibitors; wogonin is a potent MAO-A inhibitor (IC$_{50}$ 6.35\,$\mu$M, K$_i$ 0.56\,$\mu$M) with MAO-B inhibition at 20.8\,$\mu$M in a second assay system --- the inter-assay discrepancy is acknowledged in the primary papers and reproduced here rather than smoothed.\src{185}\src{196} In LPS-inflamed mice, baicalin and baicalein inhibited astrocytic MAO-B-driven GABA synthesis and reversed aberrant tonic neuronal suppression --- MAO-B engagement demonstrable \emph{in vivo}.\src{184} Baicalein inhibits BACE1 (IC$_{50}$ 23.71 $\pm$ 1.91\,$\mu$M), and in-silico target fishing ranked MAO-B and COMT among its top two targets in a parkinsonian framework.\src{166}\src{171} Amyloid models (not tau models --- none retrieved) round out the preclinical neuro-aging file.\src{179}

\subsection{The caloric-restriction-mimetic question}

Baicalein has been explicitly proposed as a caloric-restriction-mimetic candidate --- the conclusion of a natural-aging Wistar-rat study in which it improved pro-oxidant/antioxidant balance and membrane-transporter changes ``similar to CR.''\src{169} The honest reading: the CRM candidacy rests on redox endpoints, not on a metformin-style AMPK/SIRT1 pathway dossier, and this paper holds that label as a hypothesis rather than awarding it.\src{169}

% body2.tex — sections 6–10 of the Scutellaria longevity review.
% \input from paper/scutellaria_review.tex; compiled together.

\section{Senescence, Telomeres, and the Aging Organism}
\label{sec:senes}

\subsection{Senomorphic, not senolytic}

The senescence literature splits cleanly by context, and the split is the finding. In normal stressed cells, the constituents act as senescence-delays and SASP suppressors. Baicalein protected human and rat astrocytes from H$_2$O$_2$-induced senescence, suppressing the SASP secretome (IL-6, IL-8, TNF-$\alpha$, CXCL1, MMP-1) through JAK2/STAT1/NF-$\kappa$B inhibition --- the cleanest senescence-delay mechanism paper in the corpus.\src{213} In bleomycin-induced lung fibrosis, baicalein (100 mg/kg/day oral) attenuated fibroblast senescence markers (p53, p21, p16, MCP-1, PAI-1) and cut collagen deposition from 27.3\% to 4.1\% of lung area through SIRT3 restoration, with Sirt3 siRNA abolishing the effect.\src{177} Wogonin protected against the same fibrosis by inhibiting CDK9/p53-mediated senescence, and separately attenuated chondrocyte senescence in osteoarthritic mice with STAT3 as its direct target.\src{221}\src{227} Whole-plant extract delayed low-glucose-induced astrocyte senescence with the same SASP-marker signature.\src{215}

The most useful twist for a formulary: the strongest SASP inhibitors in the root are \emph{not} the three headline flavones. A screen of seventeen flavonoids from \textit{Scutellaria} radix found five potent SASP inhibitors --- oroxylin A and 2$'$-hydroxylated flavones --- acting through I$\kappa$B$\zeta$/C/EBP$\beta$ without touching the cell-cycle arrest itself; textbook senomorphic action.\src{211} One of them, 2$'$,5,7-trihydroxyflavone, suppressed SASP and I$\kappa$B$\zeta$ gene expression in the kidneys of 21-month-old rats at 2--4 mg/kg/day oral --- an aged-animal SASP result, not a young-rat result.\src{211} For a clinician deciding between an isolated baicalein product and a whole-root extract standardized to its full flavone profile, this is the strongest evidence-based argument for the extract.\src{211}

What does \emph{not} exist matters equally. No experiment demonstrates senolytic activity --- selective killing of senescent cells --- for any \textit{Scutellaria} constituent; baicalein appears on one in-silico candidate list whose authors explicitly call for validation.\src{217} And in activated stellate cells and cancer cells the same molecules \emph{induce} senescence as therapy: baicalein drove stellate-cell senescence via CEBPZ/p53-mediated glycolysis inhibition to attenuate liver fibrosis; baicalin induced colon-cancer senescence via DEPP/Ras-Raf-MEK-ERK with p16/Rb activation, and by telomerase-reverse-transcriptase inhibition in parallel work --- while in other colon-cancer lines baicalin \emph{suppressed} CDKN2A protein.\src{225}\src{206}\src{208}\src{220} The pro- versus anti-senescence direction is cell-state-dependent, exactly as a gene-expression rheostat should behave; ``\textit{Scutellaria} fights senescence'' is the oversimplification this section exists to prevent.

\subsection{Telomeres}

The telomere file is oncology-only: wogonin and aqueous root extract inhibited hTERT, hTP1, and c-myc mRNA and down-regulated telomerase activity in HL-60 leukemia cells --- and that same paper documented the $\beta$-glucuronidase conversion of baicalin and wogonoside to their aglycones that oral dosing depends on.\src{199} No study of any constituent examines telomere-length maintenance in normal or aging cells; that is a gap, stated, not a silence.\src{199}

\subsection{Organismal aging models}

In SAMP8 senescence-accelerated mice, baicalein reduced the senescence grading score, improved spatial learning, object recognition, and olfactory memory, suppressed cortical IL-6, IL-1$\beta$, and TNF-$\alpha$ with correlated microbiome shifts, and in the follow-up transcriptomic study prevented cortical JAK2/STAT1 activation while lowering A$\beta$1-42 and RAGE.\src{207}\src{124} In D-galactose aging models baicalein improved novel-object recognition and water-maze performance while inhibiting proinflammatory mediators, and a standardized \textit{S. baicalensis}--\textit{Acacia catechu} composition protected mice against induced immunosenescence.\src{204}\src{222} In naturally aged Wistar rats, baicalein improved pro-oxidant/antioxidant profiles, membrane-transporter function, and eryptosis --- the basis of its CRM candidacy.\src{169}

The organ-level aging models add texture. Skin: oral baicalein suppressed UVB-induced skin thickening and tumor genesis across 23 weeks of thrice-weekly irradiation in hairless mice; topical baicalin reduced epidermal thickening and raised collagen I/III while lowering UVB-induced senescence markers in human dermal fibroblasts; and whole extract up-regulated the clock gene BMAL1 against photoaging through REV-ERB$\alpha$ antagonism.\src{224}\src{201}\src{216} Joints: baicalin reduced IL-1$\beta$-induced catabolic and inflammatory markers in human chondrocytes with in-vivo confirmation.\src{205} Bone sounds a caution: oral baicalein at 10 mg/kg for 13 weeks showed no musculoskeletal benefit in ovariectomized rats --- the earlier positive bone findings used subcutaneous dosing, so route and dose decide the outcome.\src{214}

Two boundaries close the section. No mammalian lifespan study exists for any constituent --- rodent work ends at biomarkers, senescence-accelerated strains, and induced-aging models; the only organismal lifespan data are the invertebrate results of section~\ref{sec:pathways}.\src{167}\src{226} And the epigenetic--senescence bridge remains unclosed in a single experiment: DNMT1/SCARA5, AHCY/H3K4me3, TET2/5hmC, and LSD1/p65 are all demonstrated --- but in fibrosis, cancer, and osteoarthritis contexts, never inside a senescence or aging model; no study tests p16/CDKN2A promoter methylation after any \textit{Scutellaria} exposure.\src{157}\src{156}\src{223}\src{155}

\section{Human Evidence: Trials, Pharmacokinetics, the Microbiome}
\label{sec:human}

\subsection{The ceiling, stated first}

No human study of \textit{S. baicalensis}, baicalin, baicalein, or any blend has measured a longevity endpoint --- not mortality, not healthspan, not a geroscience biomarker panel. The longest human exposures on record are a 444-day uncontrolled liver-function monitoring study in 17 patients taking 1{,}335 mg dried root daily, a one-year open-label flavocoxid study in 17 Duchenne muscular dystrophy boys, and a six-month negative multi-herb asthma trial in children; every controlled efficacy trial runs twelve weeks or less.\src{249}\src{239}\src{237} That is the honest ceiling on every longevity claim this paper makes, and it does not move.

\subsection{The flavocoxid program: the best-powered human evidence, with an asterisk}

Flavocoxid (Limbrel) is a proprietary blend of baicalin from \textit{S. baicalensis} root with catechins from \textit{Acacia catechu}, standardized above 90\% --- the closest thing this plant has to a pharmaceutical development program, and every skullcap osteoarthritis efficacy trial used it.\src{198}\src{240} The program reads well on its face: a one-month double-blind pilot (n = 103) found flavocoxid 500 mg twice daily as effective as naproxen 500 mg twice daily for knee osteoarthritis; the flagship 220-subject, multicenter, noninferiority RCT confirmed no efficacy difference from naproxen with significantly fewer upper-gastrointestinal and edema adverse events; an independent placebo-controlled safety RCT (n = 59, 12 weeks) found no laboratory or blood-pressure signal; a dedicated coagulation study found no thromboxane, platelet, or bleeding-time effects and no potentiation of warfarin in warfarinized patients; a 1{,}067-patient open-label post-marketing study improved physician global assessments; and the one-year Duchenne study added a pediatric exposure datapoint.\src{198}\src{241}\src{238}\src{242}\src{244}\src{239}

The asterisk is the provenance and the injury signal, and a clinician should know both. Several of these trials share manufacturer-affiliated authors; the flagship was powered for noninferiority against an active comparator, not superiority over placebo; and the post-marketing study was uncontrolled.\src{241}\src{244} Against that, the Drug-Induced Liver Injury Network prospectively captured four adults with acute hepatocellular injury beginning one to three months after starting flavocoxid (mean peak ALT 1{,}268 U/L, all recovering after cessation) --- the signal that would eventually pull the product from the US market (section~\ref{sec:safety}).\src{240} A claims-data cohort of 3{,}337 flavocoxid initiators matched to 6{,}674 NSAID users put hospitalized liver injury at 3.2 versus 2.4 per 1{,}000 person-years (adjusted HR 1.16, 95\% CI 0.23--6.01) --- rare and statistically indistinguishable from NSAIDs --- with a significantly \emph{lower} gastrointestinal-bleeding rate.\src{210}

The smaller blend trials round out the record with an honest negative: the same species pair improved post-influenza-vaccination IgA and IgG responses in 50 adults aged 40--80 over 56 days, and relieved acute joint discomfort over one week versus naproxen in 79 adults --- but a four-week, 85-subject cognition trial was negative, with time-related gains in both arms and no advantage over placebo.\src{87}\src{255}\src{86}

\subsection{Single-herb and constituent trials}

The single-herb human file is small, recent, and mostly Asian. In a crossover trial, patients with type 2 diabetes already on metformin added 3.52 g/day of \textit{S. baicalensis} for eight weeks: glucose tolerance and TNF-$\alpha$ improved, and gut \textit{Lactobacillus} and \textit{Akkermansia} --- a genus the longevity field watches closely --- increased remarkably; some subjects developed elevated liver enzymes, a flag worth remembering.\src{195} A 374-patient, 12-week trial of purified baicalin 500 mg/day in coexisting coronary artery disease and rheumatoid arthritis improved triglycerides, LDL, apolipoproteins, cardiotrophin-1, and hs-CRP with good or moderate EULAR responses in 71\% versus 53\% --- striking numbers that must be held loosely, because the trial is single-center, unregistered, and unreplicated.\src{187} The mood file is genuinely new: a 2025 six-week factorial trial in 180 adults with mild-to-moderate depressive symptoms found a positive scutellaria effect on mood regulation with a synergistic interaction with saffron --- the first human demonstration of a mood effect for this plant --- and a separate 15-day crossover trial of a \textit{S. baicalensis}--\textit{Crataegus} combination improved mood, trait anxiety, and working memory under stress in 43 adults.\src{252}\src{39} Periodontal evidence is practical: a \textit{S. baicalensis} toothpaste reduced experimental gingivitis indices in 40 adults, and a compound gargle improved plaque, probing depth, bleeding, and periodontal pathogens in 110 orthodontic patients.\src{251}\src{248}

\subsection{Pharmacokinetics: the microbiome is the delivery system}

Human pharmacokinetics explain both the promise and the dosing puzzle. Baicalin as such is poorly absorbed; the colonic microbiota's $\beta$-glucuronidases hydrolyze it to baicalein, which is absorbed and immediately re-conjugated --- in an in-vitro colon model, baicalin converted to baicalein only when volunteer-derived microbiota were present, and human fecal flora convert the root's glucuronides to their aglycones.\src{235}\src{245} What circulates is therefore overwhelmingly conjugate: after oral Sho-saiko-to, baicalein 6-O-sulfate was identified in human plasma, with secondary plasma peaks at 12 and 24 hours --- the double-peak signature of enterohepatic recycling --- and after oral Shuanghuanglian, thirteen baicalin-derived components were detected in plasma, with parent baicalin in the low microgram range and free baicaleen measured in nanograms.\src{229}\src{254}\src{83} Urinary recovery is minor --- about 7.2\% of a 5.2 g root dose as baicalein conjugates, sulfates predominating.\src{253}

The purified-constituent program is the strongest clinical-pharmacology asset this plant owns. Baicalein chewable tablets completed single-ascending-dose testing from 100 to 2{,}800 mg in 72 healthy Chinese volunteers --- safe and well tolerated, with median Tmax 0.75--3.5 h, under 1\% urinary clearance, and roughly 27\% fecal recovery --- and multiple-ascending-dose testing at 200--800 mg/day, reaching steady state by day 8 with accumulation indices of 1.66--2.45 and no serious accumulation; a second tablet program reported Cmax of 280--845 ng/mL with t$_{1/2}$ 7.8--14.9 h and less-than-dose-proportional exposure.\src{228}\src{38}\src{37}\src{230} Oroxylin A tablets completed phase I in 2026 (400--2{,}400 mg, dose-proportional Cmax, 1.6-fold food effect).\src{246} (The often-repeated ``Korean'' phase I trial does not exist; the published programs are Chinese, a correction this review carries deliberately.\src{228}\src{37})

Two translation facts finish the pharmacology. In mice, all six root flavones cross the blood--brain barrier --- oroxylin A highest, at 7.9--224 pmol/g brain --- but no human BBB data exist.\src{243} And the nanoformulation/phospholipid literature, however voluminous, remains preclinical: reviews of the field concede that no well-designed human absorption trial of a delivery-optimized \textit{Scutellaria} product has been completed.\src{247}\src{231}

\subsection{The concentration-translation problem, quantified}

The mechanistic sections above rest substantially on 25--100\,$\mu$M in-vitro exposures. Human oral baicalein at 200--800 mg/day yields total plasma concentrations of roughly 280--845 ng/mL ($\sim$1--2\,$\mu$M), mostly conjugated, with free aglycone a fraction of that.\src{37}\src{228} The gap between the epigenetic-enzyme literature and achievable human exposure is therefore real --- one to two orders of magnitude for the aglycone --- and the one datum that keeps this from being a pure dismissal is that at genuine clinical doses these molecules still engage targets measurably: baicalin 50 mg three times daily for 14 days reduced rosuvastatin AUC by 47\% in humans through OATP1B1, so at real-world doses the constituents demonstrably reach and modulate transporter function in the liver.\src{271} Whether epigenetic engagement occurs at oral doses in humans is exactly the study nobody has done (section~\ref{sec:synth}).\src{37}

\begin{table}[ht]
\caption{Principal human trials and exposure studies. Longevity endpoints: none exist. Duration column is the honest ceiling on efficacy claims.}
\label{tab:human}
\footnotesize
\begin{tabular}{>{\raggedright\arraybackslash}p{3.6cm} >{\raggedright\arraybackslash}p{2.6cm} >{\raggedright\arraybackslash}p{2.2cm} >{\raggedright\arraybackslash}p{5.7cm}}
\toprule
Trial / study & Population, n & Duration & Result and caveat \\
\midrule
Flavocoxid vs naproxen, pilot RCT & knee OA, 103 & 1 month & Equal to naproxen on signs and symptoms; active comparator only\src{198} \\
Flavocoxid vs naproxen, multicenter RCT & knee OA, 220 & 12 weeks & Noninferior; fewer upper-GI and edema events; manufacturer-affiliated\src{241} \\
Flavocoxid safety RCT & OA, 59 & 12 weeks & No adverse laboratory or BP signal; independent\src{238} \\
Flavocoxid coagulation study & volunteers and warfarinized OA patients & acute & No platelet/bleeding effect; did not potentiate warfarin\src{242} \\
Flavocoxid post-marketing & 1{,}067 patients & 8 weeks & VAS improved; open-label, uncontrolled\src{244} \\
Flavocoxid in Duchenne dystrophy & 17 boys & 1 year & No dropouts; cytokine falls only with steroid co-therapy\src{239} \\
Vaccine-response RCT & 50 adults, 40--80 y & 8 weeks & Higher post-vaccination IgA/IgG (blend)\src{87} \\
Cognition RCT & 85 healthy adults & 4 weeks & Negative vs placebo (blend)\src{86} \\
Type 2 diabetes crossover & on metformin & 8 weeks & Glucose tolerance, TNF-$\alpha$, microbiome improved; some LFT elevations\src{195} \\
Baicalin in CAD + RA & 374 patients & 12 weeks & Lipids, hs-CRP, EULAR response improved; single-center, unregistered, unreplicated\src{187} \\
Mood trials (2025) & 180 and 43 adults & 6 and 0.5 weeks & First human mood effects; saffron synergy; small trials\src{252}\src{39} \\
Baicalein phase I & 72 $+$ healthy volunteers & acute and 14 days & Safe to 2{,}800 mg; Cmax 280--845 ng/mL at 200--800 mg\src{228}\src{37} \\
Hepatic monitoring & 17 patients & mean 444 days & No liver-test change on 1{,}335 mg dried root/day; uncontrolled\src{249} \\
\bottomrule
\end{tabular}
\end{table}

\section{Safety: The Liver Question and the Interaction Map}
\label{sec:safety}

\subsection{Four separable questions}

``Is skullcap hepatotoxic?'' is four questions wearing one coat, and the review separates them: (1) concentrated blends (flavocoxid/Limbrel); (2) the historical ``skullcap'' injury record, largely American skullcap adulterated with germander; (3) authenticated \textit{S. baicalensis} root itself; and (4) interactions. Each has a different answer.

\subsection{Flavocoxid: a real, rare, idiosyncratic signal}

The flavocoxid DILI record is well characterized because it was captured prospectively. The DILIN series attributed four of 877 enrolled DILI cases to flavocoxid --- all women, onset averaging 11.2 weeks, mean peak ALT 1{,}268 U/L, hepatocellular or mixed pattern, all recovering fully after cessation, no liver failure.\src{240} LiverTox assigns likelihood score C, notes premarketing aminotransferase elevations above 3$\times$ ULN in only 1--2\% of patients, and records no reported deaths.\src{257} The FDA received 194 adverse-event reports for Limbrel products (DILI, pancreatitis, hypersensitivity pneumonitis) and requested a nationwide recall in December 2017; distribution ceased the same month.\src{258}\src{259} The claims cohort's incidence --- 3.2 versus 2.4 per 1{,}000 person-years against NSAIDs, confidence interval spanning unity --- frames the signal as rare rather than epidemic.\src{210}

\subsection{The germander confounding}

The classic ``skullcap hepatotoxicity'' record belongs substantially to adulteration. LiverTox's skullcap monograph states that in most liver-injury cases multiple herbals were involved and that phytochemical analysis has identified germander adulteration or mislabeling in some instances; analytical methods detect \textit{T. chamaedrys} in \textit{S. baicalensis} material down to 0.5\% w/w, and hyperspectral screening was proposed precisely because adulteration-driven hepatotoxicity is a documented market reality.\src{64}\src{82}\src{291}

\subsection{Authenticated root: a small genuine signal inside wide margins}

A small genuine case literature for \textit{S. baicalensis} itself does exist: two Naranjo-probable hepatotoxicity cases from Chinese skullcap within a combination supplement; a case with an inadvertent positive rechallenge; a 51-year-old man developing jaundice ten days after a 400 mg \textit{S. baicalensis} supplement, bilirubin peaking at 24 mg/dL before full resolution; and recurrent DILI on re-administration of two Kampo formulas containing scutellaria root.\src{264}\src{265}\src{266}\src{263} The first population-level estimate appeared in 2026: confirmed \textit{S. baicalensis}-attributed DILI at 0.095\% of 75{,}564 herbal-medicine outpatients across Shanghai hospitals, with risk significantly elevated above the 10 g/day pharmacopoeial maximum (adjusted OR 2.21) and the ABCB1 rs2032582 C/C genotype overrepresented among cases (58.3\% versus 16.3\%) --- a pharmacogenetic handle on an idiosyncratic reaction.\src{261}\src{262}

The counterweights are equally real, and the prescribing question sits between them. Seventeen patients took 1{,}335 mg dried root daily for a mean 444 days with no significant change in any liver test, and the authors noted that every previously published case involved at least one co-supplement with its own hepatotoxicity association; LiverTox's annotation of that series concludes there is no convincing evidence that \textit{S. baicalensis} causes liver injury.\src{249}\src{64} Preclinically the root reads as protective: baicalin after acetaminophen challenge prevented hepatic necrosis and cut mortality from 43\% to 0\% in mice, and inhibited aflatoxin-B1-induced hepatocyte apoptosis in vitro.\src{269}\src{270} The GLP-compliant 26-week rat toxicology of a root extract found no target organ of toxicity with a NOAEL of 2{,}000 mg/kg/day --- roughly 75 g/day of raw herb for a 60-kg adult, seven to twenty-five times the 3--10 g/day pharmacopoeial range --- with an acute minimum lethal dose above 5{,}000 mg/kg; baicalin's own subacute NOAEL is 2{,}000 mg/kg with reversible changes only at 4{,}000.\src{267}\src{45}\src{268} The synthesis that fits all of it: a rare, idiosyncratic, largely extract- and blend-associated, dose- and possibly genotype-modulated signal, inside wide margins for authenticated root at pharmacopoeial doses --- and a 2024 meta-analysis of skullcap-associated HILI, 14.2\% mortality among published cases, that its own editor tempers with publication bias before anyone else has to.\src{293}

\subsection{Interactions: one that matters, several that are theoretical}

Only two direct human interaction studies exist, and one of them matters clinically: baicalin 50 mg three times daily for 14 days reduced rosuvastatin AUC(0--72) by 47.0\% in an OATP1B1-haplotype-dependent manner --- a loss-of-efficacy interaction for the most commonly prescribed statin class member, genotype-dependent, and the single best reason to ask about statins before recommending this plant long-term.\src{271} The single 500 mg baicalin dose study found no clinically significant effect on cyclosporine pharmacokinetics in 16 volunteers, but rodent data are bidirectional (decoction cut oral cyclosporine AUC up to 82\%; multi-dose baicalin lowered bioavailability 25--38\% via intestinal P-glycoprotein induction), so transplant patients on calcineurin inhibitors warrant level monitoring rather than reassurance.\src{232}\src{272}\src{273} The in-vitro enzymology is broad --- baicaleen and wogonin competitively inhibit CYP2C9 (K$_i$ $\leq$ 2.2\,$\mu$M, a plausible warfarin-interaction basis), baicalin inhibits CYP1A2, baicaleen inhibits CYP3A4 (IC$_{50}$ 12\,$\mu$M) and tripled simvastatin AUC in rats --- but no warfarin or statin case reports exist, and the flavocoxid coagulation study was reassuring in humans; the honest status of most of this column is mechanistic caution, not clinical event.\src{275}\src{274}\src{276}\src{242} Pharmacodynamically, all three flavones are antithrombotic in vitro and in rodents (baicalin prolonged in-vivo bleeding time), wogonin and baicaleen inhibit MAO-A and MAO-B at low-micromolar constants, and baicaleen is hypoglycemic in diabetic rodents --- so the counseling points are bleeding additivity with antithrombotics, serotonergic vigilance, and glucose awareness, each low-certainty and worth one sentence in a chart note rather than an alarm.\src{279}\src{278}\src{280}\src{196}\src{282}

\subsection{Special populations and quality}

Pregnancy is where tradition and modern data diverge most visibly: huangqin has millennia of traditional use ``to clear heat, protect the fetus, and avoid a miscarriage,'' while the modern literature is limited to disease models; the developmental toxology shows a modest, reversible fetal skeletal variation (lumbar rib) in rats at roughly 12.5--25 g/kg of water-extract concentrate, resolving postnatally and judged a transient variation at doses a hundred-fold above typical human intake --- and no human pregnancy, lactation, or pediatric safety data exist at all.\src{283}\src{284}\src{285}\src{286} Quality closes the loop opened in section~\ref{sec:identity}: in a 386-herb pharmacopoeial survey, the pyrrolizidine-alkaloid-positive \textit{Scutellaria} species was \textit{S. barbata} (61.7 $\mu$g/kg, margins-of-exposure acceptable at pharmacopoeial dosing) --- not \textit{S. baicalensis}, which the survey did not flag; aflatoxin B1 appeared in two of six commercial samples below national limits; a heavy-metal survey of medicinal species including \textit{S. baicalensis} concluded no chronic health risk.\src{287}\src{289}\src{288}

\section{Synthesis: The Longevity Formulary View}
\label{sec:synth}

The unifying model this evidence supports is compact enough to carry into a consultation. \textit{S. baicalensis} is a gene-expression botanical: its constituents engage the epigenetic machinery (DNMT1, HDAC1/8, p300, LSD1, AHCY, TET2, FTO, miRNA processing) directly; the downstream consequences are program-level shifts in the aging-relevant circuitry (NF-$\kappa$B down, Nrf2 up in non-cancer contexts, AMPK/sirtuin/FOXO modulation, autophagy and mitophagy induction); the organismal phenotype in the one model that measured it was longer life, and the one genetic experiment that asked how found the answer in an epigenetic acetyltransferase.\src{142}\src{131}\src{143}\src{117}\src{156}\src{155}\src{167}\src{109} The chain from epigenome to lifespan is complete --- in a worm, for one constituent.\src{167}\src{109}

What can be said today, in tiers:

\begin{itemize}[leftmargin=1.4em]
\item \textbf{Defensible now:} a two-thousand-year traditional record; a mapped multi-layer epigenetic pharmacology; replicated invertebrate lifespan extension for baicalein; senomorphic SASP suppression strongest from whole-root minor constituents; human safety at pharmacopoeial and phase I doses; one human drug interaction (rosuvastatin) worth active management.\src{101}\src{167}\src{211}\src{228}\src{271}
\item \textbf{Promising but unproven:} mammalian aging-biomarker benefits; mood effects; metabolic and lipid effects; microbiome reshaping including \textit{Akkermansia}; the caloric-restriction-mimetic candidacy --- each with single-study or unreplicated status.\src{207}\src{252}\src{195}\src{169}
\item \textbf{Not established:} any human longevity endpoint; senolytic activity; telomere maintenance in normal cells; mammalian lifespan extension; epigenetic-clock effects (no study has ever looked).\src{217}\src{199}\src{167}
\end{itemize}

Practically: prefer authenticated root or a full-spectrum extract standardized to baicalin against a pharmacopoeial marker over isolated-constituent products --- the SASP data argue for the extract, while the lifespan and phase I data argue for the aglycone.\src{211}\src{167}\src{228} Anchor dosing to the 3--10 g/day decoction range of the pharmacopoeias, which the 2026 pharmacovigilance data identify as the safe side of a real dose--risk gradient.\src{261} Obtain baseline and periodic liver panels for long-term users; ask about statins, anticoagulants, and transplant immunosuppression before starting; treat pregnancy as a data gap despite the traditional endorsement.\src{261}\src{271}\src{283}

And the order of enquiry, cheapest first: (1) a methylation-clock substudy piggybacked on any existing long-term \textit{Scutellaria} trial --- the single highest-value unmeasured endpoint in this literature; (2) senescence-marker endpoints (p16, SASP cytokines) in a flavocoxid-style human RCT; (3) a registered, multi-center replication of the 374-patient baicalin trial; (4) PK-bridged epigenetic sampling --- one DNMT1 or LSD1 readout in human lymphocytes at phase I doses, closing the two-order-of-magnitude question; (5) a mammalian lifespan study of baicalein, which would convert five worm papers into a testable mammalian hypothesis.\src{37}\src{187}

\section{Limitations}
\label{sec:limits}

The evidence-quality skew is the first limitation and it is structural: the epigenetic layer is documented overwhelmingly in oncology and fibrosis models, because that is where epigenetic pharmacology gets funded; this review has labeled the exceptions rather than pretending they are the rule.\src{142}\src{154}\src{177} The concentration-translation gap --- 25--100\,$\mu$M in vitro against $\sim$1--2\,$\mu$M mostly-conjugated plasma at oral human doses --- is unbridged by any PK-linked epigenetic study, in either direction.\src{37} The lifespan evidence is invertebrate-only for one constituent, with contested pathway attribution across laboratories; the one mammalian-relevant direct test of the other two constituents was null.\src{167}\src{226} The human trial record is short (12 weeks or less for all controlled efficacy trials), thin (no monopreparation osteoarthritis or liver-disease RCT), partly industry-authored, and in the case of the largest single-constituent trial, unregistered and unreplicated.\src{241}\src{187} No study of any \textit{Scutellaria} constituent has measured an epigenetic age or DNA-methylation clock --- the central test of this paper's own thesis has never been run.\src{156} Two safety-relevant records could not be fully verified and are carried as open signals rather than facts: a rat study reporting chronic liver toxicity of a huangqin ethanol extract whose PubMed record returns a mismatched abstract, and a ``16 suspected cases'' figure cited secondhand inside the GLP toxicology paper.\src{45} And this review's own method, while paranoid about citation integrity, is a narrative synthesis with parallel streams rather than a formally registered systematic review; recency weighting (2019--2026) deliberately shapes the evidence base, and the traditional-use register is historical data of a kind that cannot be randomized.\src{45}\src{293}

\section*{Sources}

\footnotesize
\noindent
Every numbered citation above resolves through a retrieval ledger built mechanically during this review --- 293 sources fetched, registered, and archived, with page text saved on disk; the subset cited here is listed below with the URL each id points to, and most entries carry verbatim evidence sentences captured from the fetched text. Entries without quoted evidence were registered from retrieved pages whose text did not yield a clean quotable sentence; the ledger, the fetched pages, and the full evidence chain live in the project repository (\texttt{citations/ledger.json}, \texttt{research/*/pages/}).

\vspace{0.5em}
% GENERATED by scripts/gen_sources_block.py -- do not edit by hand.
% 182 cited sources; 181 carry captured verbatim quotes;
% sources without quotes were registered from retrieved pages whose text
% did not yield a clean quotable sentence -- cite them from the fetched page.
\noindent [37] Dong, Ruihua, Li, Lijun, Gao, Hongzhi, Lou, Kun, Luo, Hongmei, Hao, Sheng, Yuan, Jing, Liu, Zeyuan (2021). Safety, tolerability, pharmacokinetics, and food effect of baicalein .... \textit{Journal of ethnopharmacology} :. \href{https://doi.org/10.1016/j.jep.2021.114052}{doi:10.1016/j.jep.2021.114052}. \url{https://pubmed.ncbi.nlm.nih.gov/33753147} (accessed 2026-11-14).
\noindent\hspace*{2em}\small\textit{Evidence:} ``After a single administration of baicalein tablets (200 mg, 400 mg, 600 mg, or 800 mg), Cmax were 280.44, 628.80, 845.20, 489.55 ng/mL; AUC0-∞ were 2035.57, 2939.31, 4494.88, and 3754.43 hng/mL, respectively.''
\noindent\hspace*{2em}\small\textit{Evidence:} ``The exposure of baicalein and its metabolites increased in a less than dose-proportional manner.''
\noindent\hspace*{2em}\small\textit{Evidence:} ``CONCLUSION: Baicalein tablets within the studied dose range were safe and well-tolerated in healthy Chinese subjects with no serious or severe adverse effects.''
\noindent [38] Pang H; Xue W; Shi A; Li M; Li Y; Cao G; Yan B; Dong F; Xiao W; He G; Du G; Hu X; Cheng G (2016). Multiple-Ascending-Dose Pharmacokinetics and Safety .... \textit{Clin Drug Investig} 36(9):713-724. \href{https://doi.org/10.1007/s40261-016-0418-7}{doi:10.1007/s40261-016-0418-7}. \url{https://pubmed.ncbi.nlm.nih.gov/27352310} (accessed 2026-11-14).
\noindent\hspace*{2em}\small\textit{Evidence:} ``CONCLUSION: In dose range of 200-800 mg, multiple-dose oral baicalein administration was safe and well tolerated, dose proportionality was inconclusive, and no serious accumulation of baicalein was observed.''
\noindent\hspace*{2em}\small\textit{Evidence:} ``A steady-state concentration of analytes in plasma was achieved on day 8 after repeated dosing.''
\noindent [39] Dodd, Fiona, Weishaupt, Ramon, Katumba, Philipp K M, Elcoate, Rian, Wightman, Emma (2025). Mesh Terms. \textit{Journal of psychopharmacology (Oxford, England)} :. \href{https://doi.org/10.1177/02698811251381261}{doi:10.1177/02698811251381261}. \url{https://pubmed.ncbi.nlm.nih.gov/41194549} (accessed 2026-11-14).
\noindent\hspace*{2em}\small\textit{Evidence:} ``METHODS: Forty-three participants (35 analysed) aged 18-75 years received the herbal supplement and a placebo for 15 days.''
\noindent\hspace*{2em}\small\textit{Evidence:} ``Chronic supplementation improved mood and anxiety, reducing total mood disturbance, anger/hostility, and Trait anxiety scores.''
\noindent [45] Jungbin Song; Hocheol Kim (2026). Safety Evaluation of Scutellaria baicalensis: Acute and 26-Week Repeated-Dose Oral Toxicity Studies in Rats. \textit{Toxics} 14(10):919. \href{https://doi.org/10.3390/toxics14100919}{doi:10.3390/toxics14100919}. \url{https://pmc.ncbi.nlm.nih.gov/articles/PMC13709419} (accessed 2026-11-14).
\noindent\hspace*{2em}\small\textit{Evidence:} ``Potential safety concerns have also emerged from clinical observations. Sporadic cases of clinically apparent liver injury have been reported following the use of \_S. bai-calensis\_-containing products \textbackslash\{\} [11\textbackslash\{\}]. A recent literature review identified 16 suspected cases of \_S. baicalensis\_-associated drug-induced liver injury;''
\noindent\hspace*{2em}\small\textit{Evidence:} ``This raw-material-equivalent amount is approximately 7.5–25 times the daily raw-material dose of 3–10 g described in the Chinese Pharmacopoeia.''
\noindent [54] Zhao K; Zhang J; Zhou L; Sun Z (2024). Scutellaria baicalensis and its flavonoids in the treatment of ... - PMC. \textit{Front Pharmacol}. \href{https://doi.org/10.3389/fphar.2024.1483785}{doi:10.3389/fphar.2024.1483785}. \url{https://pmc.ncbi.nlm.nih.gov/articles/PMC11625591} (accessed 2026-11-14).
\noindent\hspace*{2em}\small\textit{Evidence:} ``Baicalin (baicalein-7-O-glucuronide, C21H18O11) is metabolized to baicalein in animals ( Kang MJ et al., 2014).''
\noindent\hspace*{2em}\small\textit{Evidence:} ``\#\# 2 Application of \_Scutellaria baicalensis\_ and its flavonoids in digestive system tumors First recorded in the Shennong Ben Cao Jing (Classic of the Materia Medica of the Divine Husbandman) in China in the 2nd century B.C., \_S. baicalensis\_ is said to have the efficacy of “treating all kinds of fever, jaundice, diarrhea, edema, and sores”.''
\noindent\hspace*{2em}\small\textit{Evidence:} ``Baicalin (baicalein-7-O-glucuronide, C21H18O11) is metabolized to baicalein in animals ( Kang MJ et al., 2014). The Chinese Pharmacopoeia sets baicalin content of not less than 8\% as the standard for qualification of herbs.''
\noindent [64] National Institute of Diabetes and Digestive and Kidney Diseases (2020). Skullcap - LiverTox - NCBI Bookshelf. \textit{NCBI Bookshelf, National Library of Medicine} NBK548757 (last update 2020 Mar 28). \url{https://www.ncbi.nlm.nih.gov/books/NBK548757} (accessed 2026-11-14).
\noindent\hspace*{2em}\small\textit{Evidence:} ``Skullcap has been implicated in rare instances of clinically apparent liver injury, although in most cases multiple herbal medications were being taken and the role of skullcap in the hepatic damage was unclear.''
\noindent\hspace*{2em}\small\textit{Evidence:} ``in some instances phytochemical analysis has identified significant adulterants (germander) or mislabeling in cases of suspected skullcap hepatotoxicity.''
\noindent\hspace*{2em}\small\textit{Evidence:} ``Recovery was rapid once the herbal was discontinued, but some cases have resulted in acute liver failure.''
\noindent [72] Awad R; Arnason JT; Trudeau V; Bergeron C; Budzinski JW; Foster BC; Merali Z (2003). Phytochemical and biological analysis of skullcap (Scutellaria lateriflora L.): a medicinal plant with anxiolytic properties.. \textit{Phytomedicine : international journal of phytotherapy and phytopharmacology} 10(8):640-9. \href{https://doi.org/10.1078/0944-7113-00374}{doi:10.1078/0944-7113-00374}. \url{https://pubmed.ncbi.nlm.nih.gov/14692724} (accessed 2026-11-14).
\noindent\hspace*{2em}\small\textit{Evidence:} ``The identification and quantification of the flavonoid, baicalin in a 50\% EtOH extract (40 mg/g) and its aglycone baicalein in a 95\% EtOH extract (33 mg/g), as well as the amino acids GABA in H2O and EtOH extracts (approximately 1.6 mg/g) and glutamine in a H2O extract (31 mg/g), was performed using HPLC.''
\noindent [73] Bai C; Yang J; Cao B; Xue Y; Gao P; Liang H; Li G (2020). Growth years and post-harvest processing methods have critical roles on the contents of medicinal active ingredients of Scutellaria baicalensis.. \textit{Industrial crops and products} 158:112985. \href{https://doi.org/10.1016/j.indcrop.2020.112985}{doi:10.1016/j.indcrop.2020.112985}. \url{https://pubmed.ncbi.nlm.nih.gov/33162677} (accessed 2026-11-14).
\noindent\hspace*{2em}\small\textit{Evidence:} ``Baikal Skullcap (Scutellaria baicalensis Georgi) is a perennial herb in the Lamiaceae family and its dried root is used as a famous traditional Chinese medicine (TCM).''
\noindent\hspace*{2em}\small\textit{Evidence:} ``The contents of eight active ingredients (baicalin, wogonoside, baicalein, wogonin, scutellarin, scutellarein, apigenin, and chrysin) in roots of S. baicalensis of different growth years (ranging from 1 year to 15 years) were estimated using high performance liquid chromatography (HPLC) and further analyzed to determine the optimal harvest period.''
\noindent\hspace*{2em}\small\textit{Evidence:} ``The results showed that the best harvesting period for S. baicalensis should be determined as 2-3 years based on comprehensive evaluation of active ingredient content, annual yield increment, and land use efficiency. The contents of active ingredients including baicalin, wogonoside, baicalein, and wogonin in cortex were significantly higher than those in stele (P ≤ 0.05).''
\noindent [75] Costine B; Zhang M; Chhajed S; Pearson B; Chen S; Nadakuduti SS (2022). Exploring native Scutellaria species provides insight into differential accumulation of flavones with medicinal properties.. \textit{Scientific reports} 12(1):13201. \href{https://doi.org/10.1038/s41598-022-17586-1}{doi:10.1038/s41598-022-17586-1}. \url{https://pubmed.ncbi.nlm.nih.gov/35915209} (accessed 2026-11-14).
\noindent\hspace*{2em}\small\textit{Evidence:} ``4'-deoxyflavones accumulated in relatively high concentrations in root tissues compared to aerial tissues in all species except S. lateriflora. Baicalin, the most abundant 4'-deoxyflavone detected, was localized in the roots of S. baicalensis and leaves of S. lateriflora, indicating differential accumulation patterns be''
\noindent\hspace*{2em}\small\textit{Evidence:} ``Collectively, results presented in this study suggest an evolutionary divergence of flavonoid metabolic pathway in the Scutellaria genus of Lamiaceae.''
\noindent [76] Scutellaria baicalensis in Flora of China. \textit{Flora of China @ efloras.org (Missouri Botanical Garden)}. \url{https://www.efloras.org/florataxon.aspx?flora_id=2&taxon_id=200020285} (accessed 2026-11-14).
\noindent\hspace*{2em}\small\textit{Evidence:} ``Scutellaria baicalensis Georgi, Bemerk. Reise Russ. Reich. 1: 223. 1775. 黄芩 huang qin''
\noindent [79] Chen H; Wei Y; Song J; Yang Y; Chen Y; Sun J; Bai D; Sun Z (2025). Preparation, immunological and pharmacological effects of flavonoids in Scutellariae radix: a review.. \textit{Front Pharmacol}. \href{https://doi.org/10.3389/fphar.2025.1732322}{doi:10.3389/fphar.2025.1732322}. \url{https://pubmed.ncbi.nlm.nih.gov/41684516} (accessed 2026-11-14).
\noindent\hspace*{2em}\small\textit{Evidence:} ``Scutellariae radix\textbackslash\{\}] (SR) is bitter and cold in nature. It enters the lung, gallbladder, spleen, large intestine, and small intestine meridians. It clears heat and dries dampness, purges fire and detoxifies, stops bleeding, and stabilizes pregnancy. It excels at clearing lung fire and upper-body heat.''
\noindent\hspace*{2em}\small\textit{Evidence:} ``SR is officially recognized in the \_Pharmacopoeia of the People’s Republic of China\_ ( \_Chinese Pharmacopoeia\_), as well as in the \_European Pharmacopoeia\_ and the \_British Pharmacopoeia\_''
\noindent [80] Gafner S; Bergeron C; Batcha LL; Angerhofer CK; Sudberg S; Sudberg EM; Guinaudeau H; Gauthier R (2003). Analysis of Scutellaria lateriflora and its adulterants Teucrium canadense and Teucrium chamaedrys by LC-UV/MS, TLC, and digital photomicroscopy.. \textit{Journal of AOAC International} 86(3):453-60. \href{https://pubmed.ncbi.nlm.nih.gov/12852558/}{PMID:12852558}. \url{https://pubmed.ncbi.nlm.nih.gov/12852558} (accessed 2026-11-14).
\noindent\hspace*{2em}\small\textit{Evidence:} ``Methods using liquid chromatography with UV detection (LC-UV), thin-layer chromatography (TLC), and digital photomicroscopy were developed to distinguish between the different species of Scutellaria lateriflora L. and its adulterants Teucrium canadense L. and T. chamaedrys L. Chemically, the 70\% ethanol extract of S. lateriflora is characterized by the presence of flavonoids--predominantly baicalin, lateriflorin, dihydrobaicalin, and baicalein.''
\noindent [81] Zhao Q; Chen XY; Martin C (2016). Scutellaria baicalensis, the golden herb from the garden of Chinese medicinal plants.. \textit{Sci Bull (Beijing)}. \href{https://doi.org/10.1007/s11434-016-1136-5}{doi:10.1007/s11434-016-1136-5}. \url{https://pubmed.ncbi.nlm.nih.gov/27730005} (accessed 2026-11-14).
\noindent\hspace*{2em}\small\textit{Evidence:} ``Huang-Qin is an important ingredient of Xiaochai Hutang (Chinese) or Sho-saiko-to (SST, Japanese) preparations, first described in \_Shanghan Lun\_ (On Cold Damage), written by Zhang Zhongjing around 200 AD \textbackslash\{\} [9\textbackslash\{\}].''
\noindent\hspace*{2em}\small\textit{Evidence:} ``More than 30 types of flavones can be found in its roots (Fig. 3), including baicalin, baicalein, chrysin, oroxylin A, oroxylin A 7-O-glucuronide, wogonin and wogonoside''
\noindent\hspace*{2em}\small\textit{Evidence:} ``Xiaochai Hutang | \_Scutellaria baicalensis\_, \_Bupleurum falcatum\_, \_Pinellia ternate\_, \_Panax ginseng\_, \_Glycyrrhiza uralensis\_, \_Zingiber officinale\_, \_Ziziphus jujuba\_''
\noindent [82] Hong T; Jeong ML; Zahn M; Fay BA; Lee K; Hwangbo H; Park E; Kim M et al. (2009). Detection of the potential adulterant Teucrium chamaedrys in Scutellaria baicalensis raw material and extract by high-performance thin-layer chromatography.. \textit{Journal of AOAC International} 92(3):785-8. \href{https://pubmed.ncbi.nlm.nih.gov/19610368/}{PMID:19610368}. \url{https://pubmed.ncbi.nlm.nih.gov/19610368} (accessed 2026-11-14).
\noindent\hspace*{2em}\small\textit{Evidence:} ``Teucrium chamaedrys (Gemander) has been reported as an adulterant of Scutellaria lateriflora (American skullcap) herbal preparations and is also known to be hepatotoxic. A quick and simple high-performance thin-layer chromatographic (HPTLC) method was developed for the detection of T. chamaedrys (Germander) in S. baicalensis (Chinese skullcap) extract''
\noindent\hspace*{2em}\small\textit{Evidence:} ``The method is sensitive and can detect T. chamaedrys at levels as low as 0.5\% (w/w).''
\noindent [83] Huang T; Liu Y; Zhang C (2019). Pharmacokinetics and Bioavailability Enhancement of Baicalin: A Review.. \textit{European journal of drug metabolism and pharmacokinetics} 44(2):159-168. \href{https://doi.org/10.1007/s13318-018-0509-3}{doi:10.1007/s13318-018-0509-3}. \url{https://pubmed.ncbi.nlm.nih.gov/30209794} (accessed 2026-11-14).
\noindent\hspace*{2em}\small\textit{Evidence:} ``Baicalin displays a distinct pharmacokinetic profile including gastrointestinal hydrolysis, enterohepatic recycling, carrier-mediated transport, and complicated metabolism.''
\noindent [84] Jiang Y; Zhu C; Wang S; Wang F; Sun Z (2023). Identification of three cultivated varieties of Scutellaria baicalensis using the complete chloroplast genome as a super-barcode.. \textit{Scientific reports} 13(1):5602. \href{https://doi.org/10.1038/s41598-023-32493-9}{doi:10.1038/s41598-023-32493-9}. \url{https://pubmed.ncbi.nlm.nih.gov/37019975} (accessed 2026-11-14).
\noindent\hspace*{2em}\small\textit{Evidence:} ``Overall, this study showed that the complete chloroplast genome could be used as a super-barcode to identify three cultivated varieties of S. baicalensis and provide biological information, and it also contributes to bioprospecting.''
\noindent\hspace*{2em}\small\textit{Evidence:} ``The three new varieties cultivated could not be distinguished by morphology before flowering.''
\noindent [86] Krieger JM; Hagele AM; Mumford PW; De Gregorio D; Gaige CJ; Hoffman ER; Gross KN; Holley KF et al. (2025). Effects of Acacia catechu and Scutellaria baicalensis extract on cognitive function in a healthy adult population: a double-blind, randomized, placebo-controlled trial.. \textit{Nutritional neuroscience} 28(8):975-986. \href{https://doi.org/10.1080/1028415X.2024.2442425}{doi:10.1080/1028415X.2024.2442425}. \url{https://pubmed.ncbi.nlm.nih.gov/39701593} (accessed 2026-11-14).
\noindent\hspace*{2em}\small\textit{Evidence:} ``Four-week daily TP supplementation significantly enhanced cognitive function without difference from placebo.''
\noindent\hspace*{2em}\small\textit{Evidence:} ``DISCUSSION: Four-week daily TP supplementation significantly enhanced cognitive function without difference from placebo.''
\noindent [87] Lewis ED; Crowley DC; Guthrie N; Evans M (2023). Role of Acacia catechu and Scutellaria baicalensis in Enhancing Immune Function Following Influenza Vaccination of Healthy Adults: A Randomized, Triple-Blind, Placebo-Controlled Clinical Trial.. \textit{Journal of the American Nutrition Association} 42(7):678-690. \href{https://doi.org/10.1080/27697061.2022.2145525}{doi:10.1080/27697061.2022.2145525}. \url{https://pubmed.ncbi.nlm.nih.gov/36413261} (accessed 2026-11-14).
\noindent\hspace*{2em}\small\textit{Evidence:} ``In the post-vaccination period, total IgA and IgG levels increased in participants supplemented with UP446 vs. those on Placebo (p ≤ 0.026).''
\noindent\hspace*{2em}\small\textit{Evidence:} ``Fifty healthy adults 40-80 years of age who had not received their flu vaccine were randomized to either UP446 or Placebo.''
\noindent\hspace*{2em}\small\textit{Evidence:} ``RESULTS: In the post-vaccination period, total IgA and IgG levels increased in participants supplemented with UP446 vs.''
\noindent [88] Li C; Lin G; Zuo Z (2011). Pharmacological effects and pharmacokinetics properties of Radix Scutellariae and its bioactive flavones.. \textit{Biopharmaceutics \& drug disposition} 32(8):427-45. \href{https://doi.org/10.1002/bdd.771}{doi:10.1002/bdd.771}. \url{https://pubmed.ncbi.nlm.nih.gov/21928297} (accessed 2026-11-14).
\noindent\hspace*{2em}\small\textit{Evidence:} ``Six flavones are proven to be the major bioactive flavones in Radix Scutellariae existing in the forms of aglycones (baicalein, wogonin, oroxylin A) and glycosides (baicalin, wogonoside, oroxylin A-7-glucuronide).''
\noindent [89] Li Z; Duan B; Zhou Z; Fang H; Yang M; Xia C; Zhou Y; Wang J (2024). Comparative analysis of medicinal plants Scutellaria baicalensis and common adulterants based on chloroplast genome sequencing.. \textit{BMC genomics} 25(1):39. \href{https://doi.org/10.1186/s12864-023-09920-2}{doi:10.1186/s12864-023-09920-2}. \url{https://pubmed.ncbi.nlm.nih.gov/38191291} (accessed 2026-11-14).
\noindent\hspace*{2em}\small\textit{Evidence:} ``Four highly variable regions (matK-rps16, ndhC-trnV-UAC, psbE-petL, and rps16-trnQ-UUG) may function as potential molecular markers for differentiating S. baicalensis from its adulterants.''
\noindent\hspace*{2em}\small\textit{Evidence:} ``Scutellaria baicalensis Georgi has been extensively used as a medicinal herb in China for over 2000 years.''
\noindent [92] Lu F; Zhang WJ; Zhai S; Sun YW; Chen QJ; Yang XL; Zhang CF; Wang CZ et al. (2020). Anti-migraine effect of wine-processed Radix scutellariae: Pharmacodynamic verification in nitroglycerin-induced rats and correlation study between compounds dissolution and the fractal dimension.. \textit{Journal of ethnopharmacology} 263:113131. \href{https://doi.org/10.1016/j.jep.2020.113131}{doi:10.1016/j.jep.2020.113131}. \url{https://pubmed.ncbi.nlm.nih.gov/32730879} (accessed 2026-11-14).
\noindent\hspace*{2em}\small\textit{Evidence:} ``The data indicated that wine-processing would affect the dissolution of compounds by changing the pore structure of RS.''
\noindent\hspace*{2em}\small\textit{Evidence:} ``The compounds dissolution of total flavonoids and five major components containing baicalein, baicalin, scutellarin, wogonin and wogonoside was determined by UV-Vis spectrophotometry and HPLC separately.''
\noindent\hspace*{2em}\small\textit{Evidence:} ``wine-processed RS could significantly reduce the head-scratching times of NTG-induced rat models (p < 0.01) and markedly decrease the levels of NO (p < 0.01), CGRP (p < 0.05) and the expression of c-Fos (p < 0.01) compared with model group.''
\noindent [93] Ma W; Liu T; Ogaji OD; Li J; Du K; Chang Y (2024). Recent advances in Scutellariae radix: A comprehensive review on ethnobotanical uses, processing, phytochemistry, pharmacological effects, quality control and influence factors of biosynthesis.. \textit{Heliyon} 10(16):e36146. \href{https://doi.org/10.1016/j.heliyon.2024.e36146}{doi:10.1016/j.heliyon.2024.e36146}. \url{https://pubmed.ncbi.nlm.nih.gov/39262990} (accessed 2026-11-14).
\noindent\hspace*{2em}\small\textit{Evidence:} ``As early as the Eastern Han Dynasty, there were many prescriptions involving SR recorded in \_Shanghan Zabing Lun\_, such as Gegenqinlian decoction, Huangqin decoction, Huanglian Ejiao decoction, and Banxia Xiexin decoction \textbackslash\{\} [15\textbackslash\{\}].''
\noindent [94] Makino T; Hishida A; Goda Y; Mizukami H (2008). Comparison of the major flavonoid content of S. baicalensis, S. lateriflora, and their commercial products.. \textit{Journal of natural medicines} 62(3):294-9. \href{https://doi.org/10.1007/s11418-008-0230-7}{doi:10.1007/s11418-008-0230-7}. \url{https://pubmed.ncbi.nlm.nih.gov/18404307} (accessed 2026-11-14).
\noindent\hspace*{2em}\small\textit{Evidence:} ``These flavonoids were found in the root of S. baicalensis; its aerial parts, however, did not contain them. On the other hand, the amounts of those flavonoids in the aerial parts of S. lateriflora were larger than in the root. Herbal tea and dietary supplements of S. lateriflora obtained commercially also contained those flavonoids, and the dietary supplements contained amounts of them comparable with that in kampo medicine.''
\noindent\hspace*{2em}\small\textit{Evidence:} ``Safety Bureau, Ministry of Health Labour and Welfare of Japan, the roots of Scutellaria baicalensis (Chinese skullcap) and S. lateriflora (skullcap) are classified as "the raw materials exclusively used as pharmaceuticals", but their aerial parts are classified as "non-pharmaceuticals" so, in principle, there are no health claims for these materials and no descriptions of drug-like dosages or administration directions.''
\noindent\hspace*{2em}\small\textit{Evidence:} ``s. Dried root of S. baicalensis is also registered in Japanese Pharmacopoeia XV as scutellaria root''
\noindent [98] Scutellaria baicalensis Georgi | Plants of the World Online | Kew Science. \textit{Kew Science - Plants of the World Online}. \url{https://powo.science.kew.org/taxon/458155-1} (accessed 2026-11-14).
\noindent\hspace*{2em}\small\textit{Evidence:} ``First published in Bemerk. Reise Russ. Reich 2: 223 (1775) This species is accepted The native range of this species is Siberia to Russian Far East and Vietnam. It is a perennial and grows primarily in the temperate biome.''
\noindent [99] Song JW; Long JY; Xie L; Zhang LL; Xie QX; Chen HJ; Deng M; Li XF (2020). Applications, phytochemistry, pharmacological effects, pharmacokinetics, toxicity of Scutellaria baicalensis Georgi. and its probably potential therapeutic effects on COVID-19: a review.. \textit{Chinese medicine} 15:102. \href{https://doi.org/10.1186/s13020-020-00384-0}{doi:10.1186/s13020-020-00384-0}. \url{https://pubmed.ncbi.nlm.nih.gov/32994803} (accessed 2026-11-14).
\noindent\hspace*{2em}\small\textit{Evidence:} ``Skullcapflavone I (5,2′-Dihydroxy-7,8-dimethoxyflavone) | C17H14O6 | 314 | Root<br>Hairy Root''
\noindent\hspace*{2em}\small\textit{Evidence:} ``Skullcapflavone II (5,6′-Dihydroxy-6,7,8,2′-tetramethoxyflavone) | C19H18O8 | 374 | Root<br>Hairy Root''
\noindent\hspace*{2em}\small\textit{Evidence:} ``The best harvest time of SB is determined by the highest content of baicalin, wogonoside, baicalein and scutellarin, most of which are concentrated in September''
\noindent [100] Shi YH; Huang YF; Wang WY; Yang L; Zhou H; Sang Z (2020). Analysis on the current quality standards of Chinese materia Medica used in COVID-19 prevention and treatment.. \textit{Pharmacol Res}. \href{https://doi.org/10.1016/j.phrs.2020.105074}{doi:10.1016/j.phrs.2020.105074}. \url{https://pubmed.ncbi.nlm.nih.gov/32653649} (accessed 2026-11-14).
\noindent\hspace*{2em}\small\textit{Evidence:} ``Assay | Baicalin: ≥9.0 \% | Baicalin: ≥10.0 \% | ≥10.0 \% (total of Baicalin, baicalein and wogonin) | Baicalin: ≥9.0 \% | ≥11.0 \% (total of Baicalin and wogonin-7-O-glucuronide);<br>≥3.50 \% (total of Baicalin and wogonin) | Baicalin: ≥9.0 \% | Baicalin: ≥8.0 \%''
\noindent\hspace*{2em}\small\textit{Evidence:} ``The single marker compound of baicalin is used for quantitative analysis in ChP, EP, JP, ISO, and WHO standards, while apart from baicalin, multi-marker compounds are referred in KP and USP.''
\noindent\hspace*{2em}\small\textit{Evidence:} ``Heavy metals | — | √ | √ | √ | √ | √ | √''
\noindent [101] Wang ZL; Wang S; Kuang Y; Hu ZM; Qiao X; Ye M (2018). A comprehensive review on phytochemistry, pharmacology, and flavonoid biosynthesis of Scutellaria baicalensis.. \textit{Pharmaceutical biology} 56(1):465-484. \href{https://doi.org/10.1080/13880209.2018.1492620}{doi:10.1080/13880209.2018.1492620}. \url{https://pubmed.ncbi.nlm.nih.gov/31070530} (accessed 2026-11-14).
\noindent\hspace*{2em}\small\textit{Evidence:} ``Huang-Qin was first recorded in \_Shennong’s Classic of Materia Medica\_ (Shennong Bencao Jing in Chinese) in around 200 AD.''
\noindent\hspace*{2em}\small\textit{Evidence:} ``Chinese Pharmacopoeia requires the content of baicalin in Huang-Qin should be no less than 9\% (Chinese Pharmacopoeia Commission 2015).''
\noindent [102] Ye F; Wang H; Jiang S; Wu J; Shao J; Cheng X; Tu Y; Zhang DY (2004). Quality evaluation of commercial extracts of Scutellaria baicalensis.. \textit{Nutrition and cancer} 49(2):217-22. \href{https://doi.org/10.1207/s15327914nc4902\_14}{doi:10.1207/s15327914nc4902\_14}. \url{https://pubmed.ncbi.nlm.nih.gov/15489215} (accessed 2026-11-14).
\noindent\hspace*{2em}\small\textit{Evidence:} ``However, significant product-to-product and batch-to-batch variation of these marker components was observed with 4 products containing no baicalin at all and baicalein concentration ranging from 0 to 52.3 g/mg. The 50\% growth inhibition concentration of the extracts ranged from 0.18 to 2.0 mg/ml, more than an 11-fold variation. PGE2 levels varied from 19.5 to 111.1 pg/106 cells, more than a 5.7-fold difference.''
\noindent [104] Zheng Y; Zhou S; Zhang H; Lu Z; Deng R; Feng Y; Liu P (2023). Comparative Study of the Flavonoid Content in Radix Scutellaria from Different Cultivation Areas in China.. \textit{International journal of analytical chemistry} 2023:3754549. \href{https://doi.org/10.1155/2023/3754549}{doi:10.1155/2023/3754549}. \url{https://pubmed.ncbi.nlm.nih.gov/36820244} (accessed 2026-11-14).
\noindent\hspace*{2em}\small\textit{Evidence:} ``baicalin, wogonoside, and baicalein were the three main flavonoids with the highest contents in Radix Scutellaria. The content of flavonoids in different cultivated areas was very different, which had significant regionality and was closely related to the natural conditions of various places.''
\noindent [105] Y C Chen; S C Shen; L G Chen; T J Lee; L L Yang (2001). Wogonin, baicalin, and baicalein inhibition of inducible nitric oxide synthase and cyclooxygenase-2 gene expressions induced by nitric oxide synthase inhibitors and lipopolysaccharide.. \textit{Biochemical pharmacology} 61; 1417-27. \href{https://doi.org/10.1016/s0006-2952(01)00594-9}{doi:10.1016/s0006-2952(01)00594-9}. \url{https://pubmed.ncbi.nlm.nih.gov/11331078} (accessed 2026-11-14).
\noindent\hspace*{2em}\small\textit{Evidence:} ``Wogonin, but not baicalin or baicalein, similarly inhibited PGE2 production and COX-2 protein expression in NLA/LPS or L-NAME/LPS-co-treated RAW 264.7 cells.''
\noindent\hspace*{2em}\small\textit{Evidence:} ``The results indicated that these three polyphenolic compounds inhibited LPS-induced NO production in a concentration-dependent manner without a notable cytotoxic effect on these cells.''
\noindent [108] Melissa D Yang; Yi-Ming Chiang; Reiichi Higashiyama; Kinji Asahina; Derek A Mann; Jelena Mann; Clay C C Wang; Hidekazu Tsukamoto (2012). Rosmarinic acid and baicalin epigenetically derepress peroxisomal proliferator-activated receptor γ in hepatic stellate cells for their antifibrotic effect.. \textit{Hepatology (Baltimore, Md.)} 55; 1271-81. \href{https://doi.org/10.1002/hep.24792}{doi:10.1002/hep.24792}. \url{https://pubmed.ncbi.nlm.nih.gov/22095555} (accessed 2026-11-14).
\noindent\hspace*{2em}\small\textit{Evidence:} ``Here we demonstrate YGW prevents and reverses HSC activation by way of epigenetic derepression of Pparγ involving reductions in MeCP2 expression and its recruitment to Pparγ promoter, suppressed expression of PRC2 methyltransferase EZH2, and consequent reduction of H2K27di-methylation at the 3' exon.''
\noindent [109] Alex Lublin; Fumiko Isoda; Harshil Patel; Kelvin Yen; Linda Nguyen; Daher Hajje; Marc Schwartz; Charles Mobbs (2011). FDA-approved drugs that protect mammalian neurons from glucose toxicity slow aging dependent on cbp and protect against proteotoxicity.. \textit{PloS one} 6; e27762. \href{https://doi.org/10.1371/journal.pone.0027762}{doi:10.1371/journal.pone.0027762}. \url{https://pubmed.ncbi.nlm.nih.gov/22114686} (accessed 2026-11-14).
\noindent\hspace*{2em}\small\textit{Evidence:} ``These protective effects were blocked by RNAi inhibition of cbp-1 in adults only, which also blocks protective effects of dietary restriction.''
\noindent\hspace*{2em}\small\textit{Evidence:} ``Subsequent screening demonstrated that 6 of these 30 drugs increase lifespan in C. elegans: caffeine, ciclopirox olamine, tannic acid, acetaminophen, bacitracin, and baicalein.''
\noindent [110] Si Qin; Jihua Chen; Shunsuke Tanigawa; De-Xing Hou (2012). Gene expression profiling and pathway network analysis of hepatic metabolic enzymes targeted by baicalein.. \textit{Journal of ethnopharmacology} 140; 131-40. \href{https://doi.org/10.1016/j.jep.2011.12.046}{doi:10.1016/j.jep.2011.12.046}. \url{https://pubmed.ncbi.nlm.nih.gov/22265932} (accessed 2026-11-14).
\noindent\hspace*{2em}\small\textit{Evidence:} ``Moreover, the ARE-reporter gene assay demonstrated that baicalein stimulated Nrf2-mediated ARE transactivation.''
\noindent [111] Yan Gao; Jia Lu; Yu Zhang; Yafen Chen; Zhenlun Gu; Xiaogang Jiang (2013). Baicalein attenuates bleomycin-induced pulmonary fibrosis in rats through inhibition of miR-21.. \textit{Pulmonary pharmacology \& therapeutics} 26; 649-54. \href{https://doi.org/10.1016/j.pupt.2013.03.006}{doi:10.1016/j.pupt.2013.03.006}. \url{https://pubmed.ncbi.nlm.nih.gov/23523661} (accessed 2026-11-14).
\noindent\hspace*{2em}\small\textit{Evidence:} ``Baicalein remarkably decreased miR-21 levels and inhibited the increased expression of TGF-β1 and p-Smad-2/3 in bleomycin-treated rats.''
\noindent [113] Xin-mei Chen; Yang Bai; Yu-jian Zhong; Xiao-lin Xie; Han-wu Long; Yu-yin Yang; Shi-gen Wu; Qiang Jia (2013). Wogonin has multiple anti-cancer effects by regulating c-Myc/SKP2/Fbw7α and HDAC1/HDAC2 pathways and inducing apoptosis in human lung adenocarcinoma cell line A549.. \textit{PloS one} 8; e79201. \href{https://doi.org/10.1371/journal.pone.0079201}{doi:10.1371/journal.pone.0079201}. \url{https://pubmed.ncbi.nlm.nih.gov/24265759} (accessed 2026-11-14).
\noindent\hspace*{2em}\small\textit{Evidence:} ``The present results suggest that wogonin has multiple anti-cancer effects associated with degradation of c-Myc, SKP2, HDAC1 and HDAC2.''
\noindent [115] Yan Zhang; Guangtao Ma; Chenlong Li; Zhenyu Cao; Fuzhong Qie; Xinmiao Xu (2015). Baicaleininhibits VSMCs proliferation via regulating LncRNAAK021954 gene expression.. \textit{International journal of clinical and experimental medicine} 8; 22129-38. \href{https://pubmed.ncbi.nlm.nih.gov/26885187/}{PMID:26885187}. \url{https://pubmed.ncbi.nlm.nih.gov/26885187} (accessed 2026-11-14).
\noindent\hspace*{2em}\small\textit{Evidence:} ``Microarray analysis showed that long non-coding RNA (lncRNA) AK021954 gene expression was significantly increased in the baicalein treated group compared with the control group.''
\noindent [116] Susannah Havermann; Yvonni Chovolou; Hans-Ulrich Humpf; Wim Wätjen (2016). Modulation of the Nrf2 signalling pathway in Hct116 colon carcinoma cells by baicalein and its methylated derivative negletein.. \textit{Pharmaceutical biology} 54; 1491-502. \href{https://doi.org/10.3109/13880209.2015.1104703}{doi:10.3109/13880209.2015.1104703}. \url{https://pubmed.ncbi.nlm.nih.gov/27143122} (accessed 2026-11-14).
\noindent\hspace*{2em}\small\textit{Evidence:} ``No change in the ubiquitinylation of Nrf2, proteasomal activity and transcription of the NRF2 gene were detectable.''
\noindent\hspace*{2em}\small\textit{Evidence:} ``Baicalein decreased the phosphorylation of Nrf2 (IC50-value approximately 20 μM) suggesting an inhibitory effect of the flavonoid on protein kinases.''
\noindent\hspace*{2em}\small\textit{Evidence:} ``RESULTS: No change in the ubiquitinylation of Nrf2, proteasomal activity and transcription of the NRF2 gene were detectable.''
\noindent [117] Yi-Chao Zheng; Dan-Dan Shen; Meng Ren; Xue-Qi Liu; Zhi-Ru Wang; Ying Liu; Qian-Na Zhang; Li-Juan Zhao (2016). Baicalin, a natural LSD1 inhibitor.. \textit{Bioorganic chemistry} 69; 129-131. \href{https://doi.org/10.1016/j.bioorg.2016.10.004}{doi:10.1016/j.bioorg.2016.10.004}. \url{https://pubmed.ncbi.nlm.nih.gov/27814566} (accessed 2026-11-14).
\noindent\hspace*{2em}\small\textit{Evidence:} ``In our study, baicalin was first characterized as a LSD1 inhibitor with an IC50 of 3.01μM and showed strong LSD1 inhibitory effect in cells.''
\noindent [119] Beibei Bie; Jin Sun; Jun Li; Ying Guo; Wei Jiang; Chen Huang; Jun Yang; Zongfang Li (2017). Baicalein, a Natural Anti-Cancer Compound, Alters MicroRNA Expression Profiles in Bel-7402 Human Hepatocellular Carcinoma Cells.. \textit{Cellular physiology and biochemistry : international journal of experimental cellular physiology, biochemistry, and pharmacology} 41; 1519-1531. \href{https://doi.org/10.1159/000470815}{doi:10.1159/000470815}. \url{https://pubmed.ncbi.nlm.nih.gov/28351032} (accessed 2026-11-14).
\noindent\hspace*{2em}\small\textit{Evidence:} ``Putative target genes for differentially expressed miRNAs could be enriched in terms of cell proliferation regulation, cell cycle arrest and were mainly involved in MAPK, PI3K-Akt, Wnt, Hippo and mTOR signaling pathways.''
\noindent [120] Chunya Lu; Huaqi Wang; Shanshan Chen; Rui Yang; Hui Li; Guojun Zhang (2018). Baicalein inhibits cell growth and increases cisplatin sensitivity of A549 and H460 cells via miR-424-3p and targeting PTEN/PI3K/Akt pathway.. \textit{Journal of cellular and molecular medicine} 22; 2478-2487. \href{https://doi.org/10.1111/jcmm.13556}{doi:10.1111/jcmm.13556}. \url{https://pubmed.ncbi.nlm.nih.gov/29392841} (accessed 2026-11-14).
\noindent\hspace*{2em}\small\textit{Evidence:} ``Taken together, we concluded that baicalein inhibits cell growth and increases cisplatin sensitivity to A549 and H460 cells via down-regulation of miR-424-3p and targeting the PTEN/PI3K/Akt pathway.''
\noindent [121] Xinjian Cui; Xionghua Sun; Fanqing Lu; Xiaogang Jiang (2018). Baicalein represses TGF-β1-induced fibroblast differentiation through the inhibition of miR-21.. \textit{Toxicology and applied pharmacology} 358; 35-42. \href{https://doi.org/10.1016/j.taap.2018.09.007}{doi:10.1016/j.taap.2018.09.007}. \url{https://pubmed.ncbi.nlm.nih.gov/30201452} (accessed 2026-11-14).
\noindent\hspace*{2em}\small\textit{Evidence:} ``Baicalein diminished the expression of miR-21, and miR-21 mimics partially antagonized the effects of baicalein.''
\noindent [122] Haiyang Yu; Fangfang Zhang; Xidong Guan (2019). Baicalin reverse depressive-like behaviors through regulation SIRT1-NF-kB signaling pathway in olfactory bulbectomized rats.. \textit{Phytotherapy research : PTR} 33; 1480-1489. \href{https://doi.org/10.1002/ptr.6340}{doi:10.1002/ptr.6340}. \url{https://pubmed.ncbi.nlm.nih.gov/30848526} (accessed 2026-11-14).
\noindent\hspace*{2em}\small\textit{Evidence:} ``Moreover, in lipopolysaccharides-induced BV-2 cells, the levels of inflammatory factors (IL-1β), p65 acetylation at lysine 310, and SIRT1 expression were different in the group treated with both baicalin and nicotinamide compared with the group treated with baicalin, which suggests that baicalin regulates SIRT1 and thereby inhibits p65 acetylation.''
\noindent\hspace*{2em}\small\textit{Evidence:} ``Furthermore, baicalin regulated Sirtuin 1 (SIRT1) and decreased the levels of p65 acetylation (ac-p65) in the hippocampus and hypothalamus.''
\noindent [123] Xiaoliang Yang; Chunyan Zhang; Jinjie Jiang; Yinghao Li (2019). Baicalein retards proliferation and collagen deposition by activating p38MAPK-JNK via microRNA-29.. \textit{Journal of cellular biochemistry} 120; 15625-15634. \href{https://doi.org/10.1002/jcb.28829}{doi:10.1002/jcb.28829}. \url{https://pubmed.ncbi.nlm.nih.gov/31081145} (accessed 2026-11-14).
\noindent\hspace*{2em}\small\textit{Evidence:} ``Mechanically, BAI activates the p38MAPK/JNK pathway by promoting miR-29.''
\noindent [124] Jiaqi Li; Yuzhi Zhou; Guanhua Du; Xuemei Qin; Li Gao (2019). Integration of transcriptomics and network analysis deciphers the mechanisms of baicalein in improving learning and memory impairment in senescence-accelerated mouse prone 8 (SAMP8).. \textit{European journal of pharmacology} 865; 172789. \href{https://doi.org/10.1016/j.ejphar.2019.172789}{doi:10.1016/j.ejphar.2019.172789}. \url{https://pubmed.ncbi.nlm.nih.gov/31715136} (accessed 2026-11-14).
\noindent\hspace*{2em}\small\textit{Evidence:} ``Our results suggest that baicalein prevented the activation of JAK2/STAT1 signaling pathway and decreased the levels of Aβ1-42 and RAGE in the cortex of SAMP8 mice.''
\noindent [126] Xiaoxuan Yu; Hui Li; Mengyuan Zhu; Po Hu; Xiao Liu; Yingjie Qing; Xiangyuan Wang; Hongzheng Wang (2020). Involvement of p53 Acetylation in Growth Suppression of Cutaneous T-Cell Lymphomas Induced by HDAC Inhibition.. \textit{The Journal of investigative dermatology} 140; 2009-2022.e4. \href{https://doi.org/10.1016/j.jid.2019.12.041}{doi:10.1016/j.jid.2019.12.041}. \url{https://pubmed.ncbi.nlm.nih.gov/32119867} (accessed 2026-11-14).
\noindent\hspace*{2em}\small\textit{Evidence:} ``Moreover, we discovered that Baicalein, a natural product, exhibited an inhibitory effect on HDAC1 and HDAC8.''
\noindent [131] Xiaoxuan Yu; Hui Li; Po Hu; Yingjie Qing; Xiangyuan Wang; Mengyuan Zhu; Hongzheng Wang; Zhanyu Wang (2020). Natural HDAC-1/8 inhibitor baicalein exerts therapeutic effect in CBF-AML.. \textit{Clinical and translational medicine} 10; e154. \href{https://doi.org/10.1002/ctm2.154}{doi:10.1002/ctm2.154}. \url{https://pubmed.ncbi.nlm.nih.gov/32898337} (accessed 2026-11-14).
\noindent\hspace*{2em}\small\textit{Evidence:} ``Baicalein showed HDAC-1/8 inhibition to trigger growth suppression and differentiation induction of AML cell lines and primary AML cells.''
\noindent [136] Danhui Ma; Sinuo Chen; Heming Wang; Jiayi Wei; Hao Wu; Hong Gao; Xinlai Cheng; Taotao Liu (2021). Baicalein Induces Apoptosis of Pancreatic Cancer Cells by Regulating the Expression of miR-139-3p and miR-196b-5p.. \textit{Frontiers in oncology} 11; 653061. \href{https://doi.org/10.3389/fonc.2021.653061}{doi:10.3389/fonc.2021.653061}. \url{https://pubmed.ncbi.nlm.nih.gov/33996574} (accessed 2026-11-14).
\noindent\hspace*{2em}\small\textit{Evidence:} ``Quantitative PCR analysis confirmed that miR-139-3p was the most up-regulated miRNA after baicalein treatment, while miR-196b-5p was the most down-regulated miRNA.''
\noindent [138] Wenqu Zhou; Mi Gao; Chunxiao Liang; Biting Lin; Qinghua Wu; Ruikun Chen; Xiaoxiao Xiong; Xing Chen (2021). Systematic Understanding of the Mechanism of Baicalin against Gastric Cancer Using Transcriptome Analysis.. \textit{BioMed research international} 2021; 5521058. \href{https://doi.org/10.1155/2021/5521058}{doi:10.1155/2021/5521058}. \url{https://pubmed.ncbi.nlm.nih.gov/34337018} (accessed 2026-11-14).
\noindent\hspace*{2em}\small\textit{Evidence:} ``A total of 68 lncRNAs, 20 miRNAs, and 1648 mRNAs were differentially expressed in baicalin-treated SGC-7901 GC cells.''
\noindent [140] Jiafeng Li; Huixin Liu; Qiwang Lin; Huajiao Chen; Liya Liu; Hongjuan Liao; Ying Cheng; Xiuli Zhang (2022). Baicalin suppresses the migration and invasion of breast cancer cells via the TGF-β/lncRNA-MALAT1/miR-200c signaling pathway.. \textit{Medicine} 101; e29328. \href{https://doi.org/10.1097/MD.0000000000029328}{doi:10.1097/MD.0000000000029328}. \url{https://pubmed.ncbi.nlm.nih.gov/36401368} (accessed 2026-11-14).
\noindent\hspace*{2em}\small\textit{Evidence:} ``Additionally, baicalin treatment significantly downregulated the expression of lncRNA-MALAT1 and upregulated that of miR-200c.''
\noindent [141] Rong Cai; Yan-Ping Zhou; Yun-Hai Li; Jin-Jin Zhang; Zuo-Wei Hu (2023). Baicalin Blocks Colon Cancer Cell Cycle and Inhibits Cell Proliferation through miR-139-3p Upregulation by Targeting CDK16.. \textit{The American journal of Chinese medicine} 51; 189-203. \href{https://doi.org/10.1142/S0192415X23500118}{doi:10.1142/S0192415X23500118}. \url{https://pubmed.ncbi.nlm.nih.gov/36599649} (accessed 2026-11-14).
\noindent\hspace*{2em}\small\textit{Evidence:} ``The results showed that baicalin increased miR-139-3p expression while also decreasing CDK16 levels, blocking the cell cycle, and inhibiting cell proliferation in colon cancer cells.''
\noindent [142] Xuefen Xu; Shufan Ji; Yuan Chen; Siwei Xia; Yang Li; Li Chen; Yujia Li; Feng Zhang (2023). Induction of DNMT1-dependent demethylation of SHP-1 by the natural flavonoid compound Baicalein overcame Imatinib-resistance in CML CD34+ cells.. \textit{Cell communication and signaling : CCS} 21; 47. \href{https://doi.org/10.1186/s12964-023-01049-9}{doi:10.1186/s12964-023-01049-9}. \url{https://pubmed.ncbi.nlm.nih.gov/36869331} (accessed 2026-11-14).
\noindent\hspace*{2em}\small\textit{Evidence:} ``Baicalein induced DNMT1-mediated demethylation of the SHP-1 promoter region, and subsequently activated SHP-1 re-expression, which resulted in an inhibition of JAK2/STAT5 signaling in resistant CML CD34+ cells.''
\noindent [143] Haitao Zhang; Jinfeng Cai; Chunna Li; Lisi Deng; Hongqiong Zhu; Ting Huang; Jiacong Zhao; Jiasheng Zhou (2023). Wogonin inhibits latent HIV-1 reactivation by downregulating histone crotonylation.. \textit{Phytomedicine : international journal of phytotherapy and phytopharmacology} 116; 154855. \href{https://doi.org/10.1016/j.phymed.2023.154855}{doi:10.1016/j.phymed.2023.154855}. \url{https://pubmed.ncbi.nlm.nih.gov/37172478} (accessed 2026-11-14).
\noindent\hspace*{2em}\small\textit{Evidence:} ``Mechanistically, wogonin inhibited the reactivation of latent HIV-1 by inhibiting the expression of p300, a histone acetyltransferase, and decreasing the crotonylation of histone H3/H4 in the HIV-1 promoter region.''
\noindent [145] Ting Wang; Pinghan Wang; Song Wang; Yu Ma; Ziqiao Zhao; Fangyi Long (2023). Wogonin Diminishes Radioresistance of Breast Cancer via Inhibition of the Nrf2/HIF-1[Formula: see text] Pathway.. \textit{The American journal of Chinese medicine} 51; 2243-2262. \href{https://doi.org/10.1142/S0192415X23500969}{doi:10.1142/S0192415X23500969}. \url{https://pubmed.ncbi.nlm.nih.gov/37903716} (accessed 2026-11-14).
\noindent\hspace*{2em}\small\textit{Evidence:} ``Mechanistically, WG enhanced the level of the Nrf2 inhibitor Keap1 through reducing CpG methylation in the promoter region of the Keap1 gene.''
\noindent\hspace*{2em}\small\textit{Evidence:} ``Taken together, our findings demonstrate that WG can epigenetically regulate the Keap1 gene, inhibit the Nrf2/HIF-1[Formula: see text] pathway, induce apoptosis in breast cancer cells, and diminish acquired radioresistance.''
\noindent [147] Yun Dong; Gaojian He; Kun Chen; Xuefeng He; Meitong Pan; Xuemei Huang; Xiaolan Yu; Jiyi Xia (2024). Baicalein promotes KDM4E to induce BICD1 and inhibit triple-negative breast cancer progression by blocking PAR1 signaling.. \textit{Molecular carcinogenesis} 63; 1288-1302. \href{https://doi.org/10.1002/mc.23724}{doi:10.1002/mc.23724}. \url{https://pubmed.ncbi.nlm.nih.gov/38607237} (accessed 2026-11-14).
\noindent\hspace*{2em}\small\textit{Evidence:} ``KDM4E activated protein bicaudal D homolog 1 (BICD1) expression by reducing the deposition of histone H3 lysine 9 trimethylation (H3K9me3) in its promoter, whereas BICD1 promoted protease-activated receptor-1 (PAR1) endocytosis and blocked PAR1 signaling through physical interaction with PAR1.''
\noindent [149] Yang Chen; Shujun Bao; Zhongzhao Wang; Zheng Fang; Hao Tang (2024). Baicalin promotes the sensitivity of NSCLC to cisplatin by regulating ferritinophagy and macrophage immunity through the KEAP1-NRF2/HO-1 pathway.. \textit{European journal of medical research} 29; 387. \href{https://doi.org/10.1186/s40001-024-01930-4}{doi:10.1186/s40001-024-01930-4}. \url{https://pubmed.ncbi.nlm.nih.gov/39061086} (accessed 2026-11-14).
\noindent\hspace*{2em}\small\textit{Evidence:} ``BA induced accumulation of KEAP1 and reduction of NRF2 and HO-1.''
\noindent [152] Mohanna Khandan; Mohammad Amin Khazeei Tabari; Seyed Mostafa Rahimi; Mahmoud Hassani; Abouzar Bagheri (2025). The effects of flavonoid baicalein on miRNA expressions in cancer: a systematic review.. \textit{Naunyn-Schmiedeberg's archives of pharmacology} 398; 11331-11348. \href{https://doi.org/10.1007/s00210-025-04078-y}{doi:10.1007/s00210-025-04078-y}. \url{https://pubmed.ncbi.nlm.nih.gov/40153015} (accessed 2026-11-14).
\noindent\hspace*{2em}\small\textit{Evidence:} ``Fifteen studies meeting the inclusion criteria were included in the systematic review.''
\noindent [153] Lian Zhao; Gong Chen; Dan Li; Kangtao Wang; Michael Schaefer; Ingrid Herr; Bin Yan (2025). Baicalein disrupts TGF-β-induced EMT in pancreatic cancer by FTO-dependent m6A demethylation of ZEB1.. \textit{Biochimica et biophysica acta. Molecular cell research} 1872; 119969. \href{https://doi.org/10.1016/j.bbamcr.2025.119969}{doi:10.1016/j.bbamcr.2025.119969}. \url{https://pubmed.ncbi.nlm.nih.gov/40262723} (accessed 2026-11-14).
\noindent\hspace*{2em}\small\textit{Evidence:} ``It also downregulated FTO, an enzyme critical for m6A RNA demethylation.''
\noindent [154] Zhenggang Wu; Jing Liu; Deju Yin; Jing Huang; Yujing Huang; Pengfei Wang (2025). Baicalein Alleviates Lithium-Pilocarpine-Induced Status Epilepticus by Regulating DNMT1/GABRD Pathway in Rats.. \textit{Organogenesis} 21; 2519607. \href{https://doi.org/10.1080/15476278.2025.2519607}{doi:10.1080/15476278.2025.2519607}. \url{https://pubmed.ncbi.nlm.nih.gov/40569104} (accessed 2026-11-14).
\noindent\hspace*{2em}\small\textit{Evidence:} ``Mechanistically, baicalein interacted with DNMT1 to suppress GABRD promoter region methylation, thus increasing GABRD protein level in the hippocampus of rats induced by LiCl-PILO.''
\noindent [155] Zhe Cai; Chengyu Lu; Dehui Chen; Song Zhang; Juan Zhou; Shuilian Yu; Chun Kwok Wong; Zhaojun Jia (2025). Wogonin modulates macrophage polarization and inflammatory signaling through the LSD1-p65 axis to alleviate osteoarthritis.. \textit{Phytomedicine : international journal of phytotherapy and phytopharmacology} 146; 157149. \href{https://doi.org/10.1016/j.phymed.2025.157149}{doi:10.1016/j.phymed.2025.157149}. \url{https://pubmed.ncbi.nlm.nih.gov/40818185} (accessed 2026-11-14).
\noindent\hspace*{2em}\small\textit{Evidence:} ``Mechanistically, wogonin disrupted LSD1-mediated NF-κB p65 demethylation, enhancing p65 methylation and inhibiting JNK/NF-κB activation.''
\noindent\hspace*{2em}\small\textit{Evidence:} ``wogonin bound lysine-specific histone demethylase 1A (LSD1/KDM1A), suppressing LPS-induced M1 polarization, while increasing CD163+ M2 macrophages in OA joints.''
\noindent [156] Yaqiu Zheng; Yu Li; Minling Lu; Zhen Jia; Ting Wan; Zizheng Li; Yukai Huang; Chishun Zhou (2025). Baicalein specifically suppresses microsatellite instability colorectal cancer by targeting adenosylhomocysteinase to inhibit histone H3 lysine 4 trimethylation-mediated cancer stemness.. \textit{Phytomedicine : international journal of phytotherapy and phytopharmacology} 148; 157291. \href{https://doi.org/10.1016/j.phymed.2025.157291}{doi:10.1016/j.phymed.2025.157291}. \url{https://pubmed.ncbi.nlm.nih.gov/41038143} (accessed 2026-11-14).
\noindent\hspace*{2em}\small\textit{Evidence:} ``This inhibition triggered S-adenosylhomocysteine (SAH) accumulation, decreased histone H3 lysine 4 trimethylation (H3K4me3), and downregulated stemness markers OCT4 and Nanog.''
\noindent\hspace*{2em}\small\textit{Evidence:} ``Mechanistically, adenosylhomocysteinase (AHCY) was identified as a key direct target of baicalein.''
\noindent [157] Mengyuan Li; Yuhang Hu; Jinglu Yu; Mengting Shen; Weizhi Zhang; Jianjian Zheng; Zhenxu Zhou (2025). Baicalein facilitates hepatic stellate cell ferroptosis via the DNMT1/SCARA5/GPX4 axis.. \textit{Phytomedicine : international journal of phytotherapy and phytopharmacology} 149; 157554. \href{https://doi.org/10.1016/j.phymed.2025.157554}{doi:10.1016/j.phymed.2025.157554}. \url{https://pubmed.ncbi.nlm.nih.gov/41270385} (accessed 2026-11-14).
\noindent\hspace*{2em}\small\textit{Evidence:} ``The SCARA5 promoter harbored a CpG island, and its methylation was facilitated by DNMT1, which could be targeted by baicalein.''
\noindent [158] Donghao Wu; Cheng Liu; Longsheng Zhan; Min Liang; Qianqian Li; Qinghai Lian (2026). Integrative single-cell and spatial transcriptomics analysis reveals a baicalein-responsive 10-gene signature for non-small cell lung cancer.. \textit{Translational oncology} 70; 102853. \href{https://doi.org/10.1016/j.tranon.2026.102853}{doi:10.1016/j.tranon.2026.102853}. \url{https://pubmed.ncbi.nlm.nih.gov/42251782} (accessed 2026-11-14).
\noindent\hspace*{2em}\small\textit{Evidence:} ``Single‑cell analysis resolved 21 cell populations, revealing that baicalein targets were enriched in a 32‑gene core set.''
\noindent [159] Yanqiu Sun; Jian Liu; Mingyu He; Xiaolu Chen (2026). Baicalin improves the inflammatory response of RA-FLS by targeting the circ\_0000734/miR-197-5p/IKBKB axis.. \textit{Frontiers in immunology} 17; 1828868. \href{https://doi.org/10.3389/fimmu.2026.1828868}{doi:10.3389/fimmu.2026.1828868}. \url{https://pubmed.ncbi.nlm.nih.gov/42292404} (accessed 2026-11-14).
\noindent\hspace*{2em}\small\textit{Evidence:} ``Baicalin, a naturally derived small molecule, has been shown to improve inflammatory responses in RA by regulating circRNAs; however, the underlying regulatory mechanism remains unclear.''
\noindent [160] Xiuying Chen; Xin Wang; Tian Tian; Xin Deng; Mengying Ding; Ziqiao Yuan; Yizhuo Li; Hongjuan Wang (2026). Baicalin aggravates isoproterenol-induced myocardial fibrosis via LSD1/AP-1 pathways.. \textit{Phytomedicine : international journal of phytotherapy and phytopharmacology} 159; 158591. \href{https://doi.org/10.1016/j.phymed.2026.158591}{doi:10.1016/j.phymed.2026.158591}. \url{https://pubmed.ncbi.nlm.nih.gov/42501551} (accessed 2026-11-14).
\noindent\hspace*{2em}\small\textit{Evidence:} ``The data showed that baicalin acted as an independent fibrogenic factor, significantly aggravating ISO-induced ventricular remodeling, particularly myocardial fibrosis, both in vivo and in vitro.''
\noindent [161] Huan Wang; Yuhan Xu; Lei Wang; Wanting Meng; Yanwu Xu; Haidong Guo (2026). Baicalein Attenuates High Glucose and Sodium Palmitate-Induced Ferroptosis in Cardiomyocytes via the Nrf2/SLC7A11/GPX4 Signaling Pathway.. \textit{International journal of molecular sciences} 27. \href{https://doi.org/10.3390/ijms27146391}{doi:10.3390/ijms27146391}. \url{https://pubmed.ncbi.nlm.nih.gov/42511738} (accessed 2026-11-14).
\noindent\hspace*{2em}\small\textit{Evidence:} ``In conclusion, baicalein protects cardiomyocytes from HG + PA-induced injury by activating the Nrf2/SLC7A11/GPX4 signaling pathway and inhibiting ferroptosis.''
\noindent [162] Candela Cives-Losada; Cristiana Soldani; Michela Anna Polidoro; Barbara Franceschini; Flavio Milana; Aurora Gazzillo; Antonella Cammarota; Chiara Camisaschi (2026). Lysine-specific demethylase 1 (LSD1) regulates macrophage plasticity and immune checkpoint expression in colorectal liver metastases.. \textit{Molecular biomedicine} 7. \href{https://doi.org/10.1186/s43556-026-00602-5}{doi:10.1186/s43556-026-00602-5}. \url{https://pubmed.ncbi.nlm.nih.gov/42900293} (accessed 2026-11-14).
\noindent\hspace*{2em}\small\textit{Evidence:} ``Pharmacological inhibition of LSD1 using baicalin or INCB059872 reduced the expression of CD206 and the immune checkpoint molecule PD-L1, consistent with macrophage reprogramming away from a tumour-promoting state.''
\noindent [163] Sun W; Liu P; Wang T; Wang X; Zheng W; Li J (2020). Baicalein reduces hepatic fat accumulation by activating AMPK in oleic acid-induced HepG2 cells and high-fat diet-induced non-insulin-resistant mice.. \textit{Food Funct} 11(1):711-721. \href{https://doi.org/10.1039/c9fo02237f}{doi:10.1039/c9fo02237f}. \url{https://pubmed.ncbi.nlm.nih.gov/31909773} (accessed 2026-11-14).
\noindent\hspace*{2em}\small\textit{Evidence:} ``In this study, we characterized the hepatic fat-lowering activity of baicalein and found that baicalein reduced hepatic fat accumulation by activating AMPK and suppressing SREBP1 cleavage, thus consequently inhibiting the transcriptional activity of SREBP1 and the synthesis of hepatic fat in oleic acid-induced HepG2 cells and high-fat diet-induced non-insulin-resistant mice.''
\noindent [164] Limbach KE; Wen W; Xing Q; Yan J; Yim JH (2023). Baicalein activates 5' adenosine monophosphate-activated protein kinase, inhibits the mammalian target of rapamycin, and exhibits antiproliferative effects in pancreatic neuroendocrine tumors in vitro and in vivo.. \textit{Surgery} 173(1):12-18. \href{https://doi.org/10.1016/j.surg.2022.07.030}{doi:10.1016/j.surg.2022.07.030}. \url{https://pubmed.ncbi.nlm.nih.gov/36207198} (accessed 2026-11-14).
\noindent\hspace*{2em}\small\textit{Evidence:} ``RESULTS: Immunoblotting revealed that treatment of baicalein induced 5' adenosine monophosphate-activated protein kinase activation and the mammalian target of rapamycin inhibition.''
\noindent\hspace*{2em}\small\textit{Evidence:} ``Treatment with baicalein alone led to a significant decrease in the ratio of viable cells compared with controls at 72 hours at concentrations ≥5 μM (P = .021).''
\noindent [165] Aryal P; Kim K; Park PH; Ham S; Cho J; Song K (2014). Baicalein induces autophagic cell death through AMPK/ULK1 activation and downregulation of mTORC1 complex components in human cancer cells.. \textit{FEBS J} 281(20):4644-58. \href{https://doi.org/10.1111/febs.12969}{doi:10.1111/febs.12969}. \url{https://pubmed.ncbi.nlm.nih.gov/25132405} (accessed 2026-11-14).
\noindent\hspace*{2em}\small\textit{Evidence:} ``Here, we show that baicalein leads to human cancer cell death by inducing autophagy rather than apoptosis, because cell death induced by baicalein was completely reversed by suppressing the expression levels of key molecules in autophagy such as Beclin 1, vacuolar protein sorting 34 (Vps34), autophagy-related (Atg)5 and Atg7, but not by pan-caspase inhibitor.''
\noindent\hspace*{2em}\small\textit{Evidence:} ``The data suggested that baicalein induced not only autophagosome formation, but also autophagic flux.''
\noindent [166] Han J; Ji Y; Youn K; Lim G; Lee J; Kim DH; Jun M (2019). Baicalein as a Potential Inhibitor against BACE1 and AChE: Mechanistic Comprehension through In Vitro and Computational Approaches.. \textit{Nutrients} 11(11):. \href{https://doi.org/10.3390/nu11112694}{doi:10.3390/nu11112694}. \url{https://pubmed.ncbi.nlm.nih.gov/31703329} (accessed 2026-11-14).
\noindent\hspace*{2em}\small\textit{Evidence:} ``Baicalein exhibited strong BACE1 and AChE inhibitory properties (IC50 23.71 ± 1.91 µM and 45.95 ± 3.44 µM, respectively) and reacted in non-competitive and competitive manners with substrates, respectively. in Silico docking analysis was in full agreement with the in vitro results, demonstrating that the compound exhibited powerful binding interaction with target enzymes.''
\noindent [167] Zhao C; Schrapel D; Schaefer M (2025). Baicalein Enhances Longevity and Healthspan of C. elegans Through the Insulin/IGF-1 Signaling Pathway.. \textit{MedComm (2020)} 6(12):e70543. \href{https://doi.org/10.1002/mco2.70543}{doi:10.1002/mco2.70543}. \url{https://pubmed.ncbi.nlm.nih.gov/41415716} (accessed 2026-11-14).
\noindent\hspace*{2em}\small\textit{Evidence:} ``Mechanistically, baicalein suppressed the DAF-2-mediated insulin/IGF-1 signaling pathway and promoted the nuclear translocation of DAF-16, a pivotal longevity transcription factor.''
\noindent\hspace*{2em}\small\textit{Evidence:} ``The findings indicate that baicalein significantly extends lifespan and enhances health markers, including improved motility, increased oxidative stress resistance, and reduced lipofuscin accumulation.''
\noindent\hspace*{2em}\small\textit{Evidence:} ``extended the median lifespan of wild‐type N2 \_C. elegans\_ to 22 days (vs. 20 days in controls)''
\noindent [169] Tripathi SS; Kumar R; Bissoyi A; Rizvi SI (2022). Baicalein May Act as a Caloric Restriction Mimetic Candidate to Improve the Antioxidant Profile in a Natural Rodent Model of Aging.. \textit{Rejuvenation Res} 25(2):70-78. \href{https://doi.org/10.1089/rej.2021.0071}{doi:10.1089/rej.2021.0071}. \url{https://pubmed.ncbi.nlm.nih.gov/35316094} (accessed 2026-11-14).
\noindent\hspace*{2em}\small\textit{Evidence:} ``Thus, Baicalein may be a potent CRM candidate for antiaging interventions.''
\noindent\hspace*{2em}\small\textit{Evidence:} ``Based on the aforementioned observations we conclude that Baicalein has the potential to maintain extracellular reactive oxygen species levels and redox homeostasis during the aging process, an effect that is similar to CR.''
\noindent [170] Cai P; Lu Y; Yin Z; Wang X; Zhou X; Li Z (2021). Baicalein ameliorates osteoporosis via AKT/FOXO1 signaling.. \textit{Aging (Albany NY)} 13(13):17370-17379. \href{https://doi.org/10.18632/aging.203227}{doi:10.18632/aging.203227}. \url{https://pubmed.ncbi.nlm.nih.gov/34198266} (accessed 2026-11-14).
\noindent\hspace*{2em}\small\textit{Evidence:} ``In addition, glucocorticoid-induced osteoporosis in MC3T3-E1 cells was partially reversed by BN through inhibition of AKT, which, by upregulating FOXO1, enhanced expression of bone turnover markers (ALP, OCN, Runx2, and Col 1) and extracellular matrix mineralization.''
\noindent [171] Gao L; Fang JS; Bai XY; Zhou D; Wang YT; Liu AL; Du GH (2013). In silico target fishing for the potential targets and molecular mechanisms of baicalein as an antiparkinsonian agent: discovery of the protective effects on NMDA receptor-mediated neurotoxicity.. \textit{Chem Biol Drug Des} 81(6):675-87. \href{https://doi.org/10.1111/cbdd.12127}{doi:10.1111/cbdd.12127}. \url{https://pubmed.ncbi.nlm.nih.gov/23461900} (accessed 2026-11-14).
\noindent\hspace*{2em}\small\textit{Evidence:} ``The top two ranked targets catechol-O-methyltransferase (COMT) and monoamine oxidase B (MAO-B) have been proposed as targets of baicalein by literatures.''
\noindent [172] Chen M; Peng L; Gong P; Zheng X; Sun T; Zhang X; Huo J (2021). Baicalein Induces Mitochondrial Autophagy to Prevent Parkinson's Disease in Rats via miR-30b and the SIRT1/AMPK/mTOR Pathway.. \textit{Front Neurol} 12:646817. \href{https://doi.org/10.3389/fneur.2021.646817}{doi:10.3389/fneur.2021.646817}. \url{https://pubmed.ncbi.nlm.nih.gov/35237220} (accessed 2026-11-14).
\noindent\hspace*{2em}\small\textit{Evidence:} ``PD rats elicited decreased neurological score and DA secretion of the striatum, increased neuronal apoptosis, and injury, and reduced number of mitochondria and autophagy, whereas baicalein alleviated neuronal injury and partly recovered mitochondrial dysfunction, 3-methyladenine inhibited the protection of baicalein. miR-30b-5p was elevated and SIRT1 was diminished in PD rats and inhibited by baicalein. miR-30b-5p targeted SIRT1. miR-30b-5p overexpression or SIRT1 silencing annulled the protection of baicalein.''
\noindent [174] Kuwar OK; Kalia N (2025). Anti-inflammatory and antioxidant effects of baicalein: targeting Nrf2, and NFĸB in neurodegenerative disease.. \textit{Inflammopharmacology} 33(3):1303-1310. \href{https://doi.org/10.1007/s10787-025-01698-x}{doi:10.1007/s10787-025-01698-x}. \url{https://pubmed.ncbi.nlm.nih.gov/40014253} (accessed 2026-11-14).
\noindent\hspace*{2em}\small\textit{Evidence:} ``Simultaneously, baicalein enhances the Nrf2 activation, promoting the expression of antioxidant enzymes like HO-1, NQO1, GPx, and SOD, thus countering oxidative stress.''
\noindent\hspace*{2em}\small\textit{Evidence:} ``Baicalein has been shown to modulate the IĸB/NFĸB, and KEAP1/Nrf2 pathways, thereby mitigating neuroinflammation and oxidative stress while supporting mitochondrial health.''
\noindent [175] Liu N; Cui X; Guo T; Wei X; Sun Y; Liu J; Zhang Y; Ma W; Yan W; Chen L (2024). Baicalein Ameliorates Insulin Resistance of HFD/STZ Mice Through Activating PI3K/AKT Signal Pathway of Liver and Skeletal Muscle in a GLP-1R-Dependent Manner.. \textit{Antioxidants (Basel)} 13(10):. \href{https://doi.org/10.3390/antiox13101246}{doi:10.3390/antiox13101246}. \url{https://pubmed.ncbi.nlm.nih.gov/39456499} (accessed 2026-11-14).
\noindent\hspace*{2em}\small\textit{Evidence:} ``Our mechanism research significantly confirmed that baicalein mitigates hepatic and muscular IR through the PI3K/AKT signal pathway, both in vitro and in vivo.''
\noindent [176] Zhang X; Du L; Zhang W; Yang Y; Zhou Q; Du G (2017). Therapeutic effects of baicalein on rotenone-induced Parkinson's disease through protecting mitochondrial function and biogenesis.. \textit{Sci Rep} 7(1):9968. \href{https://doi.org/10.1038/s41598-017-07442-y}{doi:10.1038/s41598-017-07442-y}. \url{https://pubmed.ncbi.nlm.nih.gov/28855526} (accessed 2026-11-14).
\noindent\hspace*{2em}\small\textit{Evidence:} ``Furthermore, in rotenone-induced parkinsonian rats, baicalein treatment effectively restored mitochondrial function and improved mitobiogenesis, as determined by measuring the mitochondrial density and key regulators involved in mitobiogenesis.''
\noindent [177] Ji-Hong Y; Yu M; Ling-Hong Y; Jing-Jing G; Ling-Li X; Lv W; Yong-Mei J (2023). Baicalein attenuates bleomycin-induced lung fibroblast senescence and lung fibrosis through restoration of Sirt3 expression.. \textit{Pharm Biol} 61(1):288-297. \href{https://doi.org/10.1080/13880209.2022.2160767}{doi:10.1080/13880209.2022.2160767}. \url{https://pubmed.ncbi.nlm.nih.gov/36815239} (accessed 2026-11-14).
\noindent\hspace*{2em}\small\textit{Evidence:} ``Baicalein restored the BLM-induced downregulation of Sirt3 expression in lung tissues and silencing of Sirt3 abolished the inhibitory role of baicalein against BLM-induced lung fibrosis, fibroblast senescence and activation of TGF-β1/Smad signalling pathway.''
\noindent\hspace*{2em}\small\textit{Evidence:} ``silencing of Sirt3 abolished the inhibitory role of baicalein against BLM-induced lung fibrosis, fibroblast senescence and activation of TGF-β1/Smad signalling pathway.''
\noindent [178] Fang C; Liu X; Zhang F; Song T (2024). Baicalein Inhibits Cerebral Ischemia-Reperfusion Injury through SIRT6-Mediated FOXA2 Deacetylation to Promote SLC7A11 Expression.. \textit{eNeuro} 11(10):. \href{https://doi.org/10.1523/ENEURO.0174-24.2024}{doi:10.1523/ENEURO.0174-24.2024}. \url{https://pubmed.ncbi.nlm.nih.gov/39299807} (accessed 2026-11-14).
\noindent\hspace*{2em}\small\textit{Evidence:} ``Baicalein promoted FOXA2 deacetylation by upregulating SIRT6 expression.''
\noindent [179] Yu HY; Zhu Y; Zhang XL; Wang L; Zhou YM; Zhang FF; Zhang HT; Zhao XM (2022). Baicalin attenuates amyloid β oligomers induced memory deficits and mitochondria fragmentation through regulation of PDE-PKA-Drp1 signalling.. \textit{Psychopharmacology (Berl)} 239(3):851-865. \href{https://doi.org/10.1007/s00213-022-06076-x}{doi:10.1007/s00213-022-06076-x}. \url{https://pubmed.ncbi.nlm.nih.gov/35103832} (accessed 2026-11-14).
\noindent\hspace*{2em}\small\textit{Evidence:} ``RESULTS: Baicalin treatment significantly reversed the altered learning and memory behaviours in AD mouse model.''
\noindent\hspace*{2em}\small\textit{Evidence:} ``Moreover, treatment with baicalin reversed amyloid-β oligomer (AβO)-induced abnormalities in the succinate dehydrogenase complex iron sulphur subunit B (SDHB) and cytochrome c oxidase components I (COXI) and mitochondrial fragmentation in the hippocampus.''
\noindent [180] Fang P; Yu M; Zhang L; Wan D; Shi M; Zhu Y; Bo P; Zhang Z (2017). Baicalin against obesity and insulin resistance through activation of AKT/AS160/GLUT4 pathway.. \textit{Mol Cell Endocrinol} 448:77-86. \href{https://doi.org/10.1016/j.mce.2017.03.027}{doi:10.1016/j.mce.2017.03.027}. \url{https://pubmed.ncbi.nlm.nih.gov/28359800} (accessed 2026-11-14).
\noindent\hspace*{2em}\small\textit{Evidence:} ``We discovered that administration of baicalin decreased food intake, body weight, HOMA-IR and NT-PGC-1α levels, but enhanced GLUT4, PGC-1α, pP38MAPK, pAKT and pAS160 contents, as well as GLUT4 mRNA, PGC-1α mRNA, PPARγ mRNA, GLUT1 mRNA expression in skeletal muscles of obese mice and myotubes of C2C12 cells, and reversed high fat diet-induced glucose and insulin intolerance, hyperglycemia and insulin resistance in the mice.''
\noindent [182] Jin X; Zhu L; Lu S; Li C; Bai M; Xu E; Shen J; Li Y (2023). Baicalin ameliorates CUMS-induced depression-like behaviors through activating AMPK/PGC-1α pathway and enhancing NIX-mediated mitophagy in mice.. \textit{Eur J Pharmacol} 938:175435. \href{https://doi.org/10.1016/j.ejphar.2022.175435}{doi:10.1016/j.ejphar.2022.175435}. \url{https://pubmed.ncbi.nlm.nih.gov/36463946} (accessed 2026-11-14).
\noindent\hspace*{2em}\small\textit{Evidence:} ``In conclusion, the present study suggested that the antidepressant effect of baicalin may be related to the enhancement of NIX-mediated mitophagy through activating the AMPK/PGC-1α pathway by directly binding to AMPK.''
\noindent\hspace*{2em}\small\textit{Evidence:} ``Baicalin significantly down-regulated light chain 3(LC3)II/I, protein sequestosome 1 (P62), and translocase of the outer membrane 20 (TOM20), and up-regulated Nip-like protein (NIX), Adenylate activated protein kinase (AMPK), and Peroxisome proliferator-activated receptor-gamma coactivator (PGC)-1α.''
\noindent [183] Jiang M; Li Z; Qin X; Chen L; Zhu G (2025). Regulatory Role of Flavonoid Baicalin from Scutellaria baicalensis on AMPK: A Review.. \textit{Am J Chin Med} 53(3):771-801. \href{https://doi.org/10.1142/S0192415X25500296}{doi:10.1142/S0192415X25500296}. \url{https://pubmed.ncbi.nlm.nih.gov/40374371} (accessed 2026-11-14).
\noindent\hspace*{2em}\small\textit{Evidence:} ``In this review, we provide an overview of the regulatory effect of baicalin on AMPK and its upstream and downstream signaling pathways.''
\noindent [184] Cho J; Hong EB; Kim YS; Song J; Ju YH; Kim H; Lee H; Kim H; Nam MH (2024). Baicalin and baicalein from Scutellaria baicalensis Georgi alleviate aberrant neuronal suppression mediated by GABA from reactive astrocytes.. \textit{CNS Neurosci Ther} 30(5):e14740. \href{https://doi.org/10.1111/cns.14740}{doi:10.1111/cns.14740}. \url{https://pubmed.ncbi.nlm.nih.gov/38715318} (accessed 2026-11-14).
\noindent\hspace*{2em}\small\textit{Evidence:} ``Baicalin and baicalein exhibited significant MAO-B inhibition potential.''
\noindent\hspace*{2em}\small\textit{Evidence:} ``Additionally, LPS-induced aberrant tonic inhibition was reversed by both S. baicalensis extract and its key compounds.''
\noindent [185] Cho J; Hong EB; Kim YS; Song J; Ju YH; Kim H; Lee H; Kim H; Nam MH (2024). Baicalin and baicalein from Scutellaria baicalensis Georgi alleviate aberrant neuronal suppression mediated by GABA from reactive astrocytes.. \textit{CNS Neurosci Ther (PMC full text)} 30(5):e14740. \href{https://doi.org/10.1111/cns.14740}{doi:10.1111/cns.14740}. \url{https://pmc.ncbi.nlm.nih.gov/articles/PMC11076983} (accessed 2026-11-14).
\noindent\hspace*{2em}\small\textit{Evidence:} ``We observed that baicalin and baicalein showed potential MAO‐B inhibition effects with IC50 values of 6.262 and 21.01 μM, respectively (Figure 3B).''
\noindent\hspace*{2em}\small\textit{Evidence:} ``Although there is a potential concern regarding the accuracy of the calculated IC50 value for baicalein, it is noteworthy that our calculated value (21.01 μM) closely aligns with a previously reported value of 22.1 μM.''
\noindent\hspace*{2em}\small\textit{Evidence:} ``On the other hand, wogonin showed relatively high IC50 (79.24 μM) (Figure 3C), indicating a less potent MAO‐B inhibition effect.''
\noindent [186] Liu Z; Xiang H; Deng Q; Fu W; Li Y; Yu Z; Qiu Y; Mei Z; Xu L (2023). Baicalin and baicalein attenuate hyperuricemic nephropathy via inhibiting PI3K/AKT/NF-κB signalling pathway.. \textit{Nephrology (Carlton)} 28(6):315-327. \href{https://doi.org/10.1111/nep.14159}{doi:10.1111/nep.14159}. \url{https://pubmed.ncbi.nlm.nih.gov/36974463} (accessed 2026-11-14).
\noindent\hspace*{2em}\small\textit{Evidence:} ``RESULTS: Molecular docking results showed that BA and BAI had good binding ability with PI3K, AKT, p65 and IκBα.''
\noindent\hspace*{2em}\small\textit{Evidence:} ``BA and BAI significantly ameliorated the levels of renal function, decreased the p-PI3K, p-AKT and p-p65 expression, down-regulated the BAX/BCL2 and CASP3, and blunted the mRNA levels of tumour necrosis factor (TNF)-α, interleukin (IL)-1β, and IL-18 in both renal tissue of HN mice and HK-2 cells induced by MSU.''
\noindent [187] Hang Y; Qin X; Ren T; Cao J (2018). Baicalin reduces blood lipids and inflammation in patients with coronary artery disease and rheumatoid arthritis: a randomized, double-blind, placebo-controlled trial.. \textit{Lipids Health Dis} 17(1):146. \href{https://doi.org/10.1186/s12944-018-0797-2}{doi:10.1186/s12944-018-0797-2}. \url{https://pubmed.ncbi.nlm.nih.gov/29935544} (accessed 2026-11-14).
\noindent\hspace*{2em}\small\textit{Evidence:} ``RESULTS: After 12 week treatment, levels of triglycerides, total cholesterol, LDL-cholesterol and apolipoproteins, as well as CT-1 and hs-CRP, were all significantly improved in the baicalin group compared to the placebo group (1.12 ± 0.36 vs 1.87 ± 0.46 mmol/L, 2.87 ± 1.23 vs 3.22 ± 1.07 mmol/L, 1.38 ± 0.41 vs 1.16 ± 0.32 mmol/L, 1.31 ± 0.41 vs 1.23 ± 0.29 g/L, 42.9 ± 13.7 vs 128.4 ± 24.3 ng/mL, 1.64 ± 0.38 vs 3.9 ± 1.4 mg/dL, respectively).''
\noindent\hspace*{2em}\small\textit{Evidence:} ``Significantly higher proportion of patients in the baicalin group (71\%) reported good/moderate EULAR response than the placebo group (53\%).''
\noindent [188] Hu XL; Liang Q; Jie ZJ; Chen WS; Deng JL; Sun GC (2026). Baicalin inhibits the Akt/mTOR/ULK1 signaling pathway to activate autophagy and ameliorate pulmonary fibrosis.. \textit{Int Immunopharmacol} 168(Pt 2):115883. \href{https://doi.org/10.1016/j.intimp.2025.115883}{doi:10.1016/j.intimp.2025.115883}. \url{https://pubmed.ncbi.nlm.nih.gov/41308360} (accessed 2026-11-14).
\noindent\hspace*{2em}\small\textit{Evidence:} ``Mechanistically, autophagy inhibitors (3-MA, HCQ) counteracted Baicalin's anti-fibrotic effects.''
\noindent\hspace*{2em}\small\textit{Evidence:} ``The autophagy inhibitor 3-MA, activator rapamycin, and Akt agonist SC-79 were employed for mechanistic validation.''
\noindent [189] Fu YJ; Xu B; Huang SW; Luo X; Deng XL; Luo S; Liu C; Wang Q; Chen JY; Zhou L (2021). Baicalin prevents LPS-induced activation of TLR4/NF-κB p65 pathway and inflammation in mice via inhibiting the expression of CD14.. \textit{Acta Pharmacol Sin} 42(1):88-96. \href{https://doi.org/10.1038/s41401-020-0411-9}{doi:10.1038/s41401-020-0411-9}. \url{https://pubmed.ncbi.nlm.nih.gov/32457419} (accessed 2026-11-14).
\noindent\hspace*{2em}\small\textit{Evidence:} ``These changes were dose-dependently attenuated by pretreatment baicalin (12.5-50 μM), but not by baicalin post-treatment.''
\noindent [190] Zhang T; Yang D; Meng X (2017). Baicalin protects against gastroduodenal ulcers via the modulation of Nrf2 expression: Experimental, biochemical, and histological analyses.. \textit{Pharmacol Rep} 69(6):1154-1158. \href{https://doi.org/10.1016/j.pharep.2017.07.004}{doi:10.1016/j.pharep.2017.07.004}. \url{https://pubmed.ncbi.nlm.nih.gov/29128794} (accessed 2026-11-14).
\noindent\hspace*{2em}\small\textit{Evidence:} ``CONCLUSION: The present findings suggest that the protective effects of baicalin may be regulated via Nrf2-mediated anti-secretory actions.''
\noindent\hspace*{2em}\small\textit{Evidence:} ``Baicalin was 72\% effective versus the reference drug and 80\% effective against ethanol-induced ulcers.''
\noindent [191] Wang T; Sun X; Zhang Y; Wang Q; Cheng W; Gao Y; Shi X; Jin J (2025). Baicalin Promotes Skeletal Muscle Fiber Remodeling by Activating the p38MAPK/PGC-1α Signaling Pathway.. \textit{J Agric Food Chem} 73(11):6878-6889. \href{https://doi.org/10.1021/acs.jafc.5c00300}{doi:10.1021/acs.jafc.5c00300}. \url{https://pubmed.ncbi.nlm.nih.gov/40103396} (accessed 2026-11-14).
\noindent\hspace*{2em}\small\textit{Evidence:} ``Mechanistic studies revealed that baicalin, as a small molecule, could target and bind to the p38 MAPK protein, increase its expression and phosphorylation levels, and activate the p38 MAPK/PGC-1α signaling pathway.''
\noindent\hspace*{2em}\small\textit{Evidence:} ``The results showed that baicalin could participate in the remodeling of skeletal muscle fibers, promoting the conversion from glycolytic fibers to oxidative fibers in mice and pigs.''
\noindent [192] Xu J; Li Y; Lou M; Xia W; Liu Q; Xie G; Liu L; Liu B; Yang J; Qin M (2018). Baicalin regulates SirT1/STAT3 pathway and restrains excessive hepatic glucose production.. \textit{Pharmacol Res} 136:62-73. \href{https://doi.org/10.1016/j.phrs.2018.08.018}{doi:10.1016/j.phrs.2018.08.018}. \url{https://pubmed.ncbi.nlm.nih.gov/30144531} (accessed 2026-11-14).
\noindent\hspace*{2em}\small\textit{Evidence:} ``Fasting increased SirT1 protein in STAT3 immunoprecipitation products and less in the liver of baicalin-treated mice, indicating that baicalin blocked the binding of SirT1 to STAT3 and thus preserved STAT3 acetylation.''
\noindent\hspace*{2em}\small\textit{Evidence:} ``Baicalin reduces endogenous glucose production via suppression of hepatic gluconeogenesis and decreased SirT1 induction via reducing NAD+ accumulation in an energy-sensing way.''
\noindent [193] Sayed AM; Hassanein EHM; Salem SH; Hussein OE; Mahmoud AM (2020). Flavonoids-mediated SIRT1 signaling activation in hepatic disorders.. \textit{Life Sci} 259:118173. \href{https://doi.org/10.1016/j.lfs.2020.118173}{doi:10.1016/j.lfs.2020.118173}. \url{https://pubmed.ncbi.nlm.nih.gov/32750437} (accessed 2026-11-14).
\noindent\hspace*{2em}\small\textit{Evidence:} ``In this review, we focused on the flavonoids, (-)-epicatechin, ampelopsin, baicalin, delphinidin, fisetin, epigallocatechin-3-gallate, luteolin, pinocembrin, quercetin, silibinin, trans-chalcone and xanthohumol, to verify whether their potential promising hepatoprotective effects are related to activation of SIRT1.''
\noindent [194] Jiang Y; Zhang J; Liu W; Qian X; Zhuang X; Hu C (2025). Scutellariae Radix and Coptidis Rhizoma improve NAFLD via regulation of SIRT6/ACSL5 pathway and SCD1.. \textit{J Ethnopharmacol} 348:119834. \href{https://doi.org/10.1016/j.jep.2025.119834}{doi:10.1016/j.jep.2025.119834}. \url{https://pubmed.ncbi.nlm.nih.gov/40254111} (accessed 2026-11-14).
\noindent\hspace*{2em}\small\textit{Evidence:} ``Scutellariae Radix (SR) and its major constituent, baicalin, enhanced ACSL5 activity via SIRT6-mediated deacetylation, promoting fatty acid oxidation and intracellular lipid utilization.''
\noindent [195] Shin NR; Gu N; Choi HS; Kim H (2020). Combined effects of Scutellaria baicalensis with metformin on glucose tolerance of patients with type 2 diabetes via gut microbiota modulation.. \textit{Am J Physiol Endocrinol Metab} 318(1):E52-E61. \href{https://doi.org/10.1152/ajpendo.00221.2019}{doi:10.1152/ajpendo.00221.2019}. \url{https://pubmed.ncbi.nlm.nih.gov/31770016} (accessed 2026-11-14).
\noindent\hspace*{2em}\small\textit{Evidence:} ``These results suggest that SB with metformin treatment may improve the glucose tolerance and inflammation and influence the gut microbiota community in T2D.''
\noindent\hspace*{2em}\small\textit{Evidence:} ``Some subjects showed high liver enzyme levels after SB treatment, and their microbiota composition at baseline differed with subjects whose liver enzymes were not affected.''
\noindent\hspace*{2em}\small\textit{Evidence:} ``Glucose tolerance was lower in the SB group than the placebo group.''
\noindent [196] Lee HW; Ryu HW; Kang MG; Park D; Lee H; Shin HM; Oh SR; Kim H (2017). Potent inhibition of monoamine oxidase A by decursin from Angelica gigas Nakai and by wogonin from Scutellaria baicalensis Georgi.. \textit{Int J Biol Macromol} 97:598-605. \href{https://doi.org/10.1016/j.ijbiomac.2017.01.080}{doi:10.1016/j.ijbiomac.2017.01.080}. \url{https://pubmed.ncbi.nlm.nih.gov/28109809} (accessed 2026-11-14).
\noindent\hspace*{2em}\small\textit{Evidence:} ``Wogonin (11) effectively inhibited MAO-A and MAO-B (IC50=6.35 and 20.8μM, respectively).''
\noindent\hspace*{2em}\small\textit{Evidence:} ``Furthermore, compounds 5 (decursinol angelate) and 10 (baicalein) were observed to selectively and moderately inhibit MAO-A.''
\noindent [198] Levy RM; Saikovsky R; Shmidt E; Khokhlov A; Burnett BP (2009). Flavocoxid is as effective as naproxen for managing the signs and symptoms of osteoarthritis of the knee in humans: a short-term randomized, double-blind pilot study. \textit{Nutrition research (New York, N.Y.)} 29(5): 298-304. \href{https://doi.org/10.1016/j.nutres.2009.04.003}{doi:10.1016/j.nutres.2009.04.003}. \url{https://pubmed.ncbi.nlm.nih.gov/19555810} (accessed 2026-11-14).
\noindent\hspace*{2em}\small\textit{Evidence:} ``flavocoxid was as effective as naproxen in controlling the signs and symptoms of OA of the knee''
\noindent\hspace*{2em}\small\textit{Evidence:} ``103 subjects were randomly assigned to receive either flavocoxid [500 mg twice daily (BID)] or naproxen (500 mg BID) in a 1-month onset of action trial.''
\noindent\hspace*{2em}\small\textit{Evidence:} ``There were no statistically detectable differences between the flavocoxid and naproxen groups with respect to any of the outcome variables.''
\noindent [199] Huang ST; Wang CY; Yang RC; Chu CJ; Wu HT; Pang JH (2010). Wogonin, an active compound in Scutellaria baicalensis, induces apoptosis and reduces telomerase activity in the HL-60 leukemia cells. \textit{Phytomedicine : international journal of phytotherapy and phytopharmacology} 17(1): 47-54. \href{https://doi.org/10.1016/j.phymed.2009.06.005}{doi:10.1016/j.phymed.2009.06.005}. \url{https://pubmed.ncbi.nlm.nih.gov/19577445} (accessed 2026-11-14).
\noindent\hspace*{2em}\small\textit{Evidence:} ``treatment of HL-60 cells with wogonin or S. baicalensis led to the inhibition of human telomerase reverse transcriptase (hTERT), human telomerase-associated protein 1 (hTP1) and c-myc messenger ribonucleic acid (m-RNA) expression.''
\noindent\hspace*{2em}\small\textit{Evidence:} ``Two of the compounds, baicalin and wogonoside, are converted to baicalein and wogonin, respectively, after treatment with beta-glucuronidase.''
\noindent [200] Havermann S; Rohrig R; Chovolou Y; Humpf HU; Wätjen W (2013). Molecular effects of baicalein in Hct116 cells and Caenorhabditis elegans : activation of the Nrf2 signaling pathway and prolongation of lifespan. \textit{Journal of agricultural and food chemistry} 61(9): 2158-64. \href{https://doi.org/10.1021/jf304553g}{doi:10.1021/jf304553g}. \url{https://pubmed.ncbi.nlm.nih.gov/23339711} (accessed 2026-11-14).
\noindent\hspace*{2em}\small\textit{Evidence:} ``Lifespan analysis revealed that baicalein extends the mean, median, and maximum lifespans of the nematode by 45, 57 and 24\%, respectively.''
\noindent [201] Zhang JA; Yin Z; Ma LW; Yin ZQ; Hu YY; Xu Y; Wu D; Permatasari F; Luo D; Zhou BR (2014). The protective effect of baicalin against UVB irradiation induced photoaging: an in vitro and in vivo study. \textit{PloS one} 9(6): e99703. \href{https://doi.org/10.1371/journal.pone.0099703}{doi:10.1371/journal.pone.0099703}. \url{https://pubmed.ncbi.nlm.nih.gov/24949843} (accessed 2026-11-14).
\noindent\hspace*{2em}\small\textit{Evidence:} ``the irradiated fibroblasts additionally treated with baicalin demonstrated a decrease in the expression of SA-β-gal, a increase in the cell viability, a decrease in the G1 phase cell proportion, a downregulation in the level of senescence-associated and γ-H2AX proteins.''
\noindent [203] Havermann S; Humpf HU; Wätjen W (2016). Baicalein modulates stress-resistance and life span in C. elegans via SKN-1 but not DAF-16. \textit{Fitoterapia} 113: 123-7. \href{https://doi.org/10.1016/j.fitote.2016.06.018}{doi:10.1016/j.fitote.2016.06.018}. \url{https://pubmed.ncbi.nlm.nih.gov/27370100} (accessed 2026-11-14).
\noindent\hspace*{2em}\small\textit{Evidence:} ``baicalein enhances the resistance of C. elegans against lethal thermal and sodium arsenite stress and dose-dependently prolongs the life span of the nematode.''
\noindent\hspace*{2em}\small\textit{Evidence:} ``we conclude, that baicalein increases stress-resistance and life span in C. elegans via SKN-1 but not DAF-16.''
\noindent [204] Duan DD; Wang KX; Zhou YZ; Qin XM; Gao L; Du GH (2017). Baicalein Exerts Beneficial Effects in d-Galactose-Induced Aging Rats Through Attenuation of Inflammation and Metabolic Dysfunction. \textit{Rejuvenation research} 20(6): 506-516. \href{https://doi.org/10.1089/rej.2017.1919}{doi:10.1089/rej.2017.1919}. \url{https://pubmed.ncbi.nlm.nih.gov/28548620} (accessed 2026-11-14).
\noindent\hspace*{2em}\small\textit{Evidence:} ``Our findings suggest that baicalein significantly attenuated memory decline in d-gal-induced aging model, as manifested by increasing recognition index in novel object recognition test, shortening latency time, and increasing platform crossings in Morris water maze test.''
\noindent [205] Xing D; Gao H; Liu Z; Zhao Y; Gong M (2017). Baicalin Inhibits Inflammatory Responses to Interleukin-1β Stimulation in Human Chondrocytes. \textit{Journal of interferon \& cytokine research : the official journal of the International Society for Interferon and Cytokine Research} 37(9): 398-405. \href{https://doi.org/10.1089/jir.2017.0030}{doi:10.1089/jir.2017.0030}. \url{https://pubmed.ncbi.nlm.nih.gov/28829242} (accessed 2026-11-14).
\noindent\hspace*{2em}\small\textit{Evidence:} ``BA treatment led to remarkable reductions of OA-related proinflammatory gene expressions, including interleukin (IL)-6, tumor necrosis factor (TNF), chemokine (C-X-C motif) ligand 1 (CXCL1), and CXCL10.''
\noindent [206] Wang Z; Ma L; Su M; Zhou Y; Mao K; Li C; Peng G; Zhou C; Shen B; Dou J (2018). Baicalin induces cellular senescence in human colon cancer cells via upregulation of DEPP and the activation of Ras/Raf/MEK/ERK signaling. \textit{Cell death \& disease} 9(2): 217. \href{https://doi.org/10.1038/s41419-017-0223-0}{doi:10.1038/s41419-017-0223-0}. \url{https://pubmed.ncbi.nlm.nih.gov/29440765} (accessed 2026-11-14).
\noindent\hspace*{2em}\small\textit{Evidence:} ``we found that baicalin treatment significantly induced senescence in colon cancer cells.''
\noindent\hspace*{2em}\small\textit{Evidence:} ``baicalin upregulated the expression of decidual protein induced by progesterone (DEPP) in HCT116 colon cancer cells, which accompanied with the activation of Ras/Raf/MEK/ERK and p16INK4A/Rb signaling pathways.''
\noindent [207] Gao L; Li J; Zhou Y; Huang X; Qin X; Du G (2018). Effects of Baicalein on Cortical Proinflammatory Cytokines and the Intestinal Microbiome in Senescence Accelerated Mouse Prone 8. \textit{ACS chemical neuroscience} 9(7): 1714-1724. \href{https://doi.org/10.1021/acschemneuro.8b00074}{doi:10.1021/acschemneuro.8b00074}. \url{https://pubmed.ncbi.nlm.nih.gov/29668250} (accessed 2026-11-14).
\noindent\hspace*{2em}\small\textit{Evidence:} ``treatment with baicalein ameliorated the senescence status of the SAMP8 mice, as manifested by reducing the grading score of senescence.''
\noindent\hspace*{2em}\small\textit{Evidence:} ``baicalein significantly inhibited the release of proinflammatory cytokines such as interleukin-6 (IL-6), interleukin-1 beta (IL-1β), and tumor necrosis factor-α (TNF-α) in the brain cortex of SAMP8 mice.''
\noindent [208] Dou J; Wang Z; Ma L; Peng B; Mao K; Li C; Su M; Zhou C; Peng G (2018). Baicalein and baicalin inhibit colon cancer using two distinct fashions of apoptosis and senescence. \textit{Oncotarget} 9(28): 20089-20102. \href{https://doi.org/10.18632/oncotarget.24015}{doi:10.18632/oncotarget.24015}. \url{https://pubmed.ncbi.nlm.nih.gov/29732005} (accessed 2026-11-14).
\noindent\hspace*{2em}\small\textit{Evidence:} ``baicalin-induced senescence in tumor cells is due to its inhibition of telomerase reverse transcriptase expression in tumor cells,''
\noindent [209] Duan DD; Gao L; Wang KX; Qin XM; Zhou YZ; Du GH (2016). [Baicalein prolongs the lifespan of Drosophila melanogaster through antioxidation activity]. \textit{Yao xue xue bao = Acta pharmaceutica Sinica} 51(9): 1401-6. \href{https://pubmed.ncbi.nlm.nih.gov/29924517/}{PMID:29924517}. \url{https://pubmed.ncbi.nlm.nih.gov/29924517} (accessed 2026-11-14).
\noindent\hspace*{2em}\small\textit{Evidence:} ``Baicalein at concentrations of 0.04 mg·m L-1 and 0.2 mg·m L-1 extended the mean and maximum lifespans in the natural aging model, and effectively reduced the damages of oxidative stress by H2O2 and paraquat.''
\noindent [210] Curtis JR; Owensby JK; Xie F (2020). Comparative safety of flavocoxid vs prescription NSAIDs among osteoarthritis patients. \textit{Osteoarthritis and cartilage} 28(7): 917-923. \href{https://doi.org/10.1016/j.joca.2020.03.017}{doi:10.1016/j.joca.2020.03.017}. \url{https://pubmed.ncbi.nlm.nih.gov/32305527} (accessed 2026-11-14).
\noindent\hspace*{2em}\small\textit{Evidence:} ``A lower rate of GI bleed was observed, IR: 5.3 (1.7-12.3) for flavocoxid and 10.2 (5.9-16.3) for NSAIDs, aHR 0.49 (0.18-1.68).''
\noindent\hspace*{2em}\small\textit{Evidence:} ``Flavocoxid users had a significantly lower risk for hospitalized GI bleeding.''
\noindent\hspace*{2em}\small\textit{Evidence:} ``For hospitalized liver injury, it was 3.2 (95\% CI 0.7-9.3) for flavocoxid and 2.4 (95\% CI 0.7-6.1) for NSAIDs, aHR = 1.16, 95\% CI 0.23-6.01.''
\noindent [211] Lim H; Kwon YS; Kim D; Lee J; Kim HP (2020). Flavonoids from Scutellaria baicalensis inhibit senescence-associated secretory phenotype production by interrupting IκBζ/C/EBPβ pathway: Inhibition of age-related inflammation. \textit{Phytomedicine : international journal of phytotherapy and phytopharmacology} 76: 153255. \href{https://doi.org/10.1016/j.phymed.2020.153255}{doi:10.1016/j.phymed.2020.153255}. \url{https://pubmed.ncbi.nlm.nih.gov/32554301} (accessed 2026-11-14).
\noindent\hspace*{2em}\small\textit{Evidence:} ``compounds 8 and 11, the strongest SASP inhibitors, decreased the expression of IκBζ and C/EBPβ protein without affecting either BrdU uptake or the expression of senescence markers, such as pRb and p21.''
\noindent\hspace*{2em}\small\textit{Evidence:} ``the oral administration of compound 8 to aged rats at 2 and 4 mg/kg/day for 10 days significantly inhibited the gene expression of SASP and IκBζ in kidneys.''
\noindent [212] Guerrero-Rubio MA; Hernández-García S; García-Carmona F; Gandía-Herrero F (2021). Flavonoids' Effects on Caenorhabditis elegans' Longevity, Fat Accumulation, Stress Resistance and Gene Modulation Involve mTOR, SKN-1 and DAF-16. \textit{Antioxidants (Basel, Switzerland)} 10(3):. \href{https://doi.org/10.3390/antiox10030438}{doi:10.3390/antiox10030438}. \url{https://pubmed.ncbi.nlm.nih.gov/33809299} (accessed 2026-11-14).
\noindent\hspace*{2em}\small\textit{Evidence:} ``The results showed that chrysin, 6-hydroxyflavone and baicalein prolonged C. elegans' lifespan by up to 8.5\%, 11.8\% and 18.6\%, respectively.''
\noindent\hspace*{2em}\small\textit{Evidence:} ``microarray analysis showed that baicalein downregulates important age-related genes, such as mTOR and PARP.''
\noindent\hspace*{2em}\small\textit{Evidence:} ``TJ375 worms were treated with 100 µM of each flavonoid.''
\noindent [213] Gao L; Zheng WG; Wu XK; Du GH; Qin XM (2021). Baicalein Delays H(2)O(2)-Induced Astrocytic Senescence through Inhibition of Senescence-Associated Secretory Phenotype (SASP), Suppression of JAK2/STAT1/NF-κB Pathway, and Regulation of Leucine Metabolism. \textit{ACS chemical neuroscience} 12(13): 2320-2335. \href{https://doi.org/10.1021/acschemneuro.1c00024}{doi:10.1021/acschemneuro.1c00024}. \url{https://pubmed.ncbi.nlm.nih.gov/34152720} (accessed 2026-11-14).
\noindent\hspace*{2em}\small\textit{Evidence:} ``treatment with baicalein protects T98G cells from H2O2-induced damage, delays cell senescence, inhibits the secretion of SASP (IL-6, IL-8, TNF-α, CXCL1, and MMP-1), and inhibits SASP-related pathways NF-κB and JAK2/STAT1.''
\noindent [214] Saul D; Hohl FE; Franz MK; Meyer I; Taudien S; Roch PJ; Sehmisch S; Komrakova M (2021). Inhibition of Lipoxygenases Showed No Benefit for the Musculoskeletal System in Estrogen Deficient Rats. \textit{Frontiers in endocrinology} 12: 706504. \href{https://doi.org/10.3389/fendo.2021.706504}{doi:10.3389/fendo.2021.706504}. \url{https://pubmed.ncbi.nlm.nih.gov/34354672} (accessed 2026-11-14).
\noindent\hspace*{2em}\small\textit{Evidence:} ``Oral administration of Baicalein did not confirm the reported favorable cortical effects in neither vertebra nor femur.''
\noindent\hspace*{2em}\small\textit{Evidence:} ``The combination therapy did not ameliorate OVX-induced osteoporosis but induced even more bone loss.''
\noindent [215] Gao L; Yang WY; Qi H; Sun CJ; Qin XM; Du GH (2022). Unveiling the anti-senescence effects and senescence-associated secretory phenotype (SASP) inhibitory mechanisms of Scutellaria baicalensis Georgi in low glucose-induced astrocytes based on boolean network. \textit{Phytomedicine : international journal of phytotherapy and phytopharmacology} 99: 153990. \href{https://doi.org/10.1016/j.phymed.2022.153990}{doi:10.1016/j.phymed.2022.153990}. \url{https://pubmed.ncbi.nlm.nih.gov/35202958} (accessed 2026-11-14).
\noindent\hspace*{2em}\small\textit{Evidence:} ``SGE could delay astrocytes senescence by decreasing the staining rate of SA-β-gal, reducing secretions of SASP factors (IL-6, CXCL1, MMP-1), alleviating cell cycle arrest in G0/G1 phase, decreasing the formation of punctate DNA foci and down-regulating the expression of p16INK4A, p21 and γH2A.X.''
\noindent [216] Sun G; Dang Y; Lin Y; Zeng W; Wu Z; Zhang X; Dong D; Wu B (2022). Scutellaria baicalensis Georgi regulates REV-ERBα/BMAL1 to protect against skin aging in mice. \textit{Frontiers in pharmacology} 13: 991917. \href{https://doi.org/10.3389/fphar.2022.991917}{doi:10.3389/fphar.2022.991917}. \url{https://pubmed.ncbi.nlm.nih.gov/36249807} (accessed 2026-11-14).
\noindent\hspace*{2em}\small\textit{Evidence:} ``SBG was identified as an antagonist of REV-ERBα and thus an inducer of BMAL1 expression.''
\noindent [217] Barrera-Vázquez OS; Magos-Guerrero GA; Escobar-Ramírez JL; Gomez-Verjan JC (2023). Natural Products as a Major Source of Candidates for Potential Senolytic Compounds obtained by in silico Screening. \textit{Medicinal chemistry (Shariqah (United Arab Emirates))} 19(7): 653-668. \href{https://doi.org/10.2174/1573406419666221019153537}{doi:10.2174/1573406419666221019153537}. \url{https://pubmed.ncbi.nlm.nih.gov/36278464} (accessed 2026-11-14).
\noindent\hspace*{2em}\small\textit{Evidence:} ``Six drug-like candidates (3,4'-dihydroxypropiophenone, baicalein, α, β-dehydrocurvularin, lovastatin, luteolin, and phloretin) were identified.''
\noindent [218] Mohanty SK; Suchiang K (2023). Baicalein mitigates oxidative stress and enhances lifespan through modulation of Wnt ligands and GATA factor: ELT-3 in Caenorhabditis elegans. \textit{Life sciences} 329: 121946. \href{https://doi.org/10.1016/j.lfs.2023.121946}{doi:10.1016/j.lfs.2023.121946}. \url{https://pubmed.ncbi.nlm.nih.gov/37463652} (accessed 2026-11-14).
\noindent\hspace*{2em}\small\textit{Evidence:} ``it was observed that BAI could enhance the worm's lifespan via ELT-3 and SKN-1 transcription factors''
\noindent\hspace*{2em}\small\textit{Evidence:} ``BAI could downregulate different Wnt ligands mRNA expressions in C. elegans, resulting in enhanced expression of GATA transcription factor ELT-3 and antiaging gene Klotho.''
\noindent [220] Li GG; Chu XF; Xing YM; Xue X; Ihtisham B; Liang XF; Xu JX; Mi Y; Zheng PY (2024). Baicalin Prevents Colon Cancer by Suppressing CDKN2A Protein Expression. \textit{Chinese journal of integrative medicine} 30(11): 1007-1017. \href{https://doi.org/10.1007/s11655-024-4109-6}{doi:10.1007/s11655-024-4109-6}. \url{https://pubmed.ncbi.nlm.nih.gov/38941045} (accessed 2026-11-14).
\noindent\hspace*{2em}\small\textit{Evidence:} ``In vitro, the expressions of CDKN2A, MAPK, and p-AKT were suppressed by baicalin in MC38 and CT26 cells.''
\noindent [221] Wang L; Lin F; Liu Y; Li W; Ding Q; Duan X; Yang L; Bai Z; Zhang M; Guo Y (2024). Wogonin protects against bleomycin-induced mouse pulmonary fibrosis via the inhibition of CDK9/p53-mediated cell senescence. \textit{Frontiers in pharmacology} 15: 1407891. \href{https://doi.org/10.3389/fphar.2024.1407891}{doi:10.3389/fphar.2024.1407891}. \url{https://pubmed.ncbi.nlm.nih.gov/39040475} (accessed 2026-11-14).
\noindent\hspace*{2em}\small\textit{Evidence:} ``wogonin attenuated the expressions of DNA damage marker γ-H2AX and senescence-related markers including phosphorylated p53, p21, retinoblastoma protein (pRB), and senescence-associated β-galactosidase (SA-β-gal).''
\noindent\hspace*{2em}\small\textit{Evidence:} ``wogonin, as a direct and selective inhibitor of cyclin-dependent kinase 9 (CDK9)''
\noindent [222] Yimam M; Horm T; O'Neal A; Chua P; Jiao P; Hong M; Jia Q (2024). Botanical Bioflavonoid Composition from Scutellaria baicalensis- and Acacia catechu-Protected Mice against D-Galactose-Induced Immunosenescence, and Cyclophosphamide Induced Immune Suppression. \textit{Nutrients} 16(18):. \href{https://doi.org/10.3390/nu16183144}{doi:10.3390/nu16183144}. \url{https://pubmed.ncbi.nlm.nih.gov/39339746} (accessed 2026-11-14).
\noindent\hspace*{2em}\small\textit{Evidence:} ``UP446 was administered orally at doses of 100 mg/kg and 200 mg/kg starting in the fifth week of immunosenescence induction.''
\noindent [223] Yang SY; Jhang JS; Huang WL; Tsai LHJ; Tsai MC; Chan CP; Lin RI; Lin HY; Li C; Yeh CC; Chan MWY (2025). Wogonin Inhibits Ovarian Cancer by Activating the AMPK-TET2-5hmC Axis. \textit{Molecular carcinogenesis} 64(3): 440-449. \href{https://doi.org/10.1002/mc.23856}{doi:10.1002/mc.23856}. \url{https://pubmed.ncbi.nlm.nih.gov/39620386} (accessed 2026-11-14).
\noindent\hspace*{2em}\small\textit{Evidence:} ``These results indicate that wogonin promotes DNA demethylation by stabilizing TET2.''
\noindent [224] Kimura Y; Sumiyoshi M; Taniguchi M (2025). Baicalein prevents skin damage, tumorigenesis and tumor growth in chronic ultraviolet B-irradiated hairless mice. \textit{Photochemical \& photobiological sciences : Official journal of the European Photochemistry Association and the European Society for Photobiology} 24(3): 479-497. \href{https://doi.org/10.1007/s43630-025-00700-3}{doi:10.1007/s43630-025-00700-3}. \url{https://pubmed.ncbi.nlm.nih.gov/40095355} (accessed 2026-11-14).
\noindent\hspace*{2em}\small\textit{Evidence:} ``Baicalein (at doses of 10 and 30 mg/kg) suppressed UVB-induced increases in skin thickness, improved skin elasticity, and reduced the number and growth of skin tumors.''
\noindent [225] Chen S; Li R; Zhang C; Guo J; Yang A; Yang D; Mao T; Liu H; Ren H; Liang Z; Chen W; Yi C; Xiong Q; Yin D; Wang X; Yang Y (2025). Baicalein Induces Hepatic Stellate Cell Senescence and Attenuates Liver Fibrosis via CEBPZ/p53/HK2-Mediated Glycolysis Inhibition. \textit{Phytomedicine : international journal of phytotherapy and phytopharmacology} 148: 157222. \href{https://doi.org/10.1016/j.phymed.2025.157222}{doi:10.1016/j.phymed.2025.157222}. \url{https://pubmed.ncbi.nlm.nih.gov/41075517} (accessed 2026-11-14).
\noindent\hspace*{2em}\small\textit{Evidence:} ``BA suppressed cellular activation while inducing senescence through G1-phase cell cycle arrest (P < 0.001) and glycolysis inhibition (lactate decreasing: P < 0.001).''
\noindent [226] Long DM; Martinez J; Soumyanath A; Kretzschmar D (2026). Scutellaria lateriflora Extract Supplementation Provides Resilience to Age-Related Phenotypes in Drosophila melanogaster. \textit{International journal of molecular sciences} 27(1):. \href{https://doi.org/10.3390/ijms27010461}{doi:10.3390/ijms27010461}. \url{https://pubmed.ncbi.nlm.nih.gov/41516334} (accessed 2026-11-14).
\noindent\hspace*{2em}\small\textit{Evidence:} ``the beneficial effects on mobility were mainly due to baicalin, whereas sleep was improved by a wogonin mix.''
\noindent\hspace*{2em}\small\textit{Evidence:} ``Since neither the extracts nor the compounds extend the lifespan, this suggests that they improve neuronal health and function and do not generally slow down the aging process.''
\noindent [227] Yuan M; Zhou X; Xu W; Liu J; Han T; Huang Y; Liao F; Chang J; Miao C (2026). Wogonin ameliorates osteoarthritis by targeting the STAT3/PIM1 axis to attenuate chondrocyte senescence. \textit{Phytomedicine : international journal of phytotherapy and phytopharmacology} 153: 157986. \href{https://doi.org/10.1016/j.phymed.2026.157986}{doi:10.1016/j.phymed.2026.157986}. \url{https://pubmed.ncbi.nlm.nih.gov/41722125} (accessed 2026-11-14).
\noindent\hspace*{2em}\small\textit{Evidence:} ``concomitantly retarded SA-β-gal, senescence-associated secretory phenotype (SASP), γ-H2A.X, P16, P21 and P53.''
\noindent\hspace*{2em}\small\textit{Evidence:} ``Wogonin could ameliorate articular cartilage injury, inhibit synovitis and osteophyte formation in DMM mice''
\noindent [228] Li, Min, Shi, Aixin, Pang, Hongxian, Xue, Wei, Li, Yang, Cao, Guoying, Yan, Bei, Dong, Fan, Li, Kexin, Xiao, Wei, He, Guorong, Du, Guanhua, Hu, Xin (2014). Safety, tolerability, and pharmacokinetics of a single ascending dose of baicalein chewable tablets in healthy subjects.. \textit{Journal of ethnopharmacology} :. \href{https://doi.org/10.1016/j.jep.2014.08.031}{doi:10.1016/j.jep.2014.08.031}. \url{https://pubmed.ncbi.nlm.nih.gov/25219601} (accessed 2026-11-14).
\noindent\hspace*{2em}\small\textit{Evidence:} ``MATERIALS AND METHODS: This was a Phase I, randomized, double-blind, single-dose trial of baicalein (100-2800 mg) in 72 healthy adults.''
\noindent\hspace*{2em}\small\textit{Evidence:} ``CONCLUSIONS: Single oral doses of 100-2800 mg of baicalein were safe and well tolerated by healthy subjects.''
\noindent\hspace*{2em}\small\textit{Evidence:} ``The total urinary clearance of baicalein and baicalin was <1\%.''
\noindent [229] Muto, R, Motozuka, T, Nakano, M, Tatsumi, Y, Sakamoto, F, Kosaka, N (1998). [The chemical structure of new substance as the metabolite of baicalin and time profiles for the plasma concentration after oral administration of sho-saiko-to in human].. \textit{Yakugaku zasshi : Journal of the Pharmaceutical Society of Japan} :. \href{https://doi.org/10.1248/yakushi1947.118.3\_79}{doi:10.1248/yakushi1947.118.3\_79}. \url{https://pubmed.ncbi.nlm.nih.gov/9549431} (accessed 2026-11-14).
\noindent\hspace*{2em}\small\textit{Evidence:} ``The metabolite was identified as baicalein 6-O-sulfate (BS) by comparing its retention time in HPLC and electrospray ionization mass spectra (ESI-MS)/MS methods with that of an authentic sample.''
\noindent\hspace*{2em}\small\textit{Evidence:} ``BG was detected after 1 h, reached a maximum level at 5 h and then decreased to the level less than the quantitative limit after 36 h, and the plasma level of BG showed two peaks at 12 h and 24 h.''
\noindent [230] Li, Lijun, Gao, Hongzhi, Lou, Kun, Luo, Hongmei, Hao, Sheng, Yuan, Jing, Liu, Zeyuan, Dong, Ruihua (2021). Safety, tolerability, and pharmacokinetics of oral baicalein tablets in healthy Chinese subjects: A single-center, randomized, double-blind, placebo-controlled multiple-ascending-dose study.. \textit{Clinical and translational science} :. \href{https://doi.org/10.1111/cts.13063}{doi:10.1111/cts.13063}. \url{https://pubmed.ncbi.nlm.nih.gov/34156161} (accessed 2026-11-14).
\noindent\hspace*{2em}\small\textit{Evidence:} ``In this multiple-ascending-dose placebo-controlled trial, 36 healthy subjects were randomized to receive 200, 400, and 600 mg of baicalein tablet or placebo once daily on day 1 and day 10, 3 times daily on days 4-9.''
\noindent\hspace*{2em}\small\textit{Evidence:} ``Oral baicalein tablets were rapidly absorbed with peak plasma levels being reached within 2 h after multiple administration.''
\noindent [231] Kumari, Alpana, Kabra, Atul (2026). Baicalin as a multifunctional flavone glycoside: pharmacological mechanisms, ADME profile, and therapeutic potential.. \textit{Naunyn-Schmiedeberg's archives of pharmacology} :. \href{https://doi.org/10.1007/s00210-026-05381-y}{doi:10.1007/s00210-026-05381-y}. \url{https://pubmed.ncbi.nlm.nih.gov/42105002} (accessed 2026-11-14).
\noindent\hspace*{2em}\small\textit{Evidence:} ``ing clinical evidence suggests that baicalin is generally well tolerated and has translational potential in inflammatory disorders, neurodegenerative diseases, and cardiometabolic conditions, whereas nanodelivery systems such as liposomes and metal-organic frameworks help overcome bioavailability limitations.''
\noindent\hspace*{2em}\small\textit{Evidence:} ``This review underscores the need for well-designed clinical trials to establish standardized dosing, efficacy, and long-term safety, while positioning baicalin as a multifunctional bioactive molecule with substantial promise for future drug development.''
\noindent [232] Liu, Ruijuan, Li, Xia, Wei, Jingyao, Liu, Shuaibing, Chang, Yuanyuan, Zhang, Jiali, Zhang, Ji, Zhang, Xiaojian, Fuhr, Uwe, Taubert, Max, Tian, Xin (2019). A Single Dose of Baicalin Has No Clinically Significant Effect on the Pharmacokinetics of Cyclosporine A in Healthy Chinese Volunteers.. \textit{Frontiers in pharmacology} :. \href{https://doi.org/10.3389/fphar.2019.00518}{doi:10.3389/fphar.2019.00518}. \url{https://pubmed.ncbi.nlm.nih.gov/31156436} (accessed 2026-11-14).
\noindent\hspace*{2em}\small\textit{Evidence:} ``In a two-period study, 16 healthy volunteers received a single 200 mg oral CsA dose alone (reference period) or in combination with 500 mg BG (test period).''
\noindent\hspace*{2em}\small\textit{Evidence:} ``No clinically relevant effect of 500 mg BG co-administration on CsA pharmacokinetics was identified and both treatments were well tolerated.''
\noindent [234] Noh, Keumhan, Kang, Youra, Nepal, Mahesh Raj, Jeong, Ki Sun, Oh, Do Gyeong, Kang, Mi Jeong, Lee, Sangkyu, Kang, Wonku, Jeong, Hye Gwang, Jeong, Tae Cheon (2016). Role of Intestinal Microbiota in Baicalin-Induced Drug Interaction and Its Pharmacokinetics.. \textit{Molecules (Basel, Switzerland)} :. \href{https://doi.org/10.3390/molecules21030337}{doi:10.3390/molecules21030337}. \url{https://pubmed.ncbi.nlm.nih.gov/26978333} (accessed 2026-11-14).
\noindent\hspace*{2em}\small\textit{Evidence:} ``Although there have been numerous reports regarding the pharmacological effects of baicalin and baicalein in vivo and in vitro, some reports indicated that the glycoside form would hardly be absorbed in the intestine and that it should be hydrolyzed to baicalein in advance for absorption.''
\noindent\hspace*{2em}\small\textit{Evidence:} ``In addition, baicalin contains a glucuronide moiety in its structure, by which baicalin and baicalein show complex pharmacokinetic behaviors, due to the interconversion between them by phase II enzymes in the body.''
\noindent [235] Seradj, D-S, Neubauer, L, Senekowitsch, S, Weitschies, W, Schick, P (2024). Metabolism of baicalin by different microbiota determined by MimiCol.. \textit{Die Pharmazie} :. \href{https://doi.org/10.1691/ph.2024.4014}{doi:10.1691/ph.2024.4014}. \url{https://pubmed.ncbi.nlm.nih.gov/39152558} (accessed 2026-11-14).
\noindent\hspace*{2em}\small\textit{Evidence:} ``Enzymes of the colonic microbiota, namely β-glucuronidases, hydrolyze baicalin to the aglycone baicalein.''
\noindent\hspace*{2em}\small\textit{Evidence:} ``While during the experiments with added standard microbiota a metabolism of baicalin to baicalein could be detected, this was not the case in a series of experiments without added microbiota.''
\noindent [237] Wong, Eliza L Y, Sung, Rita Yn Tz, Leung, Ting Fan, Wong, Yeuk Oi, Li, Albert M C, Cheung, Kam Lau, Wong, Chun Kwok, Fok, Tai Fai, Leung, Ping Chung (2009). Randomized, double-blind, placebo-controlled trial of herbal therapy for children with asthma.. \textit{Journal of alternative and complementary medicine (New York, N.Y.)} :. \href{https://doi.org/10.1089/acm.2008.0626}{doi:10.1089/acm.2008.0626}. \url{https://pubmed.ncbi.nlm.nih.gov/19821718} (accessed 2026-11-14).
\noindent\hspace*{2em}\small\textit{Evidence:} ``PATIENTS AND METHODS: In a double-blind, placebo-controlled prospective trial, 85 children with asthma aged 7-15 years were randomly assigned to receive either a daily oral herbal formula of 0.619-g CUF2 capsule of dried aqueous extract with an equal weight of five herbs (Astragalus mongholius Bunge, Cordyceps sinensis Sacc., Radix stemonae, Bulbus fritillariae cirrhosae, and Radix scutellariae) or placebo for 6 months.''
\noindent\hspace*{2em}\small\textit{Evidence:} ``CONCLUSIONS: This study provides no evidence to support the use of the herbal formula of CUF2 in children with asthma.''
\noindent [238] Morgan, Sarah L, Baggott, Joseph E, Moreland, Larry, Desmond, Renee, Kendrach, Angela C (2009). The safety of flavocoxid, a medical food, in the dietary management of knee osteoarthritis.. \textit{Journal of medicinal food} :. \href{https://doi.org/10.1089/jmf.2008.0244}{doi:10.1089/jmf.2008.0244}. \url{https://pubmed.ncbi.nlm.nih.gov/19857081} (accessed 2026-11-14).
\noindent\hspace*{2em}\small\textit{Evidence:} ``The 12-week, randomized, double-blind, placebo-controlled trial in an academic medical center enrolled 59 patients with moderate osteoarthritis of at least one knee who were recruited who were classified as having "below average" to "a moderately above average cardiovascular risk" with a Framingham-based scoring tool.''
\noindent\hspace*{2em}\small\textit{Evidence:} ``There were no intra- or inter-group differences in any of the laboratory parameters from study baseline to completion.''
\noindent [239] Vita, Gian Luca, Sframeli, Maria, Licata, Norma, Bitto, Alessandra, Romeo, Sara, Frisone, Francesca, Ciranni, Annamaria, Pallio, Giovanni, Mannino, Federica, Aguennouz, M'Hammed, Rodolico, Carmelo, Squadrito, Francesco, Toscano, Antonio, Messina, Sonia, Vita, Giuseppe (2021). A Phase 1/2 Study of Flavocoxid, an Oral NF-κB Inhibitor, in Duchenne Muscular Dystrophy.. \textit{Brain sciences} :. \href{https://doi.org/10.3390/brainsci11010115}{doi:10.3390/brainsci11010115}. \url{https://pubmed.ncbi.nlm.nih.gov/33467104} (accessed 2026-11-14).
\noindent\hspace*{2em}\small\textit{Evidence:} ``Thirty-four patients were recruited: 17 were treated with flavocoxid at an oral dose of 250 or 500 mg, according to body weight, for one year; 17 did not receive flavocoxid and served as controls.''
\noindent\hspace*{2em}\small\textit{Evidence:} ``The treatment was well tolerated and nobody dropped out.''
\noindent [240] Chalasani, Naga, Vuppalanchi, Raj, Navarro, Victor, Fontana, Robert, Bonkovsky, Herbert, Barnhart, Huiman, Kleiner, David E, Hoofnagle, Jay H (2012). Acute liver injury due to flavocoxid (Limbrel), a medical food for osteoarthritis: a case series.. \textit{Annals of internal medicine} :. \href{https://doi.org/10.7326/0003-4819-156-12-201206190-00006}{doi:10.7326/0003-4819-156-12-201206190-00006}. \url{https://pubmed.ncbi.nlm.nih.gov/22711078} (accessed 2026-11-14).
\noindent\hspace*{2em}\small\textit{Evidence:} ``RESULTS: Among 877 patients enrolled in the prospective study, 4 had liver injury suspected to have been caused by flavocoxid.''
\noindent\hspace*{2em}\small\textit{Evidence:} ``CONCLUSION: Flavocoxid can cause clinically significant liver injury, which seems to resolve within weeks after cessation.''
\noindent [241] Levy, Robert M, Khokhlov, Alexander, Kopenkin, Sergey, Bart, Boris, Ermolova, Tatiana, Kantemirova, Raiasa, Mazurov, Vadim, Bell, Marjorie, Caldron, Paul, Pillai, Lakshmi, Burnett, Bruce P (2010). Efficacy and safety of flavocoxid, a novel therapeutic, compared with naproxen: a randomized multicenter controlled trial in subjects with osteoarthritis of the knee.. \textit{Advances in therapy} :. \href{https://doi.org/10.1007/s12325-010-0064-z}{doi:10.1007/s12325-010-0064-z}. \url{https://pubmed.ncbi.nlm.nih.gov/20845002} (accessed 2026-11-14).
\noindent\hspace*{2em}\small\textit{Evidence:} ``METHODS: In this randomized, multicenter, double-blind study, 220 subjects were assigned to receive either flavocoxid (500 mg twice daily) or naproxen (500 mg twice daily) for 12 weeks.''
\noindent\hspace*{2em}\small\textit{Evidence:} ``There were no statistically significant differences in efficacy between the flavocoxid and naproxen groups when the entire intent-to-treat population was analyzed.''
\noindent\hspace*{2em}\small\textit{Evidence:} ``The flavocoxid group had significantly fewer upper gastrointestinal (UGI) and renal (edema) adverse events (AEs) as well as a strong trend toward fewer respiratory AEs.''
\noindent [242] Pillai, Lakshmi, Levy, Robert M, Yimam, Mesfin, Zhao, Yuan, Jia, Qi, Burnett, Bruce P (2010). Flavocoxid, an anti-inflammatory agent of botanical origin, does not affect coagulation or interact with anticoagulation therapies.. \textit{Advances in therapy} :. \href{https://doi.org/10.1007/s12325-010-0040-7}{doi:10.1007/s12325-010-0040-7}. \url{https://pubmed.ncbi.nlm.nih.gov/20556559} (accessed 2026-11-14).
\noindent\hspace*{2em}\small\textit{Evidence:} ``Flavocoxid did not significantly inhibit thromboxane production or platelet aggregation, and did not increase bleeding times in healthy volunteers.''
\noindent\hspace*{2em}\small\textit{Evidence:} ``Finally, flavocoxid did not inhibit or potentiate the anticoagulant effect of warfarin.''
\noindent [243] Fong, Sophia Yui Kau, Li, Chenrui, Ho, Yiu Cheong, Li, Rui, Wang, Qian, Wong, Yin Cheong, Xue, Hong, Zuo, Zhong (2017). Brain Uptake of Bioactive Flavones in Scutellariae Radix and Its Relationship to Anxiolytic Effect in Mice.. \textit{Molecular pharmaceutics} :. \href{https://doi.org/10.1021/acs.molpharmaceut.7b00029}{doi:10.1021/acs.molpharmaceut.7b00029}. \url{https://pubmed.ncbi.nlm.nih.gov/28426226} (accessed 2026-11-14).
\noindent\hspace*{2em}\small\textit{Evidence:} ``Our results indicated that all six SR flavones including three aglycons (baicalein, wogonin, and oroxylin A) and three glucuronides (baicalin, wogonoside, and oroxyloside) could pass through the blood-brain barrier, with brain concentrations ranging from 7.9 to 224.0 pmol/g.''
\noindent\hspace*{2em}\small\textit{Evidence:} ``extract contains flavones that are positive modulators of the benzodiazepine binding site of GABAA receptors (baicalein, wogonin, and baicalin), our behavioral study for the first time indicated that SR extract (a mixture of six flavones) did not elicit significant anxiolytic effect at the studied doses.''
\noindent [244] Pillai, Lakshmi, Burnett, Bruce P, Levy, Robert M, GOAL Study Cooperative Group (2010). GOAL: multicenter, open-label, post-marketing study of flavocoxid, a novel dual pathway inhibitor anti-inflammatory agent of botanical origin.. \textit{Current medical research and opinion} :. \href{https://doi.org/10.1185/03007991003694522}{doi:10.1185/03007991003694522}. \url{https://pubmed.ncbi.nlm.nih.gov/20225990} (accessed 2026-11-14).
\noindent\hspace*{2em}\small\textit{Evidence:} ``METHODS: A total of 1067 individuals at 41 rheumatology practices were enrolled and prescribed flavocoxid, 500 mg b.i.d., for 60 days.''
\noindent\hspace*{2em}\small\textit{Evidence:} ``RESULTS: Of the 1005 patients who completed all follow-up visits, physicians recorded an average improvement in VAS scores from 60.1 +/- 18.8 at baseline to 42.5 +/- 21.9 at 8 weeks (p < 0.001) in 65.8\% of patients.''
\noindent\hspace*{2em}\small\textit{Evidence:} ``A limitation of this study was that it was open-label and not rigorously controlled.''
\noindent [245] Zuo, Feng, Zhou, Zhong-Ming, Yan, Mei-Zhen, Liu, Mei-Lan, Xiong, Yu-Lan, Zhang, Qing, Song, Hong-Yue, Ye, Wen-Hua (2002). Metabolism of constituents in Huangqin-Tang, a prescription in traditional Chinese medicine, by human intestinal flora.. \textit{Biological \& pharmaceutical bulletin} :. \href{https://doi.org/10.1248/bpb.25.558}{doi:10.1248/bpb.25.558}. \url{https://pubmed.ncbi.nlm.nih.gov/12033492} (accessed 2026-11-14).
\noindent\hspace*{2em}\small\textit{Evidence:} ``ubmed) --- Abstract --- In the course of studies on the metabolism of active components of Huangqin-Tang by human intestinal flora (HIF), Huangqin-Tang and all individual herbs in the decoctions were incubated with a human fecal suspension separately.''
\noindent\hspace*{2em}\small\textit{Evidence:} ``All the results suggested that the glycosides of many medicinal herbs could be converted to aglycones by HIF, and the metabolism of most glycosides was improved in the compound prescription.''
\noindent [246] Yang, Fan, Meng, Xianmin, Qin, Shukui, Guo, Qinglong, Wei, Libin, Zhao, Di (2026). Safety and pharmacokinetics evaluation of oroxylin A in Chinese healthy volunteers: a phase I, double-blind, placebo-controlled, single ascending dose, multiple dose, and food effect study.. \textit{Clinical \& translational oncology : official publication of the Federation of Spanish Oncology Societies and of the National Cancer Institute of Mexico} :. \href{https://doi.org/10.1007/s12094-025-04017-6}{doi:10.1007/s12094-025-04017-6}. \url{https://pubmed.ncbi.nlm.nih.gov/40903694} (accessed 2026-11-14).
\noindent\hspace*{2em}\small\textit{Evidence:} ``RESULTS: No dose-limiting toxic events (DLT), no serious adverse events (SAEs) or death occurred, and incidence of gastrointestinal AEs were higher in multiple doses than in single dose.''
\noindent\hspace*{2em}\small\textit{Evidence:} ``The Cmax of OA approximately increased in dose-proportional manner (Cmax slope 1.114), and the AUC0-t increased less than dose increasing (AUC0-t slope 0.7513).''
\noindent\hspace*{2em}\small\textit{Evidence:} ``High-fat diet can increase the Cmax of OA (1.6-fold), and its metabolites increased more significantly (p < 0.05), there was no effect of food on tmax or terminal half-life.''
\noindent [247] Wang, Jilin, Feng, Xiaojiao, Li, Ziwei, Liu, Yiting, Yang, Wenzhuo, Zhang, Tingen, Guo, Pan, Liu, Zhidong, Qi, Dongli, Pi, Jiaxin (2023). The Flavonoid Components of Scutellaria baicalensis: Biopharmaceutical Properties and their Improvement using Nanoformulation Techniques.. \textit{Current topics in medicinal chemistry} :. \href{https://doi.org/10.2174/1568026623666221128144258}{doi:10.2174/1568026623666221128144258}. \url{https://pubmed.ncbi.nlm.nih.gov/36443977} (accessed 2026-11-14).
\noindent\hspace*{2em}\small\textit{Evidence:} ``baicalensis have limited clinical application due to their low water solubility, poor permeability, and microbial transformation in vivo.''
\noindent\hspace*{2em}\small\textit{Evidence:} ``However, due to the limited number of clinical trials, doubts remain about their toxicity and improvements in human absorption as a result of nanoformulations.''
\noindent [248] Zhang, Peng, Chen, Shen Guo, Wang, Jia Ting, Wang, Jin Dong, Chen, Zai Hong, Lin, Hai Sheng (2024). A study on the impact of gargling with compound Scutellaria baicalensis Georgi on oral health and microflora changes in fixed orthodontic patients: An experimental study.. \textit{Medicine} :. \href{https://doi.org/10.1097/MD.0000000000039397}{doi:10.1097/MD.0000000000039397}. \url{https://pubmed.ncbi.nlm.nih.gov/39183390} (accessed 2026-11-14).
\noindent\hspace*{2em}\small\textit{Evidence:} ``Results: The detection levels of PLI, PD, and sBI in the experimental group were lower than those in the control group (P < .05) 3 months after orthodontic surgery (P < .05); after orthodontic correction for 3 months, there was a significant difference in coccus, bacillus, Campylobacter, Clostridium, Helicobacter, and filamentous bacteria between''
\noindent\hspace*{2em}\small\textit{Evidence:} ``CONCLUSION SUBSECTIONS: In fixed orthodontic treatment, S. baicalensis Georgi gargle can effectively inhibit oral pathogens and maintain periodontal health.''
\noindent [249] Puri, Basant K, White, Nikita, Monro, Jean A (2019). The effect of supplementation with Scutellaria baicalensis on hepatic function.. \textit{Medical hypotheses} :. \href{https://doi.org/10.1016/j.mehy.2019.109402}{doi:10.1016/j.mehy.2019.109402}. \url{https://pubmed.ncbi.nlm.nih.gov/31557595} (accessed 2026-11-14).
\noindent\hspace*{2em}\small\textit{Evidence:} ``ssays of alanine transaminase, aspartate transaminase, alkaline phosphatase and bilirubin were carried out in 17 patients (16 female) of average age 38.6 (standard error 4.4) years who had each taken 1335 mg dried root daily for an average of 444 (71) days.''
\noindent\hspace*{2em}\small\textit{Evidence:} ``These results suggest that daily intake of a relatively high level of Scutellaria baicalensis for over a year is not associated with any biochemical or clinical evidence of hepatic dysfunction.''
\noindent\hspace*{2em}\small\textit{Evidence:} ``Pre- and post-supplementation serum assays of alanine transaminase, aspartate transaminase, alkaline phosphatase and bilirubin were carried out in 17 patients (16 female) of average age 38.6 (standard error 4.4) years who had each taken 1335 mg dried root daily for an average of 444 (71) days.''
\noindent [251] Arweiler, Nicole Birgit, Pergola, Gregory, Kuenz, Jana, Hellwig, Elmar, Sculean, Anton, Auschill, Thorsten Mathias (2011). Clinical and antibacterial effect of an anti-inflammatory toothpaste formulation with Scutellaria baicalensis extract on experimental gingivitis.. \textit{Clinical oral investigations} :. \href{https://doi.org/10.1007/s00784-010-0471-1}{doi:10.1007/s00784-010-0471-1}. \url{https://pubmed.ncbi.nlm.nih.gov/20936314} (accessed 2026-11-14).
\noindent\hspace*{2em}\small\textit{Evidence:} ``Forty subjects (20 per group) discontinued all oral hygiene measures for four teeth for a period of 21 days using a shield (to generate a possible gingivitis) while they could brush the other teeth normally.''
\noindent\hspace*{2em}\small\textit{Evidence:} ``The results suggest that the new toothpaste formulation was able to significantly reduce the extent of gingivitis, plaque development, and vital flora.''
\noindent [252] Dormal, Valérie, Suchareau, Marlène, Copine, Sylvie, Simar, Laurent, Deldicque, Louise (2025). The Effects of Combined Scutellaria and Saffron Supplementation on Mood Regulation in Participants with Mild-to-Moderate Depressive Symptoms: A Randomized, Double-Blind, Placebo-Controlled Study.. \textit{Nutrients} :. \href{https://doi.org/10.3390/nu17050809}{doi:10.3390/nu17050809}. \url{https://pubmed.ncbi.nlm.nih.gov/40077679} (accessed 2026-11-14).
\noindent\hspace*{2em}\small\textit{Evidence:} ``Methods: In a randomized, double-blind, placebo-controlled trial, 180 participants with mild-to-moderate depressive symptoms were assigned to receive either scutellaria extract alone (SCUTELL'UP®), saffron extract alone (SAFFR'ACTIV®), a combination of scutellaria and saffron extracts (SAFFR'UP®), or a placebo for six weeks.''
\noindent\hspace*{2em}\small\textit{Evidence:} ``For the first time in humans, scutellaria extract demonstrated a positive effect on mood regulation.''
\noindent [253] Lai, Miao-Ying, Hsiu, Su-Lan, Chen, Chung-Chuan, Hou, Yu-Chi, Chao, Pei-Dawn Lee (2003). Urinary pharmacokinetics of baicalein, wogonin and their glycosides after oral administration of Scutellariae Radix in humans.. \textit{Biological \& pharmaceutical bulletin} :. \href{https://doi.org/10.1248/bpb.26.79}{doi:10.1248/bpb.26.79}. \url{https://pubmed.ncbi.nlm.nih.gov/12520178} (accessed 2026-11-14).
\noindent\hspace*{2em}\small\textit{Evidence:} ``Ten healthy male volunteers received a dose of 5.2 g commercial powder (comparable to 9 g crude drug), respectively.''
\noindent\hspace*{2em}\small\textit{Evidence:} ``Our results showed that the mean cumulated renal excretion of baicalein glucuronides and sulfates were 43.1+/-4.5 micromol (2.9\% of dose) and 64.8+/-6.3 micromol (4.3\% of dose), respectively, whereas wogonin glucuronides and sulfates were 21.6+/-2.0 micromol (5.9\% of dose) and 20.7+/-1.7 micromol (5.7\% of dose), respectively.''
\noindent\hspace*{2em}\small\textit{Evidence:} ``The baicalein sulfates was predominant than the corresponding glucuronides, whereas wogonin sulfates was comparable to the corresponding glucuronides.''
\noindent [254] Zhang, Shengman, Liu, Ying, Peng, Ni, Chen, Xiaoyan (2022). Qualitative and quantitative determination of the primary active components and metabolites in human plasma after oral administration of Shuanghuanglian liquid.. \textit{Journal of separation science} :. \href{https://doi.org/10.1002/jssc.202200121}{doi:10.1002/jssc.202200121}. \url{https://pubmed.ncbi.nlm.nih.gov/35713475} (accessed 2026-11-14).
\noindent\hspace*{2em}\small\textit{Evidence:} ``Thirteen components from baicalin were identified, including the parent drug baicalin and aglycone baicalein.''
\noindent\hspace*{2em}\small\textit{Evidence:} ``The validated method was successfully applied to a pharmacokinetic study of an oral administration of 60 ml Shuanghuanglian in healthy subjects.''
\noindent [255] Arjmandi, Bahram H, Ormsbee, Lauren T, Elam, Marcus L, Campbell, Sara C, Rahnama, Nader, Payton, Mark E, Brummel-Smith, Ken, Daggy, Bruce P (2014). A combination of Scutellaria baicalensis and Acacia catechu extracts for short-term symptomatic relief of joint discomfort associated with osteoarthritis of the knee.. \textit{Journal of medicinal food} :. \href{https://doi.org/10.1089/jmf.2013.0010}{doi:10.1089/jmf.2013.0010}. \url{https://pubmed.ncbi.nlm.nih.gov/24611484} (accessed 2026-11-14).
\noindent\hspace*{2em}\small\textit{Evidence:} ``The UP446 group experienced a significant decrease in perceived pain (P=.009) time dependently.''
\noindent [257] National Institute of Diabetes and Digestive and Kidney Diseases (2018). Flavocoxid [LiverTox: Clinical and Research Information on Drug-Induced Liver Injury]. \textit{NCBI Bookshelf, National Library of Medicine} NBK548129 (last update 2018 Jan 30). \url{https://www.ncbi.nlm.nih.gov/books/NBK548129} (accessed 2026-11-14).
\noindent\hspace*{2em}\small\textit{Evidence:} ``In premarketing clinical trials, serum aminotransferase elevations occurred in up to 10\% of patients on flavocoxid therapy, but elevations above 3 times the upper limit of normal occurred in only 1\% to 2\% of recipients.''
\noindent\hspace*{2em}\small\textit{Evidence:} ``The time to onset was 1 to 5 months, and the pattern of enzyme elevations was usually hepatocellular or mixed.''
\noindent\hspace*{2em}\small\textit{Evidence:} ``Most cases were moderate in severity and no instance of acute liver failure or death has been reported.''
\noindent [258] U.S. Food and Drug Administration (2017). FDA requested recall: letter to Primus Pharmaceuticals, Inc. regarding Limbrel (flavocoxid) products. \textit{U.S. Food and Drug Administration} Letter of December 18, 2017. \url{https://www.fda.gov/media/109965/download} (accessed 2026-11-14).
\noindent\hspace*{2em}\small\textit{Evidence:} ``Between January 1, 2007, and November 9, 2017, FDA received 194 adverse event reports associated with the use of Limbrel products. The adverse event reports identified a close relationship between the use of Limbrel and adverse events involving the development of drug-induced liver injury (DILI), pancreatitis, and hypersensitivity pneumonitis (HP).''
\noindent\hspace*{2em}\small\textit{Evidence:} ``This letter is to request that you cease distribution and immediately initiate a recall of all lots within expiry of the following products:''
\noindent [259] Primus Pharmaceuticals, Inc. (2018). Primus Announces a Voluntary Nationwide Recall of All Lots Within Expiry of Prescription Medical Food. \textit{FDA Recalls, Market Withdrawals \& Safety Alerts} Company Announcement, January 2018. \url{https://www.fda.gov/safety/recalls-market-withdrawals-safety-alerts/primus-announces-voluntary-nationwide-recall-all-lots-within-expiry-prescription-medical-food} (accessed 2026-11-14).
\noindent\hspace*{2em}\small\textit{Evidence:} ``Between January 1, 2007, and November 9, 2017, FDA received 30 adverse event reports of elevated liver function tests or acute hypersensitivity pneumonitis associated with the use of Limbrel products.''
\noindent\hspace*{2em}\small\textit{Evidence:} ``There have been no deaths reported with the use of Limbrel, and in all reported cases adverse effects resolved without residual effects after discontinuing use of the product.''
\noindent\hspace*{2em}\small\textit{Evidence:} ``Primus voluntarily ceased its promotion and distribution of Limbrel on December 21, 2017, and is now recalling Limbrel as FDA has requested.''
\noindent [261] \href{https://pubmed.ncbi.nlm.nih.gov/42665171/}{PMID:42665171}. \url{https://pubmed.ncbi.nlm.nih.gov/42665171} (accessed 2026-11-14).
\noindent\hspace*{2em}\small\textit{Evidence:} ``the confirmed clinical incidence of SBG-DILI was 0.095\% (72/75,564).''
\noindent\hspace*{2em}\small\textit{Evidence:} ``daily doses exceeding the pharmacopoeia recommendation (>10 g/day) significantly increased the risk of liver injury (adjusted OR = 2.210, 95\% CI = 1.185-4.122, P = 0.013).''
\noindent\hspace*{2em}\small\textit{Evidence:} ``The homozygote C/C genotype was remarkably overrepresented in SBG-DILI patients (58.3\%) compared to other-DILI cases (4.5\%) and population controls (16.3\%).''
\noindent [262] \href{https://pubmed.ncbi.nlm.nih.gov/41431018/}{PMID:41431018}. \url{https://pubmed.ncbi.nlm.nih.gov/41431018} (accessed 2026-11-14).
\noindent\hspace*{2em}\small\textit{Evidence:} ``The report describes a rare case of IDILI induced by Scutellariae Radix, and the presence of the rs2032582 variant in ABCB1 was revealed in this patient, which may represent an individual observation of potential relevance.''
\noindent [263] \href{https://pubmed.ncbi.nlm.nih.gov/29434136/}{PMID:29434136}. \url{https://pubmed.ncbi.nlm.nih.gov/29434136} (accessed 2026-11-14).
\noindent\hspace*{2em}\small\textit{Evidence:} ``Both shosaikoto and saikokeishikankyoto contain Scutellariae Radix. This case is thought to be one of recurrent drug-induced liver injury caused by the incidental readministration of a Kampo formula containing Scutellariae Radix.''
\noindent [264] \href{https://pubmed.ncbi.nlm.nih.gov/20586134/}{PMID:20586134}. \url{https://pubmed.ncbi.nlm.nih.gov/20586134} (accessed 2026-11-14).
\noindent\hspace*{2em}\small\textit{Evidence:} ``Use of the Naranjo adverse drug reaction probability scale indicated a probable relationship (score of 6 for both patients) between the patients' development of hepatotoxicity and the Move Free Advanced supplement.''
\noindent [265] \href{https://pubmed.ncbi.nlm.nih.gov/22855699/}{PMID:22855699}. \url{https://pubmed.ncbi.nlm.nih.gov/22855699} (accessed 2026-11-14).
\noindent\hspace*{2em}\small\textit{Evidence:} ``Since neither the patient nor the treating physician recognized this supplement as a potential hepatotoxin, she resumed taking the supplement and again suffered from considerable hepatotoxicity.''
\noindent [266] \href{https://pubmed.ncbi.nlm.nih.gov/30957370/}{PMID:30957370}. \url{https://pubmed.ncbi.nlm.nih.gov/30957370} (accessed 2026-11-14).
\noindent [267] \href{https://pubmed.ncbi.nlm.nih.gov/42931707/}{PMID:42931707}. \url{https://pubmed.ncbi.nlm.nih.gov/42931707} (accessed 2026-11-14).
\noindent\hspace*{2em}\small\textit{Evidence:} ``The no-observed-effect level and no-observed-adverse-effect level were 1000 and 2000 mg/kg/day, respectively, and no target organ of toxicity was identified.''
\noindent\hspace*{2em}\small\textit{Evidence:} ``Acute administration caused only transient, non-adverse clinical signs, and the minimum lethal dose was >5000 mg/kg.''
\noindent [268] \href{https://pubmed.ncbi.nlm.nih.gov/40918531/}{PMID:40918531}. \url{https://pubmed.ncbi.nlm.nih.gov/40918531} (accessed 2026-11-14).
\noindent\hspace*{2em}\small\textit{Evidence:} ``Baicalin showed high acute safety with an LD50 over 4,000 mg/kg in mice, and a subacute no-observed-adverse-effect level (NOAEL) of 2,000 mg/kg, indicating its potential as a neuroprotective agent.''
\noindent\hspace*{2em}\small\textit{Evidence:} ``However, 4,000 mg/kg doses led to reversible hepatorenal toxicity and biochemical alterations, highlighting the need to monitor organ function during extended high-dose use.''
\noindent [269] \href{https://pubmed.ncbi.nlm.nih.gov/14686800/}{PMID:14686800}. \url{https://pubmed.ncbi.nlm.nih.gov/14686800} (accessed 2026-11-14).
\noindent\hspace*{2em}\small\textit{Evidence:} ``Treatment with BA (300 mg/kg, p.o.) 0.5 h after AP administration significantly prevented an increase in plasma alanine aminotransferase and aspartate aminotransferase activities and AP-induced hepatic necrosis, and also reduced AP-induced mortality from 43\% to 0\%.''
\noindent [270] \href{https://pubmed.ncbi.nlm.nih.gov/38913091/}{PMID:38913091}. \url{https://pubmed.ncbi.nlm.nih.gov/38913091} (accessed 2026-11-14).
\noindent\hspace*{2em}\small\textit{Evidence:} ``In the normal human liver cell line L02, baicalin treatment significantly inhibited AFB1-induced c-Jun-N-terminal Kinase (JNK) activation and cell apoptosis.''
\noindent [271] \href{https://pubmed.ncbi.nlm.nih.gov/17851565/}{PMID:17851565}. \url{https://pubmed.ncbi.nlm.nih.gov/17851565} (accessed 2026-11-14).
\noindent\hspace*{2em}\small\textit{Evidence:} ``After baicalin treatment, the area under the plasma concentration-time curve (AUC)(0-72) and AUC(0-infinity) of rosuvastatin decreased by 47.0+/-11.0\% (P=0.001) and 41.9+/-7.19\% (P=0.001) in OATP1B11b/1b,''
\noindent\hspace*{2em}\small\textit{Evidence:} ``Baicalin reduces plasma concentrations of rosuvastatin in an OATP1B1 haplotype-dependent manner.''
\noindent [272] \href{https://pubmed.ncbi.nlm.nih.gov/14994190/}{PMID:14994190}. \url{https://pubmed.ncbi.nlm.nih.gov/14994190} (accessed 2026-11-14).
\noindent\hspace*{2em}\small\textit{Evidence:} ``coadministration of SR decoction at doses of 1 g/kg and 2 g/kg significantly decreased the C (max) of cyclosporine by 62.9 \% and 79.6 \% and reduced the AUC (0 - 540) by 55.2 \% and 82.0 \%, respectively.''
\noindent\hspace*{2em}\small\textit{Evidence:} ``In order to ensure the efficacy and safety of cyclosporine, the coadministration of SR and its preparations with oral cyclosporine should be avoided.''
\noindent [273] \href{https://pubmed.ncbi.nlm.nih.gov/30663770/}{PMID:30663770}. \url{https://pubmed.ncbi.nlm.nih.gov/30663770} (accessed 2026-11-14).
\noindent\hspace*{2em}\small\textit{Evidence:} ``Multiple doses of BG decreased the oral bioavailability of CsA in rats significantly, which may be mainly attributable to inhibition of absorption of CsA in intestine and induction of P-gp.''
\noindent\hspace*{2em}\small\textit{Evidence:} ``The dosage adjustment and blood concentration monitoring of CsA may be required in clinic.''
\noindent [274] \href{https://pubmed.ncbi.nlm.nih.gov/24587011/}{PMID:24587011}. \url{https://pubmed.ncbi.nlm.nih.gov/24587011} (accessed 2026-11-14).
\noindent\hspace*{2em}\small\textit{Evidence:} ``baicalin exhibited a mixed-type inhibition in pooled HLMs, with a Ki value of 25.4 µM.''
\noindent\hspace*{2em}\small\textit{Evidence:} ``In conclusion, baicalin can inhibit the activity of CYP1A2 in HLMs and exhibit large inter-individual variation that has no relationship with gene polymorphism.''
\noindent [275] \href{https://pubmed.ncbi.nlm.nih.gov/19074529/}{PMID:19074529}. \url{https://pubmed.ncbi.nlm.nih.gov/19074529} (accessed 2026-11-14).
\noindent\hspace*{2em}\small\textit{Evidence:} ``This article describes an in vitro investigation of the inhibition of cytochrome P450 (P450) 2C9 by a series of flavonoids made up of flavones (flavone, 6-hydroxyflavone, 7-hydroxyflavone, chrysin, baicalein, apigenin, luteolin, scutellarein, and wogonin) and flavonols (galangin, fisetin, kaempferol, morin, and quercetin). With the exception of flavone, all flavonoids were shown to inhibit CYP2C9-mediated diclofenac 4'-hydroxylation in the CYP2C9 RECO system, with K(i) value <or= 2.2 microM.''
\noindent [276] \href{https://pubmed.ncbi.nlm.nih.gov/34214011/}{PMID:34214011}. \url{https://pubmed.ncbi.nlm.nih.gov/34214011} (accessed 2026-11-14).
\noindent\hspace*{2em}\small\textit{Evidence:} ``Baicalein showed an inhibitory effect on the activity of CYP3A4 with the IC50 value of 12.03 μM, which is responsible for the metabolism of simvastatin.''
\noindent\hspace*{2em}\small\textit{Evidence:} ``The co-administration of baicalein and simvastatin may induce drug-drug interaction through inhibiting CYP3A4.''
\noindent [278] \href{https://pubmed.ncbi.nlm.nih.gov/25268843/}{PMID:25268843}. \url{https://pubmed.ncbi.nlm.nih.gov/25268843} (accessed 2026-11-14).
\noindent\hspace*{2em}\small\textit{Evidence:} ``In the present study, baicalein inhibited agonist-induced platelet aggregation, granule secretion markers (P-selectin expression and ATP release), [Ca(2+)]i mobilization, and integrin αIIbβ3 expression.''
\noindent [279] \href{https://pubmed.ncbi.nlm.nih.gov/24849036/}{PMID:24849036}. \url{https://pubmed.ncbi.nlm.nih.gov/24849036} (accessed 2026-11-14).
\noindent\hspace*{2em}\small\textit{Evidence:} ``Our data showed that baicalin inhibited thrombin-catalyzed fibrin polymerization and platelet functions, prolonged aPTT and PT significantly and inhibited the activities and production of thrombin and FXa.''
\noindent\hspace*{2em}\small\textit{Evidence:} ``Baicalin also prolonged in vivo bleeding time and inhibited TNF-α induced PAI-1 production.''
\noindent [280] \href{https://pubmed.ncbi.nlm.nih.gov/25020199/}{PMID:25020199}. \url{https://pubmed.ncbi.nlm.nih.gov/25020199} (accessed 2026-11-14).
\noindent\hspace*{2em}\small\textit{Evidence:} ``In addition, WGN(S) inhibited thrombin-catalyzed fibrin polymerization and platelet aggregation. WGN(S) also elicited anticoagulant effects in mice.''
\noindent [282] \href{https://pubmed.ncbi.nlm.nih.gov/42170474/}{PMID:42170474}. \url{https://pubmed.ncbi.nlm.nih.gov/42170474} (accessed 2026-11-14).
\noindent\hspace*{2em}\small\textit{Evidence:} ``TXN and BCN significantly reduced blood glucose levels, improved lipid profiles, and decreased oxidative stress in diabetic rats.''
\noindent [283] \href{https://pubmed.ncbi.nlm.nih.gov/36216728/}{PMID:36216728}. \url{https://pubmed.ncbi.nlm.nih.gov/36216728} (accessed 2026-11-14).
\noindent\hspace*{2em}\small\textit{Evidence:} ``In ancient China, S. baicalensis root was used to clear heat, protect the fetus, and avoid a miscarriage for thousands of years.''
\noindent\hspace*{2em}\small\textit{Evidence:} ``The mechanisms, potential targets and experimental models of S. baicalensis root for ameliorating pregnancy-related diseases, such as recurrent spontaneous abortion, preeclampsia, preterm birth, fetal growth restriction and gestational diabetes mellitus, are highlighted.''
\noindent [284] \href{https://pubmed.ncbi.nlm.nih.gov/8448419/}{PMID:8448419}. \url{https://pubmed.ncbi.nlm.nih.gov/8448419} (accessed 2026-11-14).
\noindent\hspace*{2em}\small\textit{Evidence:} ``There was a significant (P < 0.05) dose-dependent increase in the incidence of skeletal variations (presence of lumbar rib).''
\noindent [285] \href{https://pubmed.ncbi.nlm.nih.gov/20549833/}{PMID:20549833}. \url{https://pubmed.ncbi.nlm.nih.gov/20549833} (accessed 2026-11-14).
\noindent\hspace*{2em}\small\textit{Evidence:} ``The present data suggest that the appearance of lumbar rib induced by SR is a transient fetal variation rather than teratogenicity or maternal toxicity.''
\noindent\hspace*{2em}\small\textit{Evidence:} ``at a dose of 186 mg/kg/day, equivalent to 25 g/kg of starting material, representing a 100-fold increase over typical human intake level.''
\noindent [286] \href{https://pubmed.ncbi.nlm.nih.gov/19283861/}{PMID:19283861}. \url{https://pubmed.ncbi.nlm.nih.gov/19283861} (accessed 2026-11-14).
\noindent\hspace*{2em}\small\textit{Evidence:} ``Oral administration of aqueous extract of S. baicalensis roots at or below 32 g/kg/day to ICR mice during organogenesis did not cause significant fetal external or skeletal malformations.''
\noindent\hspace*{2em}\small\textit{Evidence:} ``However, 32 g/kg/day presented potential maternal toxicity.''
\noindent [287] Wang J; Zhang M; Chen L; Qiao Y; Ma S; Sun D; Si J; Liao Y (2021). Determination of Toxic Pyrrolizidine Alkaloids in Traditional Chinese Herbal Medicines by UPLC-MS/MS and Accompanying Risk Assessment for Human Health. \textit{Molecules} 26(6):1648. \href{https://doi.org/10.3390/molecules26061648}{doi:10.3390/molecules26061648}. \url{https://pmc.ncbi.nlm.nih.gov/articles/PMC8000276} (accessed 2026-11-14).
\noindent\hspace*{2em}\small\textit{Evidence:} ``Results showed that PAs was detected in 271 out of 386 samples with a content of 0.1–25,567.4 μg/kg, and there were 20 samples with EDI values above the baseline, 0.007 μg/kg bw/day.''
\noindent\hspace*{2em}\small\textit{Evidence:} ``HM-11 | \_Scutellaria barbata\_ | 5 (LyNO, Ly, ImNO) | 61.7 | 30 | 1.8 | 0.026 | 8958.3 | 232,917.5''
\noindent\hspace*{2em}\small\textit{Evidence:} ``the MOE values for 10 out of 271 positive samples were below 10,000.''
\noindent [288] \href{https://pubmed.ncbi.nlm.nih.gov/34075510/}{PMID:34075510}. \url{https://pubmed.ncbi.nlm.nih.gov/34075510} (accessed 2026-11-14).
\noindent\hspace*{2em}\small\textit{Evidence:} ``This research aims to determine the levels of Cr, Ni, Cu, Zn, As, Cd, Hg and Pb in four medicinal plant species (Aster tataricus L.f., Salvia miltiorrhiza Bge, Radix Aucklandiae, Scutellaria baicalensis Georgi) and their native soil.''
\noindent\hspace*{2em}\small\textit{Evidence:} ``The contents of Cu, Cd, Hg and Pb in some medicinal herbs exceeded the standards.''
\noindent\hspace*{2em}\small\textit{Evidence:} ``The comprehensive health risk assessment of heavy metals with chronic non-carcinogenic effects for human body showed that none of the four medicinal herbs can create a health risk.''
\noindent [289] \href{https://pubmed.ncbi.nlm.nih.gov/41003537/}{PMID:41003537}. \url{https://pubmed.ncbi.nlm.nih.gov/41003537} (accessed 2026-11-14).
\noindent\hspace*{2em}\small\textit{Evidence:} ``Application of the method to six commercial Scutellaria baicalensis samples revealed detectable AFB1 in two samples, though all levels were below national safety limits.''
\noindent [291] \href{https://pubmed.ncbi.nlm.nih.gov/25184892/}{PMID:25184892}. \url{https://pubmed.ncbi.nlm.nih.gov/25184892} (accessed 2026-11-14).
\noindent\hspace*{2em}\small\textit{Evidence:} ``due to intentional or unintentional adulteration, cases of hepatotoxicity have been reported.''
\noindent [293] Insights into skullcap herb-induced liver injury.. \url{https://pubmed.ncbi.nlm.nih.gov/38495279} (accessed 2026-11-14).
\noindent\hspace*{2em}\small\textit{Evidence:} ``skullcap herb-induced liver injury.''



\end{document}
